% !TeX program = lualatex
% ALIUS Bulletin native visual reconstruction source.
% Generated from off-repo reference layout data; does not include or input a pre-existing article PDF.
\ifdefined\ALIUSIssueBuild
\else
\documentclass{article}
\usepackage[paperwidth=595.2756bp,paperheight=841.8898bp,margin=0bp]{geometry}
\usepackage{graphicx}
\usepackage{xcolor}
\usepackage{tikz}
\usepackage{iftex}
\ifPDFTeX
  \usepackage[T1]{fontenc}
  \usepackage[utf8]{inputenc}
  \usepackage{textcomp}
  \usepackage{amssymb}
\else
  \usepackage{fontspec}
\fi
\pagestyle{empty}
\fi
\ifPDFTeX
  \providecommand{\ALIUSDeclareUnicodeFallbacks}{%
    \DeclareUnicodeCharacter{00AC}{\ensuremath{\neg}}%
    \DeclareUnicodeCharacter{00BE}{\ensuremath{\frac{3}{4}}}%
    \DeclareUnicodeCharacter{0101}{\=a}%
    \DeclareUnicodeCharacter{0107}{\'c}%
    \DeclareUnicodeCharacter{012B}{\=\i}%
    \DeclareUnicodeCharacter{0131}{\i}%
    \DeclareUnicodeCharacter{0142}{\l}%
    \DeclareUnicodeCharacter{015B}{\'s}%
    \DeclareUnicodeCharacter{016B}{\=u}%
    \DeclareUnicodeCharacter{025C}{q}%
    \DeclareUnicodeCharacter{0278}{t}%
    \DeclareUnicodeCharacter{02AE}{Th}%
    \DeclareUnicodeCharacter{02B0}{ff}%
    \DeclareUnicodeCharacter{02B4}{ffi}%
    \DeclareUnicodeCharacter{02BE}{ft}%
    \DeclareUnicodeCharacter{02BF}{fi}%
    \DeclareUnicodeCharacter{02C0}{fl}%
    \DeclareUnicodeCharacter{02C1}{fi}%
    \DeclareUnicodeCharacter{02D9}{Th}%
    \DeclareUnicodeCharacter{02EF}{gy}%
    \DeclareUnicodeCharacter{0394}{\ensuremath{\Delta}}%
    \DeclareUnicodeCharacter{03B2}{\ensuremath{\beta}}%
    \DeclareUnicodeCharacter{03BA}{\ensuremath{\kappa}}%
    \DeclareUnicodeCharacter{068D}{-}%
    \DeclareUnicodeCharacter{08BC}{ti}%
    \DeclareUnicodeCharacter{099E}{ti}%
    \DeclareUnicodeCharacter{1E43}{\d{m}}%
    \DeclareUnicodeCharacter{1E47}{\d{n}}%
    \DeclareUnicodeCharacter{1E63}{\d{s}}%
    \DeclareUnicodeCharacter{201C}{``}%
    \DeclareUnicodeCharacter{201D}{''}%
    \DeclareUnicodeCharacter{2260}{\ensuremath{\neq}}%
    \DeclareUnicodeCharacter{25A1}{\ensuremath{\square}}%
  }%
  \ifdefined\ALIUSIssueBuild\else\ALIUSDeclareUnicodeFallbacks\fi
  \providecommand{\ALIUSFontLatoLight}{\fontfamily{lato-LF}\fontseries{l}\selectfont}
  \providecommand{\ALIUSFontLatoRegular}{\fontfamily{lato-LF}\fontseries{m}\selectfont}
  \providecommand{\ALIUSFontLatoItalic}{\fontfamily{lato-LF}\fontseries{m}\itshape}
  \providecommand{\ALIUSFontLatoLightItalic}{\fontfamily{lato-LF}\fontseries{l}\itshape}
  \providecommand{\ALIUSFontCormorantRegular}{\fontfamily{CormorantGaramond-LF}\fontseries{m}\selectfont}
  \providecommand{\ALIUSFontCormorantLight}{\fontfamily{CormorantGaramond-LF}\fontseries{l}\selectfont}
  \providecommand{\ALIUSFontCormorantMedium}{\fontfamily{CormorantGaramond-LF}\fontseries{medium}\selectfont}
  \providecommand{\ALIUSFontCormorantBold}{\fontfamily{CormorantGaramond-LF}\fontseries{b}\selectfont}
  \providecommand{\ALIUSFontCormorantItalic}{\fontfamily{CormorantGaramond-LF}\fontseries{m}\itshape}
  \providecommand{\ALIUSFontTimes}{\fontfamily{ptm}\selectfont}
  \providecommand{\ALIUSFontCalibri}{\fontfamily{phv}\selectfont}
  \providecommand{\ALIUSFontCambria}{\fontfamily{ptm}\selectfont}
  \providecommand{\ALIUSFontSymbolFallback}{\fontfamily{ptm}\selectfont}
\else
  \providecommand{\ALIUSFontLatoLight}{\fontspec{Lato Light}}
  \providecommand{\ALIUSFontLatoRegular}{\fontspec{Lato Regular}}
  \providecommand{\ALIUSFontLatoItalic}{\fontspec{Lato Italic}}
  \providecommand{\ALIUSFontLatoLightItalic}{\fontspec{Lato Light Italic}}
  \providecommand{\ALIUSFontCormorantRegular}{\fontspec{Cormorant Garamond Regular}}
  \providecommand{\ALIUSFontCormorantLight}{\fontspec{Cormorant Garamond Light}}
  \providecommand{\ALIUSFontCormorantMedium}{\fontspec{Cormorant Garamond Medium}}
  \providecommand{\ALIUSFontCormorantBold}{\fontspec{Cormorant Garamond Bold}}
  \providecommand{\ALIUSFontCormorantItalic}{\fontspec{Cormorant Garamond Italic}}
  \providecommand{\ALIUSFontTimes}{\fontspec{Times New Roman}}
  \providecommand{\ALIUSFontCalibri}{\fontspec{Calibri}}
  \providecommand{\ALIUSFontCambria}{\fontspec{Cambria}}
  \providecommand{\ALIUSFontSymbolFallback}{\fontspec{Times New Roman}}
\fi
\definecolor{ALIUSC000000}{HTML}{000000}
\definecolor{ALIUSC0000FF}{HTML}{0000FF}
\definecolor{ALIUSC1F8135}{HTML}{1F8135}
\definecolor{ALIUSC595959}{HTML}{595959}
\definecolor{ALIUSC767171}{HTML}{767171}
\definecolor{ALIUSCBFBFBF}{HTML}{BFBFBF}
\providecommand{\ALIUSPlacedText}[6]{%
  \node[anchor=north west,inner sep=0pt,outer sep=0pt,text=#4] at (#1bp,-#2bp) {%
    \resizebox{#3bp}{!}{{#5\fontsize{#6bp}{#6bp}\selectfont #6\kern0pt}}%
  };%
}
\providecommand{\ALIUSPlacedTextContent}[7]{%
  \node[anchor=north west,inner sep=0pt,outer sep=0pt,text=#4] at (#1bp,-#2bp) {%
    \resizebox{#3bp}{!}{{#5\fontsize{#6bp}{#6bp}\selectfont #7}}%
  };%
}

\ifdefined\ALIUSIssueBuild
\else
\begin{document}
\fi

% Page 1
\thispagestyle{empty}
\noindent\begin{tikzpicture}[remember picture,overlay,x=1bp,y=1bp,shift={(current page.north west)}]
  \fill[ALIUSC1F8135] (351.591bp,-192.231bp) rectangle ++(155.760bp,-0.240bp);
  \fill[ALIUSC1F8135] (351.591bp,-252.231bp) rectangle ++(90.240bp,-0.240bp);
  \fill[ALIUSC1F8135] (351.591bp,-308.631bp) rectangle ++(73.680bp,-0.240bp);
  \fill[ALIUSCBFBFBF] (345.831bp,-169.191bp) rectangle ++(0.480bp,-88.800bp);
  \fill[ALIUSC0000FF] (96.231bp,-331.191bp) rectangle ++(163.920bp,-0.480bp);
  \fill[ALIUSCBFBFBF] (345.831bp,-257.991bp) rectangle ++(0.480bp,-2.160bp);
  \fill[ALIUSCBFBFBF] (345.831bp,-260.151bp) rectangle ++(0.480bp,-72.240bp);
  \ALIUSPlacedTextContent{90.951}{72.970}{276.116}{ALIUSC000000}{\ALIUSFontLatoRegular}{19.920}{Bodily boundaries and beyond:}
  \ALIUSPlacedTextContent{90.951}{108.970}{249.655}{ALIUSC000000}{\ALIUSFontLatoLight}{19.920}{Exploring the malleability of}
  \ALIUSPlacedTextContent{90.951}{132.970}{223.556}{ALIUSC000000}{\ALIUSFontLatoLight}{19.920}{bodily self-consciousness}
  \ALIUSPlacedTextContent{351.591}{169.094}{100.835}{ALIUSC000000}{\ALIUSFontLatoRegular}{11.040}{Bigna Lenggenhager}
  \ALIUSPlacedTextContent{351.591}{182.269}{163.127}{ALIUSC1F8135}{\ALIUSFontLatoLight}{9.120}{bigna.lenggenhager@psychologie.uzh.ch}
  \ALIUSPlacedTextContent{96.231}{188.360}{130.764}{ALIUSC000000}{\ALIUSFontLatoLight}{16.080}{An interview with}
  \ALIUSPlacedTextContent{351.591}{193.069}{83.875}{ALIUSC000000}{\ALIUSFontLatoLight}{9.120}{University of Zürich,}
  \ALIUSPlacedTextContent{351.591}{203.869}{76.495}{ALIUSC000000}{\ALIUSFontLatoLight}{9.120}{Zürich, Switzerland}
  \ALIUSPlacedTextContent{96.231}{207.597}{171.595}{ALIUSC000000}{\ALIUSFontLatoLight}{18.960}{Bigna Lenggenhager}
  \ALIUSPlacedTextContent{351.591}{229.094}{80.479}{ALIUSC000000}{\ALIUSFontLatoRegular}{11.040}{Raphaël Millière}
  \ALIUSPlacedTextContent{351.591}{242.269}{92.704}{ALIUSC1F8135}{\ALIUSFontLatoLight}{9.120}{rm3799@columbia.edu}
  \ALIUSPlacedTextContent{96.231}{242.319}{205.521}{ALIUSC000000}{\ALIUSFontLatoLight}{12.960}{by Jasmine T. Ho \& Raphaël Millière}
  \ALIUSPlacedTextContent{351.591}{253.069}{82.804}{ALIUSC000000}{\ALIUSFontLatoLight}{9.120}{Columbia University,}
  \ALIUSPlacedTextContent{96.231}{260.282}{247.115}{ALIUSC000000}{\ALIUSFontLatoLight}{10.080}{Cite as: Lenggenhager, B., Ho, J. T. \& Millière, R. (2020).}
  \ALIUSPlacedTextContent{351.591}{263.869}{60.413}{ALIUSC000000}{\ALIUSFontLatoLight}{9.120}{New York, USA}
  \ALIUSPlacedTextContent{96.231}{272.282}{247.111}{ALIUSC000000}{\ALIUSFontLatoLight}{10.080}{Bodily boundaries and beyond: Exploring the malleability}
  \ALIUSPlacedTextContent{96.231}{284.282}{247.093}{ALIUSC000000}{\ALIUSFontLatoLight}{10.080}{of bodily self-consciousness. An interview with Bigna}
  \ALIUSPlacedTextContent{351.591}{285.494}{70.695}{ALIUSC000000}{\ALIUSFontLatoRegular}{11.040}{Jasmine T. Ho}
  \ALIUSPlacedTextContent{96.231}{296.282}{249.807}{ALIUSC000000}{\ALIUSFontLatoLight}{10.080}{Lenggenhager by Jasmine T. Ho and Raphaël Millière.}
  \ALIUSPlacedTextContent{351.591}{298.669}{76.266}{ALIUSC1F8135}{\ALIUSFontLatoLight}{9.120}{jasmine.ho@uzh.ch}
  \ALIUSPlacedTextContent{96.231}{308.282}{27.503}{ALIUSC000000}{\ALIUSFontLatoLightItalic}{10.080}{ALIUS}
  \ALIUSPlacedTextContent{168.909}{308.282}{30.642}{ALIUSC000000}{\ALIUSFontLatoLightItalic}{10.080}{Bulletin}
  \ALIUSPlacedTextContent{199.502}{308.282}{4.932}{ALIUSC000000}{\ALIUSFontLatoLight}{10.080}{,}
  \ALIUSPlacedTextContent{249.697}{308.282}{5.846}{ALIUSC000000}{\ALIUSFontLatoLightItalic}{10.080}{4}
  \ALIUSPlacedTextContent{255.497}{308.282}{4.932}{ALIUSC000000}{\ALIUSFontLatoLight}{10.080}{,}
  \ALIUSPlacedTextContent{305.692}{308.282}{37.631}{ALIUSC000000}{\ALIUSFontLatoLight}{10.080}{82-105,}
  \ALIUSPlacedTextContent{351.591}{309.469}{81.432}{ALIUSC000000}{\ALIUSFontLatoLight}{9.120}{University of Zürich,}
  \ALIUSPlacedTextContent{351.591}{320.269}{74.072}{ALIUSC000000}{\ALIUSFontLatoLight}{9.120}{Zürich, Switzerland}
  \ALIUSPlacedTextContent{96.231}{320.282}{164.082}{ALIUSC0000FF}{\ALIUSFontLatoRegular}{10.080}{https://doi.org/10.34700/2vez-qq65}
  \ALIUSPlacedTextContent{90.951}{389.452}{55.719}{ALIUSC000000}{\ALIUSFontLatoRegular}{13.920}{Abstract}
  \ALIUSPlacedTextContent{90.951}{417.987}{418.270}{ALIUSC000000}{\ALIUSFontLatoLight}{12.000}{In this interview, Bigna Lenggenhager discusses her groundbreaking empirical}
  \ALIUSPlacedTextContent{90.951}{432.387}{418.272}{ALIUSC000000}{\ALIUSFontLatoLight}{12.000}{work on bodily self-consciousness, bodily disorders and bodily illusions. The}
  \ALIUSPlacedTextContent{90.951}{446.787}{418.270}{ALIUSC000000}{\ALIUSFontLatoLight}{12.000}{conversation explores issues related to the interpretation of the rubber hand}
  \ALIUSPlacedTextContent{90.951}{461.187}{415.000}{ALIUSC000000}{\ALIUSFontLatoLight}{12.000}{illusion and the full-body illusion, the nature of the relationship between self-}
  \ALIUSPlacedTextContent{90.951}{475.587}{418.266}{ALIUSC000000}{\ALIUSFontLatoLight}{12.000}{consciousness and bodily awareness, syndromes of disembodiment, as well as}
  \ALIUSPlacedTextContent{90.951}{489.987}{341.060}{ALIUSC000000}{\ALIUSFontLatoLight}{12.000}{the use of virtual reality as a therapeutic tool for bodily disorders.}
  \ALIUSPlacedTextContent{90.951}{516.387}{47.340}{ALIUSC000000}{\ALIUSFontLatoItalic}{12.000}{keywords}
  \ALIUSPlacedTextContent{138.289}{516.387}{370.686}{ALIUSC000000}{\ALIUSFontLatoLightItalic}{12.000}{: self-consciousness, bodily awareness, rubber hand illusion, full-body}
  \ALIUSPlacedTextContent{90.951}{530.787}{111.024}{ALIUSC000000}{\ALIUSFontLatoLightItalic}{12.000}{illusion, virtual reality}
  \ALIUSPlacedTextContent{90.951}{574.293}{418.451}{ALIUSC1F8135}{\ALIUSFontLatoLight}{12.357}{What sparked your interest in bodily self-consciousness? What types of}
  \ALIUSPlacedTextContent{90.951}{591.573}{323.813}{ALIUSC1F8135}{\ALIUSFontLatoLight}{12.357}{questions were you aiming to answer by entering this field?}
  \ALIUSPlacedTextContent{90.951}{626.009}{418.185}{ALIUSC000000}{\ALIUSFontCormorantRegular}{13.920}{As a child, I read, or rather tried to read, many of my father's books on yoga,}
  \ALIUSPlacedTextContent{90.951}{645.448}{418.162}{ALIUSC000000}{\ALIUSFontCormorantRegular}{13.920}{meditation, and Buddhism. Although I didn't understand much, I was}
  \ALIUSPlacedTextContent{90.951}{664.889}{418.314}{ALIUSC000000}{\ALIUSFontCormorantRegular}{13.920}{fascinated by the different understanding of the self, life, and death in these}
  \ALIUSPlacedTextContent{90.951}{684.568}{418.255}{ALIUSC000000}{\ALIUSFontCormorantRegular}{13.920}{traditions. Later during high school, my desire grew to study the human brain}
  \ALIUSPlacedTextContent{90.951}{704.009}{418.109}{ALIUSC000000}{\ALIUSFontCormorantRegular}{13.920}{in its complexity, hoping to find answers to some of the big questions about}
  \ALIUSPlacedTextContent{90.951}{723.448}{418.075}{ALIUSC000000}{\ALIUSFontCormorantRegular}{13.920}{how we perceive, think, and act, and why we are who we are. I went to}
  \ALIUSPlacedTextContent{90.951}{742.888}{418.249}{ALIUSC000000}{\ALIUSFontCormorantRegular}{13.920}{university to study psychology, psychopathology, and neurophysiology.}
  \ALIUSPlacedTextContent{90.951}{793.143}{122.397}{ALIUSC767171}{\ALIUSFontCormorantLight}{12.000}{ALIUS Bulletin n°4 (2020)}
  \ALIUSPlacedTextContent{387.435}{793.143}{120.728}{ALIUSC767171}{\ALIUSFontCormorantLight}{12.000}{aliusresearch.org/bulletin}
\end{tikzpicture}
\null
\newpage

% Page 2
\thispagestyle{empty}
\noindent\begin{tikzpicture}[remember picture,overlay,x=1bp,y=1bp,shift={(current page.north west)}]
  \ALIUSPlacedTextContent{90.951}{36.329}{258.472}{ALIUSC767171}{\ALIUSFontCormorantRegular}{13.920}{Lenggenhager – Bodily boundaries and beyond}
  \ALIUSPlacedTextContent{493.671}{36.329}{15.549}{ALIUSC000000}{\ALIUSFontCormorantRegular}{13.920}{83}
  \ALIUSPlacedTextContent{90.951}{73.048}{418.255}{ALIUSC000000}{\ALIUSFontCormorantRegular}{13.920}{During this time, I got interested in illusions and their value in revealing how}
  \ALIUSPlacedTextContent{90.951}{92.488}{418.346}{ALIUSC000000}{\ALIUSFontCormorantRegular}{13.920}{the brain works. I was intrigued by visual and multisensory illusions, but}
  \ALIUSPlacedTextContent{90.951}{111.928}{418.162}{ALIUSC000000}{\ALIUSFontCormorantRegular}{13.920}{above all by bodily illusions, such as the rubber hand illusion (Botvinick \&}
  \ALIUSPlacedTextContent{90.951}{131.609}{418.261}{ALIUSC000000}{\ALIUSFontCormorantRegular}{13.920}{Cohen, 1998) and Pinocchio illusion (Lackner, 1988). I found it fascinating}
  \ALIUSPlacedTextContent{90.951}{151.049}{418.244}{ALIUSC000000}{\ALIUSFontCormorantRegular}{13.920}{how easily the perception of the most familiar object, i.e., one's own body,}
  \ALIUSPlacedTextContent{90.951}{170.488}{418.285}{ALIUSC000000}{\ALIUSFontCormorantRegular}{13.920}{can be altered by just using a clever experimental design and without any}
  \ALIUSPlacedTextContent{90.951}{189.928}{418.269}{ALIUSC000000}{\ALIUSFontCormorantRegular}{13.920}{drugs or long meditation techniques. But definitively the strongest driver of}
  \ALIUSPlacedTextContent{90.951}{209.609}{418.187}{ALIUSC000000}{\ALIUSFontCormorantRegular}{13.920}{my interest in bodily self-consciousness were the patients I investigated}
  \ALIUSPlacedTextContent{90.951}{229.049}{418.307}{ALIUSC000000}{\ALIUSFontCormorantRegular}{13.920}{during my first internship in the Neurology Department of the University}
  \ALIUSPlacedTextContent{90.951}{248.489}{418.257}{ALIUSC000000}{\ALIUSFontCormorantRegular}{13.920}{Hospital of Zurich. Listening to their stories made me want to understand}
  \ALIUSPlacedTextContent{90.951}{267.928}{418.382}{ALIUSC000000}{\ALIUSFontCormorantRegular}{13.920}{everything about how we perceive our body and how it links to our sense of}
  \ALIUSPlacedTextContent{90.951}{287.608}{418.303}{ALIUSC000000}{\ALIUSFontCormorantRegular}{13.920}{self and consciousness. Realizing how many aspects of the perception of the}
  \ALIUSPlacedTextContent{90.951}{307.048}{418.387}{ALIUSC000000}{\ALIUSFontCormorantRegular}{13.920}{self can be altered, I wanted to learn why and how we normally perceive our}
  \ALIUSPlacedTextContent{90.951}{326.488}{418.232}{ALIUSC000000}{\ALIUSFontCormorantRegular}{13.920}{body (but not another person's body) as belonging to us, how we feel in}
  \ALIUSPlacedTextContent{90.951}{345.928}{418.212}{ALIUSC000000}{\ALIUSFontCormorantRegular}{13.920}{control of our body, and how we perceive and act from a physically embodied}
  \ALIUSPlacedTextContent{90.951}{365.608}{64.379}{ALIUSC000000}{\ALIUSFontCormorantRegular}{13.920}{perspective.}
  \ALIUSPlacedTextContent{90.951}{404.613}{418.471}{ALIUSC1F8135}{\ALIUSFontLatoLight}{12.357}{There are many disorders of bodily self-consciousness. Is there a specific}
  \ALIUSPlacedTextContent{90.951}{421.893}{285.401}{ALIUSC1F8135}{\ALIUSFontLatoLight}{12.357}{condition that elicits a particular fascination for you?}
  \ALIUSPlacedTextContent{90.951}{458.488}{414.974}{ALIUSC000000}{\ALIUSFontCormorantRegular}{13.920}{From the external perspective as a researcher, all disorders of bodily self-}
  \ALIUSPlacedTextContent{90.951}{477.928}{418.196}{ALIUSC000000}{\ALIUSFontCormorantRegular}{13.920}{consciousness are extremely fascinating to me. It is impressive how many}
  \ALIUSPlacedTextContent{90.951}{497.608}{418.223}{ALIUSC000000}{\ALIUSFontCormorantRegular}{13.920}{different aspects of bodily self-consciousness can be affected. I think this}
  \ALIUSPlacedTextContent{90.951}{517.049}{418.215}{ALIUSC000000}{\ALIUSFontCormorantRegular}{13.920}{broad range of symptoms makes bodily self-consciousness disorders so}
  \ALIUSPlacedTextContent{90.951}{536.489}{418.267}{ALIUSC000000}{\ALIUSFontCormorantRegular}{13.920}{fascinating, ranging from the feeling of being duplicated or even multiplied,}
  \ALIUSPlacedTextContent{90.951}{555.928}{418.101}{ALIUSC000000}{\ALIUSFontCormorantRegular}{13.920}{as in a specific form of heautoscopy (Brugger et al., 2006), or the feeling that}
  \ALIUSPlacedTextContent{90.951}{575.369}{237.374}{ALIUSC000000}{\ALIUSFontCormorantRegular}{13.920}{one does not exist, as in Cotard syndrome.}
  \ALIUSPlacedTextContent{90.951}{614.489}{418.242}{ALIUSC000000}{\ALIUSFontCormorantRegular}{13.920}{What's more, different disorders fascinate me for different reasons. Take for}
  \ALIUSPlacedTextContent{90.951}{633.928}{418.194}{ALIUSC000000}{\ALIUSFontCormorantRegular}{13.920}{example, out-of- body experiences: I find it extremely fascinating that these}
  \ALIUSPlacedTextContent{90.951}{653.369}{418.345}{ALIUSC000000}{\ALIUSFontCormorantRegular}{13.920}{experiences have been described in so many different contexts and cultures,}
  \ALIUSPlacedTextContent{90.951}{673.048}{418.263}{ALIUSC000000}{\ALIUSFontCormorantRegular}{13.920}{and that they might actually have shaped the way people think about the}
  \ALIUSPlacedTextContent{90.951}{692.489}{418.270}{ALIUSC000000}{\ALIUSFontCormorantRegular}{13.920}{relation between the body and the mind, and between life and death. And}
  \ALIUSPlacedTextContent{90.951}{711.928}{418.410}{ALIUSC000000}{\ALIUSFontCormorantRegular}{13.920}{the fact that such complex illusions can be induced by local brain stimulation}
  \ALIUSPlacedTextContent{90.951}{731.369}{258.934}{ALIUSC000000}{\ALIUSFontCormorantRegular}{13.920}{(Blanke et al., 2002) remains fascinating to me.}
  \ALIUSPlacedTextContent{90.951}{793.143}{122.397}{ALIUSC767171}{\ALIUSFontCormorantLight}{12.000}{ALIUS Bulletin n°4 (2020)}
  \ALIUSPlacedTextContent{387.435}{793.143}{120.728}{ALIUSC767171}{\ALIUSFontCormorantLight}{12.000}{aliusresearch.org/bulletin}
\end{tikzpicture}
\null
\newpage

% Page 3
\thispagestyle{empty}
\noindent\begin{tikzpicture}[remember picture,overlay,x=1bp,y=1bp,shift={(current page.north west)}]
  \ALIUSPlacedTextContent{90.951}{36.329}{258.472}{ALIUSC767171}{\ALIUSFontCormorantRegular}{13.920}{Lenggenhager – Bodily boundaries and beyond}
  \ALIUSPlacedTextContent{492.711}{36.329}{16.417}{ALIUSC000000}{\ALIUSFontCormorantRegular}{13.920}{84}
  \ALIUSPlacedTextContent{90.951}{73.048}{418.282}{ALIUSC000000}{\ALIUSFontCormorantRegular}{13.920}{Another example is Body Integrity Dysphoria, a disorder in which otherwise}
  \ALIUSPlacedTextContent{90.951}{92.488}{418.263}{ALIUSC000000}{\ALIUSFontCormorantRegular}{13.920}{healthy individuals feel like a part of their body does not belong to them}
  \ALIUSPlacedTextContent{90.951}{111.928}{418.240}{ALIUSC000000}{\ALIUSFontCormorantRegular}{13.920}{(Brugger et al., 2013). Here, the most fascinating aspect for me is probably the}
  \ALIUSPlacedTextContent{90.951}{131.609}{418.267}{ALIUSC000000}{\ALIUSFontCormorantRegular}{13.920}{related societal and ethical questions, which heavily depend on our}
  \ALIUSPlacedTextContent{90.951}{151.049}{293.134}{ALIUSC000000}{\ALIUSFontCormorantRegular}{13.920}{conception of how body, brain, and mind are linked.}
  \ALIUSPlacedTextContent{90.951}{187.893}{418.398}{ALIUSC1F8135}{\ALIUSFontLatoLight}{12.357}{In 2007, you published the first experimental study investigating the so-called}
  \ALIUSPlacedTextContent{90.951}{205.173}{418.369}{ALIUSC1F8135}{\ALIUSFontLatoLight}{12.357}{full-body illusion (Lenggenhager et al., 2007). Since this landmark publication,}
  \ALIUSPlacedTextContent{90.951}{222.453}{418.439}{ALIUSC1F8135}{\ALIUSFontLatoLight}{12.357}{you have done pioneering work on various versions of this illusion that have}
  \ALIUSPlacedTextContent{90.951}{239.733}{418.437}{ALIUSC1F8135}{\ALIUSFontLatoLight}{12.357}{yielded a wealth of insights on the multisensory neural mechanisms underlying}
  \ALIUSPlacedTextContent{90.951}{257.013}{418.480}{ALIUSC1F8135}{\ALIUSFontLatoLight}{12.357}{bodily self-consciousness in humans. Could you describe the basic setup of}
  \ALIUSPlacedTextContent{90.951}{274.053}{418.481}{ALIUSC1F8135}{\ALIUSFontLatoLight}{12.357}{the full-body illusion, and summarize the main insights that you have derived}
  \ALIUSPlacedTextContent{90.951}{291.333}{173.322}{ALIUSC1F8135}{\ALIUSFontLatoLight}{12.357}{from your work on this illusion?}
  \ALIUSPlacedTextContent{90.951}{327.928}{418.289}{ALIUSC000000}{\ALIUSFontCormorantRegular}{13.920}{In the mentioned study, we tried to bridge the knowledge and theories about}
  \ALIUSPlacedTextContent{90.951}{347.608}{418.201}{ALIUSC000000}{\ALIUSFontCormorantRegular}{13.920}{the fundamental mechanisms of out-of-body experience (Blanke et al., 2005)}
  \ALIUSPlacedTextContent{90.951}{367.048}{418.179}{ALIUSC000000}{\ALIUSFontCormorantRegular}{13.920}{with at that time recent insights from the plasticity of the bodily self in}
  \ALIUSPlacedTextContent{90.951}{386.488}{418.291}{ALIUSC000000}{\ALIUSFontCormorantRegular}{13.920}{healthy participants using multisensory stimulation paradigms. We tried to}
  \ALIUSPlacedTextContent{90.951}{405.928}{418.257}{ALIUSC000000}{\ALIUSFontCormorantRegular}{13.920}{extend the technique used in the famous rubber hand illusion paradigm}
  \ALIUSPlacedTextContent{90.951}{425.608}{418.247}{ALIUSC000000}{\ALIUSFontCormorantRegular}{13.920}{(Botvinick \& Cohen, 1998) using 3D video-based virtual reality systems to}
  \ALIUSPlacedTextContent{90.951}{445.048}{418.194}{ALIUSC000000}{\ALIUSFontCormorantRegular}{13.920}{create more autoscopic or out-of-body-like illusions. We had great fun in the}
  \ALIUSPlacedTextContent{90.951}{464.488}{333.439}{ALIUSC000000}{\ALIUSFontCormorantRegular}{13.920}{lab exploring and trying out many versions of such illusions.}
  \ALIUSPlacedTextContent{90.951}{503.608}{418.235}{ALIUSC000000}{\ALIUSFontCormorantRegular}{13.920}{In the final setup, participants were shown their own (or an object's or a}
  \ALIUSPlacedTextContent{90.951}{523.049}{418.068}{ALIUSC000000}{\ALIUSFontCormorantRegular}{13.920}{mannequin's) body projected two meters in front of them. Then, by touching}
  \ALIUSPlacedTextContent{90.951}{542.489}{418.427}{ALIUSC000000}{\ALIUSFontCormorantRegular}{13.920}{their back while displaying it synchronously (as compared to asynchronously)}
  \ALIUSPlacedTextContent{90.951}{561.928}{418.232}{ALIUSC000000}{\ALIUSFontCormorantRegular}{13.920}{on the body projection, we enhanced self-identification with their projected}
  \ALIUSPlacedTextContent{90.951}{581.609}{321.694}{ALIUSC000000}{\ALIUSFontCormorantRegular}{13.920}{body, which also led to a change in perceived self-location.}
  \ALIUSPlacedTextContent{90.951}{620.489}{418.300}{ALIUSC000000}{\ALIUSFontCormorantRegular}{13.920}{For me personally, the most interesting finding of this study was that we not}
  \ALIUSPlacedTextContent{90.951}{639.928}{418.174}{ALIUSC000000}{\ALIUSFontCormorantRegular}{13.920}{only altered the perception of a single body part, as during the rubber hand}
  \ALIUSPlacedTextContent{90.951}{659.369}{418.267}{ALIUSC000000}{\ALIUSFontCormorantRegular}{13.920}{illusion, but that even the perception of the whole body and self can be}
  \ALIUSPlacedTextContent{90.951}{679.048}{418.276}{ALIUSC000000}{\ALIUSFontCormorantRegular}{13.920}{changed. This is interesting from a theoretical perspective but might also}
  \ALIUSPlacedTextContent{90.951}{698.489}{417.620}{ALIUSC000000}{\ALIUSFontCormorantRegular}{13.920}{spark potential therapeutic applications (see e.g. (Pamment \& Aspell, 2017)).}
  \ALIUSPlacedTextContent{90.951}{735.333}{414.986}{ALIUSC1F8135}{\ALIUSFontLatoLight}{12.357}{In a commentary from 2010, Adrian Alsmith writes the following about full-}
  \ALIUSPlacedTextContent{90.951}{752.613}{79.330}{ALIUSC1F8135}{\ALIUSFontLatoLight}{12.357}{body illusions:}
  \ALIUSPlacedTextContent{90.951}{793.143}{122.397}{ALIUSC767171}{\ALIUSFontCormorantLight}{12.000}{ALIUS Bulletin n°4 (2020)}
  \ALIUSPlacedTextContent{387.435}{793.143}{120.728}{ALIUSC767171}{\ALIUSFontCormorantLight}{12.000}{aliusresearch.org/bulletin}
\end{tikzpicture}
\null
\newpage

% Page 4
\thispagestyle{empty}
\noindent\begin{tikzpicture}[remember picture,overlay,x=1bp,y=1bp,shift={(current page.north west)}]
  \ALIUSPlacedTextContent{90.951}{36.329}{258.472}{ALIUSC767171}{\ALIUSFontCormorantRegular}{13.920}{Lenggenhager – Bodily boundaries and beyond}
  \ALIUSPlacedTextContent{493.431}{36.329}{15.829}{ALIUSC000000}{\ALIUSFontCormorantRegular}{13.920}{85}
  \ALIUSPlacedTextContent{126.951}{73.173}{382.391}{ALIUSC1F8135}{\ALIUSFontLatoLight}{12.357}{"Unfortunately, the question of whether or not there are full-body}
  \ALIUSPlacedTextContent{126.951}{90.453}{382.392}{ALIUSC1F8135}{\ALIUSFontLatoLight}{12.357}{illusions is empirically under-determined, as putative full-body illusions}
  \ALIUSPlacedTextContent{126.951}{107.733}{382.336}{ALIUSC1F8135}{\ALIUSFontLatoLight}{12.357}{are difficult to isolate from illusions involving composite parts that do}
  \ALIUSPlacedTextContent{126.951}{125.013}{382.380}{ALIUSC1F8135}{\ALIUSFontLatoLight}{12.357}{not constitute a 'full' or 'whole' body. That is to say, a plausible}
  \ALIUSPlacedTextContent{126.951}{142.293}{382.313}{ALIUSC1F8135}{\ALIUSFontLatoLight}{12.357}{alternative is that only representations of the body parts directly}
  \ALIUSPlacedTextContent{126.951}{159.333}{382.395}{ALIUSC1F8135}{\ALIUSFontLatoLight}{12.357}{stimulated become subject to the experimentally induced bias, whilst}
  \ALIUSPlacedTextContent{126.951}{176.613}{382.327}{ALIUSC1F8135}{\ALIUSFontLatoLight}{12.357}{other parts remain relatively (perhaps even completely) unaffected."}
  \ALIUSPlacedTextContent{126.951}{193.893}{75.743}{ALIUSC1F8135}{\ALIUSFontLatoLight}{12.357}{(Smith, 2010)}
  \ALIUSPlacedTextContent{90.951}{228.453}{418.369}{ALIUSC1F8135}{\ALIUSFontLatoLight}{12.357}{In other words, Alsmith suggests that the so-called full-body illusion might in}
  \ALIUSPlacedTextContent{90.951}{245.733}{418.448}{ALIUSC1F8135}{\ALIUSFontLatoLight}{12.357}{fact be something like a "trunk illusion" in the classic setup. Do you think this}
  \ALIUSPlacedTextContent{90.951}{263.013}{418.447}{ALIUSC1F8135}{\ALIUSFontLatoLight}{12.357}{is plausible? How do you think this question could be empirically settled? For}
  \ALIUSPlacedTextContent{90.951}{280.293}{418.371}{ALIUSC1F8135}{\ALIUSFontLatoLight}{12.357}{example, if multiple body parts were stroked instead of just the trunk, would}
  \ALIUSPlacedTextContent{90.951}{297.333}{418.519}{ALIUSC1F8135}{\ALIUSFontLatoLight}{12.357}{you expect to see a change in the effect size of implicit measurements and}
  \ALIUSPlacedTextContent{90.951}{314.613}{120.843}{ALIUSC1F8135}{\ALIUSFontLatoLight}{12.357}{questionnaire reports?}
  \ALIUSPlacedTextContent{90.951}{343.768}{418.346}{ALIUSC000000}{\ALIUSFontCormorantRegular}{13.920}{I agree with Alsmith that with our experimental setup, we cannot}
  \ALIUSPlacedTextContent{90.951}{363.208}{418.115}{ALIUSC000000}{\ALIUSFontCormorantRegular}{13.920}{differentiate between these two alternative explanations. We asked the}
  \ALIUSPlacedTextContent{90.951}{382.888}{418.199}{ALIUSC000000}{\ALIUSFontCormorantRegular}{13.920}{participants in the self-location task to go back to "where they were standing}
  \ALIUSPlacedTextContent{90.951}{402.328}{418.204}{ALIUSC000000}{\ALIUSFontCormorantRegular}{13.920}{before". This question/measure does not allow us to localize the trunk in a}
  \ALIUSPlacedTextContent{90.951}{421.768}{418.179}{ALIUSC000000}{\ALIUSFontCormorantRegular}{13.920}{different place than let's say the arm. I think it would be rather easy to}
  \ALIUSPlacedTextContent{90.951}{441.208}{418.394}{ALIUSC000000}{\ALIUSFontCormorantRegular}{13.920}{experimentally test Alsmith's hypothesis, both by stimulating more and}
  \ALIUSPlacedTextContent{90.951}{460.888}{418.299}{ALIUSC000000}{\ALIUSFontCormorantRegular}{13.920}{different body parts (as in (Salomon et al., 2013)) or by using more}
  \ALIUSPlacedTextContent{90.951}{480.328}{418.279}{ALIUSC000000}{\ALIUSFontCormorantRegular}{13.920}{sophisticated self-location measures in virtual reality that could be easily}
  \ALIUSPlacedTextContent{90.951}{499.769}{418.140}{ALIUSC000000}{\ALIUSFontCormorantRegular}{13.920}{adapted to different body parts (Nakul et al., 2020). However, I do not think}
  \ALIUSPlacedTextContent{90.951}{519.208}{418.419}{ALIUSC000000}{\ALIUSFontCormorantRegular}{13.920}{that his alternative explanation is very likely. Firstly, none of the participants}
  \ALIUSPlacedTextContent{90.951}{538.889}{418.263}{ALIUSC000000}{\ALIUSFontCormorantRegular}{13.920}{has ever reported a sense of disruption or fragmentation of the body (which}
  \ALIUSPlacedTextContent{90.951}{558.328}{418.143}{ALIUSC000000}{\ALIUSFontCormorantRegular}{13.920}{would, of course, be very interesting but does not seem very common in}
  \ALIUSPlacedTextContent{90.951}{577.768}{418.168}{ALIUSC000000}{\ALIUSFontCormorantRegular}{13.920}{clinical conditions either). Secondly, a study of Ehrsson's group (Gentile et}
  \ALIUSPlacedTextContent{90.951}{597.208}{418.254}{ALIUSC000000}{\ALIUSFontCormorantRegular}{13.920}{al., 2015) looked at generalizability of ownership from the stroked body part}
  \ALIUSPlacedTextContent{90.951}{616.648}{418.275}{ALIUSC000000}{\ALIUSFontCormorantRegular}{13.920}{to the full body. Even if their setup was slightly different – with a full body}
  \ALIUSPlacedTextContent{90.951}{636.328}{418.291}{ALIUSC000000}{\ALIUSFontCormorantRegular}{13.920}{illusion from the first-person perspective – the results clearly show that the}
  \ALIUSPlacedTextContent{90.951}{655.768}{331.102}{ALIUSC000000}{\ALIUSFontCormorantRegular}{13.920}{sense of ownership spread to the non-stimulated body parts.}
  \ALIUSPlacedTextContent{108.991}{682.599}{14.160}{ALIUSC595959}{\ALIUSFontCormorantLight}{48.000}{?}
  \ALIUSPlacedTextContent{137.071}{693.307}{314.880}{ALIUSC000000}{\ALIUSFontLatoLight}{15.120}{We not only altered the perception of a single}
  \ALIUSPlacedTextContent{137.071}{711.307}{314.870}{ALIUSC000000}{\ALIUSFontLatoLight}{15.120}{body part, as during the rubber hand illusion, but}
  \ALIUSPlacedTextContent{137.071}{729.307}{314.886}{ALIUSC000000}{\ALIUSFontLatoLight}{15.120}{even the perception of the whole body and self}
  \ALIUSPlacedTextContent{462.271}{733.719}{14.160}{ALIUSC595959}{\ALIUSFontCormorantLight}{48.000}{?}
  \ALIUSPlacedTextContent{137.071}{747.307}{105.051}{ALIUSC000000}{\ALIUSFontLatoLight}{15.120}{can be changed.}
  \ALIUSPlacedTextContent{90.951}{793.143}{122.397}{ALIUSC767171}{\ALIUSFontCormorantLight}{12.000}{ALIUS Bulletin n°4 (2020)}
  \ALIUSPlacedTextContent{387.435}{793.143}{120.728}{ALIUSC767171}{\ALIUSFontCormorantLight}{12.000}{aliusresearch.org/bulletin}
\end{tikzpicture}
\null
\newpage

% Page 5
\thispagestyle{empty}
\noindent\begin{tikzpicture}[remember picture,overlay,x=1bp,y=1bp,shift={(current page.north west)}]
  \ALIUSPlacedTextContent{90.951}{36.329}{258.472}{ALIUSC767171}{\ALIUSFontCormorantRegular}{13.920}{Lenggenhager – Bodily boundaries and beyond}
  \ALIUSPlacedTextContent{492.471}{36.329}{16.613}{ALIUSC000000}{\ALIUSFontCormorantRegular}{13.920}{86}
  \ALIUSPlacedTextContent{90.951}{73.173}{418.405}{ALIUSC1F8135}{\ALIUSFontLatoLight}{12.357}{The rubber hand illusion (Botvinick \& Cohen, 1998) represents the pioneering}
  \ALIUSPlacedTextContent{90.951}{90.453}{418.347}{ALIUSC1F8135}{\ALIUSFontLatoLight}{12.357}{and most widely used experimental paradigm to temporarily alter body}
  \ALIUSPlacedTextContent{90.951}{107.733}{418.395}{ALIUSC1F8135}{\ALIUSFontLatoLight}{12.357}{ownership by transferring the sense of ownership from one's own arm to an}
  \ALIUSPlacedTextContent{90.951}{125.013}{418.374}{ALIUSC1F8135}{\ALIUSFontLatoLight}{12.357}{artificial body part. Synchronous stroking of the participant's concealed real}
  \ALIUSPlacedTextContent{90.951}{142.293}{418.340}{ALIUSC1F8135}{\ALIUSFontLatoLight}{12.357}{hand and visible rubber hand, so that the felt touch of the brush on the real}
  \ALIUSPlacedTextContent{90.951}{159.333}{418.285}{ALIUSC1F8135}{\ALIUSFontLatoLight}{12.357}{hand and the seen touch on the fake hand are closely matched, induces a}
  \ALIUSPlacedTextContent{90.951}{176.613}{418.320}{ALIUSC1F8135}{\ALIUSFontLatoLight}{12.357}{sense of ownership of the artificial limb. Broadly accepted theoretical}
  \ALIUSPlacedTextContent{90.951}{193.893}{418.387}{ALIUSC1F8135}{\ALIUSFontLatoLight}{12.357}{frameworks maintain that the rubber hand illusion reflects the role of}
  \ALIUSPlacedTextContent{90.951}{211.173}{418.387}{ALIUSC1F8135}{\ALIUSFontLatoLight}{12.357}{multimodal integration in embodiment (Suzuki et al., 2013), indicative of the}
  \ALIUSPlacedTextContent{90.951}{228.453}{400.923}{ALIUSC1F8135}{\ALIUSFontLatoLight}{12.357}{brain's attempt to rectify discrepant multisensory visuotactile information.}
  \ALIUSPlacedTextContent{90.951}{257.733}{418.420}{ALIUSC1F8135}{\ALIUSFontLatoLight}{12.357}{However, recent criticism contends that existing interpretations overlook the}
  \ALIUSPlacedTextContent{90.951}{275.013}{418.397}{ALIUSC1F8135}{\ALIUSFontLatoLight}{12.357}{role of trait differences in the ability to generate experiences that meet}
  \ALIUSPlacedTextContent{90.951}{292.293}{72.866}{ALIUSC1F8135}{\ALIUSFontLatoLight}{12.357}{expectancies}
  \ALIUSPlacedTextContent{178.440}{292.293}{104.927}{ALIUSC1F8135}{\ALIUSFontLatoLight}{12.357}{(phenomenological}
  \ALIUSPlacedTextContent{297.991}{292.293}{44.871}{ALIUSC1F8135}{\ALIUSFontLatoLight}{12.357}{control)}
  \ALIUSPlacedTextContent{357.486}{292.293}{34.779}{ALIUSC1F8135}{\ALIUSFontLatoLight}{12.357}{(Lush,}
  \ALIUSPlacedTextContent{406.889}{292.293}{38.598}{ALIUSC1F8135}{\ALIUSFontLatoLight}{12.357}{2020).}
  \ALIUSPlacedTextContent{460.111}{292.293}{49.343}{ALIUSC1F8135}{\ALIUSFontLatoLight}{12.357}{Demand}
  \ALIUSPlacedTextContent{90.951}{309.333}{418.526}{ALIUSC1F8135}{\ALIUSFontLatoLight}{12.357}{characteristics, specifically participant expectancies and stable trait}
  \ALIUSPlacedTextContent{90.951}{326.613}{418.387}{ALIUSC1F8135}{\ALIUSFontLatoLight}{12.357}{suggestibility, could account for the rubber hand illusion experiences by}
  \ALIUSPlacedTextContent{90.951}{343.893}{418.393}{ALIUSC1F8135}{\ALIUSFontLatoLight}{12.357}{"generating expectancies which are met by the voluntary top-down control of}
  \ALIUSPlacedTextContent{90.951}{361.173}{418.281}{ALIUSC1F8135}{\ALIUSFontLatoLight}{12.357}{phenomenology" (Lush, 2020, p. 1), similar to imaginative suggestibility within}
  \ALIUSPlacedTextContent{90.951}{378.453}{418.445}{ALIUSC1F8135}{\ALIUSFontLatoLight}{12.357}{the context of hypnosis. This line of argumentation rests on empirical results}
  \ALIUSPlacedTextContent{90.951}{395.733}{418.353}{ALIUSC1F8135}{\ALIUSFontLatoLight}{12.357}{evincing that measures of the rubber hand illusion are substantially related to}
  \ALIUSPlacedTextContent{90.951}{413.013}{418.412}{ALIUSC1F8135}{\ALIUSFontLatoLight}{12.357}{hypnotic and sensory suggestibility (Fiorio et al., 2020; Marotta et al., 2016;}
  \ALIUSPlacedTextContent{90.951}{430.293}{418.444}{ALIUSC1F8135}{\ALIUSFontLatoLight}{12.357}{Walsh et al., 2015), and that expectancies predict illusion scores (Lush et al.,}
  \ALIUSPlacedTextContent{90.951}{447.333}{38.575}{ALIUSC1F8135}{\ALIUSFontLatoLight}{12.357}{2019).}
  \ALIUSPlacedTextContent{90.951}{476.613}{418.432}{ALIUSC1F8135}{\ALIUSFontLatoLight}{12.357}{Do you think that measures of the rubber hand illusion may be confounded}
  \ALIUSPlacedTextContent{90.951}{493.893}{418.327}{ALIUSC1F8135}{\ALIUSFontLatoLight}{12.357}{by the active generation of phenomenological control? If a lack of control of}
  \ALIUSPlacedTextContent{90.951}{511.173}{418.381}{ALIUSC1F8135}{\ALIUSFontLatoLight}{12.357}{demand characteristics truly exists, how could rubber hand illusion measures}
  \ALIUSPlacedTextContent{90.951}{528.453}{418.295}{ALIUSC1F8135}{\ALIUSFontLatoLight}{12.357}{be improved? Would the aforementioned criticisms extend to full body}
  \ALIUSPlacedTextContent{90.951}{545.733}{418.538}{ALIUSC1F8135}{\ALIUSFontLatoLight}{12.357}{illusions as well, particularly those completed in virtual reality, which would}
  \ALIUSPlacedTextContent{90.951}{563.013}{418.315}{ALIUSC1F8135}{\ALIUSFontLatoLight}{12.357}{not entail contradictory proprioceptive information between the real and}
  \ALIUSPlacedTextContent{90.951}{580.293}{106.429}{ALIUSC1F8135}{\ALIUSFontLatoLight}{12.357}{artificial body part?}
  \ALIUSPlacedTextContent{90.951}{609.208}{418.108}{ALIUSC000000}{\ALIUSFontCormorantRegular}{13.920}{Yes, I think it is certainly worth it and important to follow up this line of}
  \ALIUSPlacedTextContent{90.951}{628.889}{418.181}{ALIUSC000000}{\ALIUSFontCormorantRegular}{13.920}{research and alternative explanations in more detail and with solid}
  \ALIUSPlacedTextContent{90.951}{648.328}{418.333}{ALIUSC000000}{\ALIUSFontCormorantRegular}{13.920}{experimental methods. Other than in classical visual illusions, like the}
  \ALIUSPlacedTextContent{90.951}{667.768}{418.230}{ALIUSC000000}{\ALIUSFontCormorantRegular}{13.920}{Müller-Lyer illusion (Müller-Lyer, 1889), which are perceived by everyone}
  \ALIUSPlacedTextContent{90.951}{687.208}{418.254}{ALIUSC000000}{\ALIUSFontCormorantRegular}{13.920}{and extremely stable between participants, in bodily illusions we typically}
  \ALIUSPlacedTextContent{90.951}{706.889}{418.292}{ALIUSC000000}{\ALIUSFontCormorantRegular}{13.920}{observe strong individual differences in how they are perceived and how}
  \ALIUSPlacedTextContent{90.951}{726.328}{418.268}{ALIUSC000000}{\ALIUSFontCormorantRegular}{13.920}{strongly people react to them. On the other hand, when you observe how}
  \ALIUSPlacedTextContent{90.951}{745.768}{418.258}{ALIUSC000000}{\ALIUSFontCormorantRegular}{13.920}{people react to walking over a virtual plank in the air, like in Vive's plank}
  \ALIUSPlacedTextContent{90.951}{793.143}{122.397}{ALIUSC767171}{\ALIUSFontCormorantLight}{12.000}{ALIUS Bulletin n°4 (2020)}
  \ALIUSPlacedTextContent{387.435}{793.143}{120.728}{ALIUSC767171}{\ALIUSFontCormorantLight}{12.000}{aliusresearch.org/bulletin}
\end{tikzpicture}
\null
\newpage

% Page 6
\thispagestyle{empty}
\noindent\begin{tikzpicture}[remember picture,overlay,x=1bp,y=1bp,shift={(current page.north west)}]
  \ALIUSPlacedTextContent{90.951}{36.329}{258.472}{ALIUSC767171}{\ALIUSFontCormorantRegular}{13.920}{Lenggenhager – Bodily boundaries and beyond}
  \ALIUSPlacedTextContent{493.191}{36.329}{16.109}{ALIUSC000000}{\ALIUSFontCormorantRegular}{13.920}{87}
  \ALIUSPlacedTextContent{90.951}{73.048}{418.314}{ALIUSC000000}{\ALIUSFontCormorantRegular}{13.920}{simulation, you can tell that even if they have insight into the experimental}
  \ALIUSPlacedTextContent{90.951}{92.488}{418.342}{ALIUSC000000}{\ALIUSFontCormorantRegular}{13.920}{setup, the body reacts as if they didn't. I agree that individual differences}
  \ALIUSPlacedTextContent{90.951}{111.928}{418.290}{ALIUSC000000}{\ALIUSFontCormorantRegular}{13.920}{could partially be caused by differences in suggestibility, and I find it}
  \ALIUSPlacedTextContent{90.951}{131.609}{418.274}{ALIUSC000000}{\ALIUSFontCormorantRegular}{13.920}{surprising myself how little research directly addresses why people react so}
  \ALIUSPlacedTextContent{90.951}{151.049}{418.321}{ALIUSC000000}{\ALIUSFontCormorantRegular}{13.920}{differently and what potential underlying mechanisms might be at work. I}
  \ALIUSPlacedTextContent{90.951}{170.488}{418.244}{ALIUSC000000}{\ALIUSFontCormorantRegular}{13.920}{think it is generally important to put more effort in understanding}
  \ALIUSPlacedTextContent{90.951}{189.928}{418.235}{ALIUSC000000}{\ALIUSFontCormorantRegular}{13.920}{differences between various individuals and also between various settings,}
  \ALIUSPlacedTextContent{90.951}{209.609}{418.160}{ALIUSC000000}{\ALIUSFontCormorantRegular}{13.920}{but do not think that demand characteristics can explain everything.}
  \ALIUSPlacedTextContent{90.951}{229.049}{418.301}{ALIUSC000000}{\ALIUSFontCormorantRegular}{13.920}{Sebastian Dieguez recently suggested, in a critical comment to one of our}
  \ALIUSPlacedTextContent{90.951}{248.489}{418.188}{ALIUSC000000}{\ALIUSFontCormorantRegular}{13.920}{papers, to call bodily illusions rather "aliefs" (Dieguez, 2018). He suggests that}
  \ALIUSPlacedTextContent{90.951}{267.928}{418.360}{ALIUSC000000}{\ALIUSFontCormorantRegular}{13.920}{such aliefs depend on motivation and expectancy of the participants, which}
  \ALIUSPlacedTextContent{90.951}{287.608}{418.129}{ALIUSC000000}{\ALIUSFontCormorantRegular}{13.920}{is very much in line with Lush's findings. But I think, even if future empirical}
  \ALIUSPlacedTextContent{90.951}{307.048}{418.261}{ALIUSC000000}{\ALIUSFontCormorantRegular}{13.920}{work would indeed show that demand characteristics are the or an}
  \ALIUSPlacedTextContent{90.951}{326.488}{418.224}{ALIUSC000000}{\ALIUSFontCormorantRegular}{13.920}{explanation of bodily illusions, it does not make the findings of the bodily}
  \ALIUSPlacedTextContent{90.951}{345.928}{418.264}{ALIUSC000000}{\ALIUSFontCormorantRegular}{13.920}{self literature less interesting (Roel Lesur et al., 2018). Over the last 20 years,}
  \ALIUSPlacedTextContent{90.951}{365.608}{418.215}{ALIUSC000000}{\ALIUSFontCormorantRegular}{13.920}{a vast amount of studies have shown that with bodily illusions, many}
  \ALIUSPlacedTextContent{90.951}{385.048}{418.324}{ALIUSC000000}{\ALIUSFontCormorantRegular}{13.920}{different physiological, emotional, and cognitive aspects can be altered. For}
  \ALIUSPlacedTextContent{90.951}{404.488}{418.271}{ALIUSC000000}{\ALIUSFontCormorantRegular}{13.920}{example, hippocampal activity underlying autobiographical memory}
  \ALIUSPlacedTextContent{90.951}{423.928}{418.232}{ALIUSC000000}{\ALIUSFontCormorantRegular}{13.920}{(Bergouignan et al., 2014), implicit peripersonal space measures (Noel et al.,}
  \ALIUSPlacedTextContent{90.951}{443.608}{418.202}{ALIUSC000000}{\ALIUSFontCormorantRegular}{13.920}{2015), and immunological response (Barnsley et al., 2011) have shown to be}
  \ALIUSPlacedTextContent{90.951}{463.048}{418.029}{ALIUSC000000}{\ALIUSFontCormorantRegular}{13.920}{modulated by bodily illusions. If these processes can all be modified by}
  \ALIUSPlacedTextContent{90.951}{482.488}{418.229}{ALIUSC000000}{\ALIUSFontCormorantRegular}{13.920}{phenomenological control in the predicted way, this would seem quite}
  \ALIUSPlacedTextContent{90.951}{501.928}{418.258}{ALIUSC000000}{\ALIUSFontCormorantRegular}{13.920}{impressive to me. Together with you, Jasmine, we are currently investigating}
  \ALIUSPlacedTextContent{90.951}{521.369}{418.322}{ALIUSC000000}{\ALIUSFontCormorantRegular}{13.920}{the link between embodiment, illusory embodiment, and placebo response,}
  \ALIUSPlacedTextContent{90.951}{541.049}{418.291}{ALIUSC000000}{\ALIUSFontCormorantRegular}{13.920}{which I think is a promising research path to continue. A further aspect that}
  \ALIUSPlacedTextContent{90.951}{560.489}{418.400}{ALIUSC000000}{\ALIUSFontCormorantRegular}{13.920}{I think might be difficult to explain with phenomenological control are the}
  \ALIUSPlacedTextContent{90.951}{579.928}{418.226}{ALIUSC000000}{\ALIUSFontCormorantRegular}{13.920}{increasing findings in animals on "bodily illusion", suggesting that there are}
  \ALIUSPlacedTextContent{90.951}{599.369}{418.210}{ALIUSC000000}{\ALIUSFontCormorantRegular}{13.920}{physiological changes specifically to the illusion condition (e.g., Shokur et al.,}
  \ALIUSPlacedTextContent{90.951}{619.048}{418.302}{ALIUSC000000}{\ALIUSFontCormorantRegular}{13.920}{2013; Wada et al., 2016). This is, by the way, another line of research that I}
  \ALIUSPlacedTextContent{90.951}{638.489}{418.219}{ALIUSC000000}{\ALIUSFontCormorantRegular}{13.920}{think would be exciting to follow up. For example: how does firing rate of}
  \ALIUSPlacedTextContent{90.951}{657.928}{418.327}{ALIUSC000000}{\ALIUSFontCormorantRegular}{13.920}{place or grid cells change after a full body illusion in rats? This would be a}
  \ALIUSPlacedTextContent{90.951}{677.369}{418.029}{ALIUSC000000}{\ALIUSFontCormorantRegular}{13.920}{much more implicit measure of self location than the one we used in our first}
  \ALIUSPlacedTextContent{90.951}{697.048}{158.948}{ALIUSC000000}{\ALIUSFontCormorantRegular}{13.920}{study you mentioned above.}
  \ALIUSPlacedTextContent{90.951}{728.489}{418.208}{ALIUSC000000}{\ALIUSFontCormorantRegular}{13.920}{Regarding the question on how we can improve current research, I think it}
  \ALIUSPlacedTextContent{90.951}{747.929}{418.238}{ALIUSC000000}{\ALIUSFontCormorantRegular}{13.920}{would be important to develop more implicit measures, but also (try) to}
  \ALIUSPlacedTextContent{90.951}{793.143}{122.397}{ALIUSC767171}{\ALIUSFontCormorantLight}{12.000}{ALIUS Bulletin n°4 (2020)}
  \ALIUSPlacedTextContent{387.435}{793.143}{120.728}{ALIUSC767171}{\ALIUSFontCormorantLight}{12.000}{aliusresearch.org/bulletin}
\end{tikzpicture}
\null
\newpage

% Page 7
\thispagestyle{empty}
\noindent\begin{tikzpicture}[remember picture,overlay,x=1bp,y=1bp,shift={(current page.north west)}]
  \ALIUSPlacedTextContent{90.951}{36.329}{258.472}{ALIUSC767171}{\ALIUSFontCormorantRegular}{13.920}{Lenggenhager – Bodily boundaries and beyond}
  \ALIUSPlacedTextContent{492.231}{36.329}{16.949}{ALIUSC000000}{\ALIUSFontCormorantRegular}{13.920}{88}
  \ALIUSPlacedTextContent{90.951}{73.048}{418.206}{ALIUSC000000}{\ALIUSFontCormorantRegular}{13.920}{replicate the many interesting implicit measures in humans and animals that}
  \ALIUSPlacedTextContent{90.951}{92.488}{418.257}{ALIUSC000000}{\ALIUSFontCormorantRegular}{13.920}{have already been described, ideally in much bigger sample sizes. This is}
  \ALIUSPlacedTextContent{90.951}{111.928}{418.310}{ALIUSC000000}{\ALIUSFontCormorantRegular}{13.920}{important, as some of these beautiful data have not yet or only rarely been}
  \ALIUSPlacedTextContent{90.951}{131.609}{418.221}{ALIUSC000000}{\ALIUSFontCormorantRegular}{13.920}{replicated. Bigger sample size would facilitate the assessment of individual}
  \ALIUSPlacedTextContent{90.951}{151.049}{418.250}{ALIUSC000000}{\ALIUSFontCormorantRegular}{13.920}{differences and to link different implicit measures with phenomenological}
  \ALIUSPlacedTextContent{90.951}{170.488}{51.808}{ALIUSC000000}{\ALIUSFontCormorantRegular}{13.920}{changes.}
  \ALIUSPlacedTextContent{90.951}{202.053}{418.408}{ALIUSC1F8135}{\ALIUSFontLatoLight}{12.357}{To follow up on the points above, the rubber hand illusion is classically}
  \ALIUSPlacedTextContent{90.951}{219.333}{418.469}{ALIUSC1F8135}{\ALIUSFontLatoLight}{12.357}{considered a measure of body ownership. Yet, accumulating evidence}
  \ALIUSPlacedTextContent{90.951}{236.613}{418.445}{ALIUSC1F8135}{\ALIUSFontLatoLight}{12.357}{suggests that body ownership and self-location constitute distinct aspects of}
  \ALIUSPlacedTextContent{90.951}{253.893}{418.397}{ALIUSC1F8135}{\ALIUSFontLatoLight}{12.357}{bodily self-consciousness (Serino et al., 2013). While explicit measures of}
  \ALIUSPlacedTextContent{90.951}{271.173}{418.409}{ALIUSC1F8135}{\ALIUSFontLatoLight}{12.357}{body ownership are assessed by means of a self-report questionnaire,}
  \ALIUSPlacedTextContent{90.951}{288.453}{418.291}{ALIUSC1F8135}{\ALIUSFontLatoLight}{12.357}{proprioceptive drift towards the rubber hand is generally considered an}
  \ALIUSPlacedTextContent{90.951}{305.733}{418.436}{ALIUSC1F8135}{\ALIUSFontLatoLight}{12.357}{implicit measure of body ownership. However, if proprioception measures the}
  \ALIUSPlacedTextContent{90.951}{323.013}{418.408}{ALIUSC1F8135}{\ALIUSFontLatoLight}{12.357}{shift of the own arm towards the rubber hand, would implicit proprioceptive}
  \ALIUSPlacedTextContent{90.951}{340.053}{418.429}{ALIUSC1F8135}{\ALIUSFontLatoLight}{12.357}{drift not more accurately constitute a representation of self-location rather}
  \ALIUSPlacedTextContent{90.951}{357.333}{418.405}{ALIUSC1F8135}{\ALIUSFontLatoLight}{12.357}{than ownership? Considering that these implicit and explicit measures of the}
  \ALIUSPlacedTextContent{90.951}{374.613}{418.388}{ALIUSC1F8135}{\ALIUSFontLatoLight}{12.357}{rubber hand illusion further seem to be substantially yet differentially related}
  \ALIUSPlacedTextContent{90.951}{391.893}{418.392}{ALIUSC1F8135}{\ALIUSFontLatoLight}{12.357}{to hypnotic and sensory suggestibility, would a reform or at least more specific}
  \ALIUSPlacedTextContent{90.951}{409.173}{383.874}{ALIUSC1F8135}{\ALIUSFontLatoLight}{12.357}{operationalization of the rubber hand illusion measures be warranted?}
  \ALIUSPlacedTextContent{90.951}{438.328}{418.247}{ALIUSC000000}{\ALIUSFontCormorantRegular}{13.920}{Yes, I would agree. A measure that might be more directly linked to the sense}
  \ALIUSPlacedTextContent{90.951}{457.768}{418.201}{ALIUSC000000}{\ALIUSFontCormorantRegular}{13.920}{of body ownership than the proprioceptive drift might be, at least}
  \ALIUSPlacedTextContent{90.951}{477.208}{418.243}{ALIUSC000000}{\ALIUSFontCormorantRegular}{13.920}{conceptually, the electrodermal activity to threat (Armel \& Ramachandran,}
  \ALIUSPlacedTextContent{90.951}{496.888}{418.271}{ALIUSC000000}{\ALIUSFontCormorantRegular}{13.920}{2003). However, there are also some problems with this method and we often}
  \ALIUSPlacedTextContent{90.951}{516.328}{418.261}{ALIUSC000000}{\ALIUSFontCormorantRegular}{13.920}{find diverging results between threat measures and self-report questionnaires}
  \ALIUSPlacedTextContent{90.951}{535.768}{418.215}{ALIUSC000000}{\ALIUSFontCormorantRegular}{13.920}{(e.g., Roel Lesur, Weijs, et al., 2020). As mentioned above, I think it will be}
  \ALIUSPlacedTextContent{90.951}{555.208}{418.259}{ALIUSC000000}{\ALIUSFontCormorantRegular}{13.920}{important to relate the different measures of the rubber hand illusion to}
  \ALIUSPlacedTextContent{90.951}{574.648}{418.281}{ALIUSC000000}{\ALIUSFontCormorantRegular}{13.920}{different phenomenological measures and also to potential confounding}
  \ALIUSPlacedTextContent{90.951}{594.328}{418.323}{ALIUSC000000}{\ALIUSFontCormorantRegular}{13.920}{differences (e.g., the aforementioned suggestibility). For that, I think it will}
  \ALIUSPlacedTextContent{90.951}{613.768}{418.226}{ALIUSC000000}{\ALIUSFontCormorantRegular}{13.920}{also be important to further develop the explicit measures we use, and I am}
  \ALIUSPlacedTextContent{90.951}{633.208}{418.192}{ALIUSC000000}{\ALIUSFontCormorantRegular}{13.920}{very optimistic that some of the measures you, Raphaël, apply, like}
  \ALIUSPlacedTextContent{90.951}{652.648}{418.288}{ALIUSC000000}{\ALIUSFontCormorantRegular}{13.920}{microphenomenological interviews combined with EEG based biomarkers}
  \ALIUSPlacedTextContent{90.951}{672.328}{357.422}{ALIUSC000000}{\ALIUSFontCormorantRegular}{13.920}{(Timmermann et al., 2019) might advance this research direction.}
  \ALIUSPlacedTextContent{90.951}{703.893}{418.361}{ALIUSC1F8135}{\ALIUSFontLatoLight}{12.357}{The notion of "first-person perspective" (1PP) is widely used in the scientific}
  \ALIUSPlacedTextContent{90.951}{721.173}{418.391}{ALIUSC1F8135}{\ALIUSFontLatoLight}{12.357}{literature on the full-body illusion and virtual reality (e.g., Blanke \& Metzinger,}
  \ALIUSPlacedTextContent{90.951}{738.453}{418.410}{ALIUSC1F8135}{\ALIUSFontLatoLight}{12.357}{2009). However, its definition is not always consistent across publications.}
  \ALIUSPlacedTextContent{90.951}{755.733}{140.361}{ALIUSC1F8135}{\ALIUSFontLatoLight}{12.357}{Here are a few examples:}
  \ALIUSPlacedTextContent{90.951}{793.143}{122.397}{ALIUSC767171}{\ALIUSFontCormorantLight}{12.000}{ALIUS Bulletin n°4 (2020)}
  \ALIUSPlacedTextContent{387.435}{793.143}{120.728}{ALIUSC767171}{\ALIUSFontCormorantLight}{12.000}{aliusresearch.org/bulletin}
\end{tikzpicture}
\null
\newpage

% Page 8
\thispagestyle{empty}
\noindent\begin{tikzpicture}[remember picture,overlay,x=1bp,y=1bp,shift={(current page.north west)}]
  \ALIUSPlacedTextContent{90.951}{36.329}{258.472}{ALIUSC767171}{\ALIUSFontCormorantRegular}{13.920}{Lenggenhager – Bodily boundaries and beyond}
  \ALIUSPlacedTextContent{492.471}{36.329}{16.613}{ALIUSC000000}{\ALIUSFontCormorantRegular}{13.920}{89}
  \ALIUSPlacedTextContent{126.951}{73.173}{382.373}{ALIUSC1F8135}{\ALIUSFontLatoLight}{12.357}{(a) "the feeling from where 'I' experience the world around me" (Pfeiffer}
  \ALIUSPlacedTextContent{126.951}{90.453}{293.319}{ALIUSC1F8135}{\ALIUSFontLatoLight}{12.357}{et al., 2014; see also Ionta et al., 2011; Blanke, 2012).}
  \ALIUSPlacedTextContent{126.951}{119.733}{382.429}{ALIUSC1F8135}{\ALIUSFontLatoLight}{12.357}{(b) "the point from which visual information from the environment is}
  \ALIUSPlacedTextContent{126.951}{137.013}{141.160}{ALIUSC1F8135}{\ALIUSFontLatoLight}{12.357}{gathered" (Maselli, 2015).}
  \ALIUSPlacedTextContent{126.951}{166.293}{382.403}{ALIUSC1F8135}{\ALIUSFontLatoLight}{12.357}{(c) "the experience of taking a first-person, body-centered, perspective}
  \ALIUSPlacedTextContent{126.951}{183.333}{244.027}{ALIUSC1F8135}{\ALIUSFontLatoLight}{12.357}{on [one's] environment" (Serino et al., 2013).}
  \ALIUSPlacedTextContent{126.951}{212.613}{382.457}{ALIUSC1F8135}{\ALIUSFontLatoLight}{12.357}{(d) "a purely geometrical feature of an egocentric model of reality}
  \ALIUSPlacedTextContent{126.951}{229.893}{382.362}{ALIUSC1F8135}{\ALIUSFontLatoLight}{12.357}{[which] includes a spatial frame of reference, plus a global body}
  \ALIUSPlacedTextContent{126.951}{247.173}{382.435}{ALIUSC1F8135}{\ALIUSFontLatoLight}{12.357}{representation, with a perspective originating within this body}
  \ALIUSPlacedTextContent{126.951}{264.453}{237.833}{ALIUSC1F8135}{\ALIUSFontLatoLight}{12.357}{representation" (Blanke \& Metzinger, 2009)}
  \ALIUSPlacedTextContent{126.951}{293.733}{382.377}{ALIUSC1F8135}{\ALIUSFontLatoLight}{12.357}{(e) "the experience of being a subject and of being directed at the word"}
  \ALIUSPlacedTextContent{126.951}{311.013}{117.592}{ALIUSC1F8135}{\ALIUSFontLatoLight}{12.357}{(Pfeiffer et al., 2016).}
  \ALIUSPlacedTextContent{90.951}{340.293}{418.354}{ALIUSC1F8135}{\ALIUSFontLatoLight}{12.357}{Many of these definitions (e.g., a, c and e) explicitly describe 1PP as an}
  \ALIUSPlacedTextContent{90.951}{357.333}{418.457}{ALIUSC1F8135}{\ALIUSFontLatoLight}{12.357}{experience, but some (e.g., b and d) do not. Furthermore, some definitions}
  \ALIUSPlacedTextContent{90.951}{374.613}{418.402}{ALIUSC1F8135}{\ALIUSFontLatoLight}{12.357}{include a reference to bodily representation, while others merely refer to the}
  \ALIUSPlacedTextContent{90.951}{391.893}{236.683}{ALIUSC1F8135}{\ALIUSFontLatoLight}{12.357}{information conveyed by visual experience.}
  \ALIUSPlacedTextContent{90.951}{421.173}{418.348}{ALIUSC1F8135}{\ALIUSFontLatoLight}{12.357}{Beyond these lexical discrepancies, it seems that there is a broad difference}
  \ALIUSPlacedTextContent{90.951}{438.453}{418.320}{ALIUSC1F8135}{\ALIUSFontLatoLight}{12.357}{between two ways of operationalizing the notion of 1PP. Some studies seem}
  \ALIUSPlacedTextContent{90.951}{455.733}{418.319}{ALIUSC1F8135}{\ALIUSFontLatoLight}{12.357}{to construe 1PP as the subject's awareness of the egocentric location of the}
  \ALIUSPlacedTextContent{90.951}{473.013}{418.415}{ALIUSC1F8135}{\ALIUSFontLatoLight}{12.357}{point of origin of her visuospatial spatial perspective – i.e., the location from}
  \ALIUSPlacedTextContent{90.951}{490.293}{418.360}{ALIUSC1F8135}{\ALIUSFontLatoLight}{12.357}{where the subject sees their environment. Other studies seem to construe}
  \ALIUSPlacedTextContent{90.951}{507.333}{418.502}{ALIUSC1F8135}{\ALIUSFontLatoLight}{12.357}{1PP as the subject's awareness of the orientation of their head or body with}
  \ALIUSPlacedTextContent{90.951}{524.613}{418.416}{ALIUSC1F8135}{\ALIUSFontLatoLight}{12.357}{respect to a geocentric frame of reference – e.g., whether the subject is}
  \ALIUSPlacedTextContent{90.951}{541.893}{293.082}{ALIUSC1F8135}{\ALIUSFontLatoLight}{12.357}{looking up or down with respect to gravitational cues.}
  \ALIUSPlacedTextContent{90.951}{571.173}{418.375}{ALIUSC1F8135}{\ALIUSFontLatoLight}{12.357}{Presumably, these two constructs are not equivalent. Blind individuals (or}
  \ALIUSPlacedTextContent{90.951}{588.453}{418.434}{ALIUSC1F8135}{\ALIUSFontLatoLight}{12.357}{sighted individuals wearing a blindfold) lack a visual perspective, but can}
  \ALIUSPlacedTextContent{90.951}{605.733}{418.315}{ALIUSC1F8135}{\ALIUSFontLatoLight}{12.357}{nonetheless be aware of their orientation with respect to gravitational cues.}
  \ALIUSPlacedTextContent{90.951}{623.013}{418.400}{ALIUSC1F8135}{\ALIUSFontLatoLight}{12.357}{Conversely, sighted individuals can become temporarily disoriented (for}
  \ALIUSPlacedTextContent{90.951}{640.293}{418.397}{ALIUSC1F8135}{\ALIUSFontLatoLight}{12.357}{example because of vertigo, zero G, or galvanic vestibular stimulation), while}
  \ALIUSPlacedTextContent{90.951}{657.333}{392.736}{ALIUSC1F8135}{\ALIUSFontLatoLight}{12.357}{remaining aware of the location from where they see their environment.}
  \ALIUSPlacedTextContent{90.951}{686.613}{418.504}{ALIUSC1F8135}{\ALIUSFontLatoLight}{12.357}{Do you think the notion of 1PP should be disambiguated, such that more}
  \ALIUSPlacedTextContent{90.951}{703.893}{418.360}{ALIUSC1F8135}{\ALIUSFontLatoLight}{12.357}{specific constructs can be investigated empirically, or do you think the}
  \ALIUSPlacedTextContent{90.951}{721.173}{418.241}{ALIUSC1F8135}{\ALIUSFontLatoLight}{12.357}{apparent discrepancies between construals of the notion can be reconciled in}
  \ALIUSPlacedTextContent{90.951}{738.453}{103.180}{ALIUSC1F8135}{\ALIUSFontLatoLight}{12.357}{a single construct?}
  \ALIUSPlacedTextContent{90.951}{793.143}{122.397}{ALIUSC767171}{\ALIUSFontCormorantLight}{12.000}{ALIUS Bulletin n°4 (2020)}
  \ALIUSPlacedTextContent{387.435}{793.143}{120.728}{ALIUSC767171}{\ALIUSFontCormorantLight}{12.000}{aliusresearch.org/bulletin}
\end{tikzpicture}
\null
\newpage

% Page 9
\thispagestyle{empty}
\noindent\begin{tikzpicture}[remember picture,overlay,x=1bp,y=1bp,shift={(current page.north west)}]
  \ALIUSPlacedTextContent{90.951}{36.329}{258.472}{ALIUSC767171}{\ALIUSFontCormorantRegular}{13.920}{Lenggenhager – Bodily boundaries and beyond}
  \ALIUSPlacedTextContent{492.711}{36.329}{16.445}{ALIUSC000000}{\ALIUSFontCormorantRegular}{13.920}{90}
  \ALIUSPlacedTextContent{90.951}{73.048}{418.373}{ALIUSC000000}{\ALIUSFontCormorantRegular}{13.920}{It is true that there are many different definitions and conceptions of the}
  \ALIUSPlacedTextContent{90.951}{92.488}{414.974}{ALIUSC000000}{\ALIUSFontCormorantRegular}{13.920}{first-person perspective. Defining the first-person perspective as a visuo-}
  \ALIUSPlacedTextContent{90.951}{111.928}{418.328}{ALIUSC000000}{\ALIUSFontCormorantRegular}{13.920}{spatial perspective, in my opinion, neglects our multisensory perception of}
  \ALIUSPlacedTextContent{90.951}{131.609}{418.229}{ALIUSC000000}{\ALIUSFontCormorantRegular}{13.920}{the world, the body, and ourselves. This, I think, reflects the classical very}
  \ALIUSPlacedTextContent{90.951}{151.049}{418.377}{ALIUSC000000}{\ALIUSFontCormorantRegular}{13.920}{strong dominance of vision and visual research in science. Thus, I would}
  \ALIUSPlacedTextContent{90.951}{170.488}{414.974}{ALIUSC000000}{\ALIUSFontCormorantRegular}{13.920}{argue that at the very least, it should be called and considered as an audio-}
  \ALIUSPlacedTextContent{90.951}{189.928}{418.213}{ALIUSC000000}{\ALIUSFontCormorantRegular}{13.920}{visuospatial perspective, since we clearly have an auditory egocentric}
  \ALIUSPlacedTextContent{90.951}{209.609}{418.175}{ALIUSC000000}{\ALIUSFontCormorantRegular}{13.920}{perspective as well. It can be experimentally manipulated as well, similar to}
  \ALIUSPlacedTextContent{90.951}{229.049}{418.279}{ALIUSC000000}{\ALIUSFontCormorantRegular}{13.920}{the visual one (e.g., Lesur et al., 2020) and can also be mismatched in rare}
  \ALIUSPlacedTextContent{90.951}{248.489}{418.226}{ALIUSC000000}{\ALIUSFontCormorantRegular}{13.920}{clinical conditions (Blanke et al., 2004). Furthermore, as you mention, the}
  \ALIUSPlacedTextContent{90.951}{267.928}{418.143}{ALIUSC000000}{\ALIUSFontCormorantRegular}{13.920}{gravitational perspective is also important and often neglected. But rather}
  \ALIUSPlacedTextContent{90.951}{287.608}{418.245}{ALIUSC000000}{\ALIUSFontCormorantRegular}{13.920}{than disambiguate the notion of the first-person perspective, I think it is}
  \ALIUSPlacedTextContent{90.951}{307.048}{418.249}{ALIUSC000000}{\ALIUSFontCormorantRegular}{13.920}{important to consider that in normal waking consciousness, our first-person}
  \ALIUSPlacedTextContent{90.951}{326.488}{418.115}{ALIUSC000000}{\ALIUSFontCormorantRegular}{13.920}{perspective is based on the integration of these different perspectives. I think}
  \ALIUSPlacedTextContent{90.951}{345.928}{418.287}{ALIUSC000000}{\ALIUSFontCormorantRegular}{13.920}{it is interesting and important to experimentally take them apart and}
  \ALIUSPlacedTextContent{90.951}{365.608}{418.228}{ALIUSC000000}{\ALIUSFontCormorantRegular}{13.920}{investigate how they affect various other aspects of self-consciousness, such}
  \ALIUSPlacedTextContent{90.951}{385.048}{418.228}{ALIUSC000000}{\ALIUSFontCormorantRegular}{13.920}{as for example self-location. To give an example, let's say the auditory}
  \ALIUSPlacedTextContent{90.951}{404.488}{418.302}{ALIUSC000000}{\ALIUSFontCormorantRegular}{13.920}{perspective is different from the visual, where would you feel localized? At}
  \ALIUSPlacedTextContent{90.951}{423.928}{418.293}{ALIUSC000000}{\ALIUSFontCormorantRegular}{13.920}{either one or the other perspective, alternating, bilocated, or somewhere in}
  \ALIUSPlacedTextContent{90.951}{443.608}{56.778}{ALIUSC000000}{\ALIUSFontCormorantRegular}{13.920}{between?}
  \ALIUSPlacedTextContent{90.951}{475.173}{418.367}{ALIUSC1F8135}{\ALIUSFontLatoLight}{12.357}{Similarly, one might wonder whether the notion of "self-location" – also widely}
  \ALIUSPlacedTextContent{90.951}{492.453}{418.395}{ALIUSC1F8135}{\ALIUSFontLatoLight}{12.357}{used in the scientific literature on the full-body illusion and virtual reality –}
  \ALIUSPlacedTextContent{90.951}{509.733}{418.458}{ALIUSC1F8135}{\ALIUSFontLatoLight}{12.357}{could benefit from disambiguation. Self-location is often defined as "the}
  \ALIUSPlacedTextContent{90.951}{526.773}{418.542}{ALIUSC1F8135}{\ALIUSFontLatoLight}{12.357}{experience of where I am in space" (e.g., Ionta et al., 2011; Blanke, 2012;}
  \ALIUSPlacedTextContent{90.951}{544.053}{221.527}{ALIUSC1F8135}{\ALIUSFontLatoLight}{12.357}{Blanke et al., 2015; Pfeiffer et al., 2014).}
  \ALIUSPlacedTextContent{90.951}{578.613}{418.433}{ALIUSC1F8135}{\ALIUSFontLatoLight}{12.357}{In the full-body illusion with front stroking, three notions of self-location can}
  \ALIUSPlacedTextContent{90.951}{595.893}{153.760}{ALIUSC1F8135}{\ALIUSFontLatoLight}{12.357}{potentially be distinguished:}
  \ALIUSPlacedTextContent{126.951}{625.173}{13.575}{ALIUSC1F8135}{\ALIUSFontLatoLight}{12.357}{1.}
  \ALIUSPlacedTextContent{162.951}{625.173}{346.479}{ALIUSC1F8135}{\ALIUSFontLatoLight}{12.357}{The subject's awareness of the location of the point of origin of}
  \ALIUSPlacedTextContent{126.951}{642.453}{349.906}{ALIUSC1F8135}{\ALIUSFontLatoLight}{12.357}{her visuospatial perspective in an egocentric frame of reference.}
  \ALIUSPlacedTextContent{126.951}{671.733}{13.575}{ALIUSC1F8135}{\ALIUSFontLatoLight}{12.357}{2.}
  \ALIUSPlacedTextContent{162.951}{671.733}{346.398}{ALIUSC1F8135}{\ALIUSFontLatoLight}{12.357}{The subject's awareness of the location of her own body – that}
  \ALIUSPlacedTextContent{126.951}{688.773}{355.018}{ALIUSC1F8135}{\ALIUSFontLatoLight}{12.357}{she sees in front of herself – in an egocentric frame of reference.}
  \ALIUSPlacedTextContent{126.951}{718.053}{13.575}{ALIUSC1F8135}{\ALIUSFontLatoLight}{12.357}{3.}
  \ALIUSPlacedTextContent{162.951}{718.053}{346.341}{ALIUSC1F8135}{\ALIUSFontLatoLight}{12.357}{The subject's awareness of her location in an allocentric frame}
  \ALIUSPlacedTextContent{126.951}{735.333}{382.433}{ALIUSC1F8135}{\ALIUSFontLatoLight}{12.357}{of reference (i.e., with respect to a map-like representation of her}
  \ALIUSPlacedTextContent{126.951}{752.613}{78.068}{ALIUSC1F8135}{\ALIUSFontLatoLight}{12.357}{environment).}
  \ALIUSPlacedTextContent{90.951}{793.143}{122.397}{ALIUSC767171}{\ALIUSFontCormorantLight}{12.000}{ALIUS Bulletin n°4 (2020)}
  \ALIUSPlacedTextContent{387.435}{793.143}{120.728}{ALIUSC767171}{\ALIUSFontCormorantLight}{12.000}{aliusresearch.org/bulletin}
\end{tikzpicture}
\null
\newpage

% Page 10
\thispagestyle{empty}
\noindent\begin{tikzpicture}[remember picture,overlay,x=1bp,y=1bp,shift={(current page.north west)}]
  \ALIUSPlacedTextContent{90.951}{36.329}{258.472}{ALIUSC767171}{\ALIUSFontCormorantRegular}{13.920}{Lenggenhager – Bodily boundaries and beyond}
  \ALIUSPlacedTextContent{494.871}{36.329}{14.415}{ALIUSC000000}{\ALIUSFontCormorantRegular}{13.920}{91}
  \ALIUSPlacedTextContent{90.951}{73.173}{418.422}{ALIUSC1F8135}{\ALIUSFontLatoLight}{12.357}{The first notion seems equivalent to one of the two notions of 1PP discussed}
  \ALIUSPlacedTextContent{90.951}{90.453}{418.432}{ALIUSC1F8135}{\ALIUSFontLatoLight}{12.357}{in the previous section. Consequently, there might be some contexts in which}
  \ALIUSPlacedTextContent{90.951}{107.733}{418.517}{ALIUSC1F8135}{\ALIUSFontLatoLight}{12.357}{"self-location" and "1PP" refer to a single construct. However, forward}
  \ALIUSPlacedTextContent{90.951}{125.013}{418.309}{ALIUSC1F8135}{\ALIUSFontLatoLight}{12.357}{proprioceptive drift, which is used in many studies on the FBI as an implicit}
  \ALIUSPlacedTextContent{90.951}{142.293}{418.395}{ALIUSC1F8135}{\ALIUSFontLatoLight}{12.357}{measure of self-location, appears to be more related to the third notion.}
  \ALIUSPlacedTextContent{90.951}{159.333}{418.385}{ALIUSC1F8135}{\ALIUSFontLatoLight}{12.357}{Indeed, this measure is obtained by moving participants around the room with}
  \ALIUSPlacedTextContent{90.951}{176.613}{418.233}{ALIUSC1F8135}{\ALIUSFontLatoLight}{12.357}{a blindfold, and asking them to walk back to their original location. Presumably,}
  \ALIUSPlacedTextContent{90.951}{193.893}{418.423}{ALIUSC1F8135}{\ALIUSFontLatoLight}{12.357}{this requires subjects to represent their current location and their original}
  \ALIUSPlacedTextContent{90.951}{211.173}{418.376}{ALIUSC1F8135}{\ALIUSFontLatoLight}{12.357}{location within an allocentric frame of reference, compute a path from one to}
  \ALIUSPlacedTextContent{90.951}{228.453}{418.386}{ALIUSC1F8135}{\ALIUSFontLatoLight}{12.357}{the other, and translate the relevant coordinates in an egocentric frame of}
  \ALIUSPlacedTextContent{90.951}{245.733}{265.933}{ALIUSC1F8135}{\ALIUSFontLatoLight}{12.357}{reference to carry out the adequate movements.}
  \ALIUSPlacedTextContent{90.951}{275.013}{418.403}{ALIUSC1F8135}{\ALIUSFontLatoLight}{12.357}{By contrast, it is perhaps more difficult to say which notion of self-location is}
  \ALIUSPlacedTextContent{90.951}{292.293}{418.379}{ALIUSC1F8135}{\ALIUSFontLatoLight}{12.357}{probed by the mental ball-dropping task, which is used as an alternative}
  \ALIUSPlacedTextContent{90.951}{309.333}{418.481}{ALIUSC1F8135}{\ALIUSFontLatoLight}{12.357}{implicit measurement of self-location in horizontal versions of the full body}
  \ALIUSPlacedTextContent{90.951}{326.613}{418.454}{ALIUSC1F8135}{\ALIUSFontLatoLight}{12.357}{illusion. This task involves asking subjects to imagine that they are dropping a}
  \ALIUSPlacedTextContent{90.951}{343.893}{418.381}{ALIUSC1F8135}{\ALIUSFontLatoLight}{12.357}{ball from their location, and to estimate how long the ball would take to hit}
  \ALIUSPlacedTextContent{90.951}{361.173}{418.380}{ALIUSC1F8135}{\ALIUSFontLatoLight}{12.357}{the floor, as a way to assess how elevated the subjects take their location to}
  \ALIUSPlacedTextContent{90.951}{378.453}{418.395}{ALIUSC1F8135}{\ALIUSFontLatoLight}{12.357}{be compared to ground level. In so far as one normally drops a ball with one's}
  \ALIUSPlacedTextContent{90.951}{395.733}{418.386}{ALIUSC1F8135}{\ALIUSFontLatoLight}{12.357}{hand – as opposed to one's mouth, for example – subjects might use the}
  \ALIUSPlacedTextContent{90.951}{413.013}{414.986}{ALIUSC1F8135}{\ALIUSFontLatoLight}{12.357}{second notion of self-location as their reference point for the mental ball-}
  \ALIUSPlacedTextContent{90.951}{430.293}{418.420}{ALIUSC1F8135}{\ALIUSFontLatoLight}{12.357}{dropping task. Furthermore, the computations required to perform the task}
  \ALIUSPlacedTextContent{90.951}{447.333}{418.366}{ALIUSC1F8135}{\ALIUSFontLatoLight}{12.357}{might not require translation of the relevant coordinates (that of the ball and}
  \ALIUSPlacedTextContent{90.951}{464.613}{418.240}{ALIUSC1F8135}{\ALIUSFontLatoLight}{12.357}{of the ground level) into an allocentric frame of reference. Indeed, subjects}
  \ALIUSPlacedTextContent{90.951}{481.893}{418.392}{ALIUSC1F8135}{\ALIUSFontLatoLight}{12.357}{should be able to estimate the egocentric distances of the fictional ball and of}
  \ALIUSPlacedTextContent{90.951}{499.173}{418.367}{ALIUSC1F8135}{\ALIUSFontLatoLight}{12.357}{the ground level with respect to the point of origin of their visuospatial}
  \ALIUSPlacedTextContent{90.951}{516.453}{418.515}{ALIUSC1F8135}{\ALIUSFontLatoLight}{12.357}{perspective, and carry out the mental task entirely in an egocentric frame of}
  \ALIUSPlacedTextContent{90.951}{533.733}{58.312}{ALIUSC1F8135}{\ALIUSFontLatoLight}{12.357}{reference.}
  \ALIUSPlacedTextContent{90.951}{563.013}{418.441}{ALIUSC1F8135}{\ALIUSFontLatoLight}{12.357}{Do you think these brief comments about the targets of the implicit}
  \ALIUSPlacedTextContent{90.951}{580.293}{418.349}{ALIUSC1F8135}{\ALIUSFontLatoLight}{12.357}{measurements used in the Full Body Illusion literature are plausible and}
  \ALIUSPlacedTextContent{90.951}{597.333}{418.404}{ALIUSC1F8135}{\ALIUSFontLatoLight}{12.357}{accurate? Specifically, do you think the proprioceptive drift task and the}
  \ALIUSPlacedTextContent{90.951}{614.613}{418.432}{ALIUSC1F8135}{\ALIUSFontLatoLight}{12.357}{mental ball-dropping task measure the same notion of "self-location"? If not,}
  \ALIUSPlacedTextContent{90.951}{631.893}{418.315}{ALIUSC1F8135}{\ALIUSFontLatoLight}{12.357}{do you have suggestions about how to disambiguate between different}
  \ALIUSPlacedTextContent{90.951}{649.173}{387.317}{ALIUSC1F8135}{\ALIUSFontLatoLight}{12.357}{notions of self-location in empirical studies using the Full Body Illusion?}
  \ALIUSPlacedTextContent{90.951}{678.328}{418.073}{ALIUSC000000}{\ALIUSFontCormorantRegular}{13.920}{Yes, I would agree with your interpretation of the different definitions and}
  \ALIUSPlacedTextContent{90.951}{697.768}{418.283}{ALIUSC000000}{\ALIUSFontCormorantRegular}{13.920}{measures of self-location. As mentioned above, in normal waking}
  \ALIUSPlacedTextContent{90.951}{717.208}{418.260}{ALIUSC000000}{\ALIUSFontCormorantRegular}{13.920}{consciousness, the multisensory first-person perspective and self-location are}
  \ALIUSPlacedTextContent{90.951}{736.888}{418.280}{ALIUSC000000}{\ALIUSFontCormorantRegular}{13.920}{typically co-located, and should also be comparable in the different implicit}
  \ALIUSPlacedTextContent{90.951}{793.143}{122.397}{ALIUSC767171}{\ALIUSFontCormorantLight}{12.000}{ALIUS Bulletin n°4 (2020)}
  \ALIUSPlacedTextContent{387.435}{793.143}{120.728}{ALIUSC767171}{\ALIUSFontCormorantLight}{12.000}{aliusresearch.org/bulletin}
\end{tikzpicture}
\null
\newpage

% Page 11
\thispagestyle{empty}
\noindent\begin{tikzpicture}[remember picture,overlay,x=1bp,y=1bp,shift={(current page.north west)}]
  \ALIUSPlacedTextContent{90.951}{36.329}{258.472}{ALIUSC767171}{\ALIUSFontCormorantRegular}{13.920}{Lenggenhager – Bodily boundaries and beyond}
  \ALIUSPlacedTextContent{493.671}{36.329}{15.395}{ALIUSC000000}{\ALIUSFontCormorantRegular}{13.920}{92}
  \ALIUSPlacedTextContent{90.951}{73.048}{418.268}{ALIUSC000000}{\ALIUSFontCormorantRegular}{13.920}{and explicit measurements of self-location, even if they require, as you}
  \ALIUSPlacedTextContent{90.951}{92.488}{418.101}{ALIUSC000000}{\ALIUSFontCormorantRegular}{13.920}{mention, different mental processing (note that in the walking task, memory}
  \ALIUSPlacedTextContent{90.951}{111.928}{418.262}{ALIUSC000000}{\ALIUSFontCormorantRegular}{13.920}{processes might also play a more important role). However, it might well be}
  \ALIUSPlacedTextContent{90.951}{131.609}{418.195}{ALIUSC000000}{\ALIUSFontCormorantRegular}{13.920}{that, if you modulate one or the other perspective, be it the visual, the}
  \ALIUSPlacedTextContent{90.951}{151.049}{418.154}{ALIUSC000000}{\ALIUSFontCormorantRegular}{13.920}{auditory, or the gravitational, this would affect the different measures}
  \ALIUSPlacedTextContent{90.951}{170.488}{418.231}{ALIUSC000000}{\ALIUSFontCormorantRegular}{13.920}{differently. One step in this direction would be to use different assessments}
  \ALIUSPlacedTextContent{90.951}{189.928}{418.291}{ALIUSC000000}{\ALIUSFontCormorantRegular}{13.920}{of self-location measures in the full body and other illusions, and relate them}
  \ALIUSPlacedTextContent{90.951}{209.609}{418.226}{ALIUSC000000}{\ALIUSFontCormorantRegular}{13.920}{also to the subjective perception during different illusions, which is}
  \ALIUSPlacedTextContent{90.951}{229.049}{418.264}{ALIUSC000000}{\ALIUSFontCormorantRegular}{13.920}{something we have recently started in collaboration with Christophe Lopez}
  \ALIUSPlacedTextContent{90.951}{248.489}{418.330}{ALIUSC000000}{\ALIUSFontCormorantRegular}{13.920}{(Nakul et al., 2020). As mentioned above, I am in general not yet very happy}
  \ALIUSPlacedTextContent{90.951}{267.928}{418.234}{ALIUSC000000}{\ALIUSFontCormorantRegular}{13.920}{with the self-location measures we are using. More implicit measures such as}
  \ALIUSPlacedTextContent{90.951}{287.608}{198.230}{ALIUSC000000}{\ALIUSFontCormorantRegular}{13.920}{place cell activity would be helpful.}
  \ALIUSPlacedTextContent{90.951}{319.173}{418.376}{ALIUSC1F8135}{\ALIUSFontLatoLight}{12.357}{Although the topic surrounding the self-awareness of animals remains beset}
  \ALIUSPlacedTextContent{90.951}{336.453}{414.986}{ALIUSC1F8135}{\ALIUSFontLatoLight}{12.357}{with a number of difficulties, a few studies have demonstrated that non-}
  \ALIUSPlacedTextContent{90.951}{353.733}{418.407}{ALIUSC1F8135}{\ALIUSFontLatoLight}{12.357}{human primates and mice may experience body ownership as well. These}
  \ALIUSPlacedTextContent{90.951}{370.773}{418.412}{ALIUSC1F8135}{\ALIUSFontLatoLight}{12.357}{illusions similarly utilize multisensory stimulations akin to the rubber hand}
  \ALIUSPlacedTextContent{90.951}{388.053}{418.367}{ALIUSC1F8135}{\ALIUSFontLatoLight}{12.357}{illusion in humans (i.e., (a)synchronous stimulation of a physical and artificial}
  \ALIUSPlacedTextContent{90.951}{405.333}{418.293}{ALIUSC1F8135}{\ALIUSFontLatoLight}{12.357}{arm or tail). Do you think such bodily illusions can effectively translate to}
  \ALIUSPlacedTextContent{90.951}{422.613}{418.436}{ALIUSC1F8135}{\ALIUSFontLatoLight}{12.357}{animal studies and validly allude to the presence of bodily self-consciousness}
  \ALIUSPlacedTextContent{90.951}{439.893}{418.540}{ALIUSC1F8135}{\ALIUSFontLatoLight}{12.357}{in animals? What type of paradigm would you employ if you were to conduct}
  \ALIUSPlacedTextContent{90.951}{457.173}{76.281}{ALIUSC1F8135}{\ALIUSFontLatoLight}{12.357}{such a study?}
  \ALIUSPlacedTextContent{90.951}{489.208}{418.441}{ALIUSC000000}{\ALIUSFontCormorantRegular}{13.920}{If you see how much difficulty healthy adult participants have in describing}
  \ALIUSPlacedTextContent{90.951}{506.249}{418.243}{ALIUSC000000}{\ALIUSFontCormorantRegular}{13.920}{what exactly they feel during bodily illusions, I think it is very hard to}
  \ALIUSPlacedTextContent{90.951}{523.049}{418.243}{ALIUSC000000}{\ALIUSFontCormorantRegular}{13.920}{imagine what animals would feel during such illusions and even harder to}
  \ALIUSPlacedTextContent{90.951}{540.089}{418.261}{ALIUSC000000}{\ALIUSFontCormorantRegular}{13.920}{imagine how you would measure such or quantify these feelings.}
  \ALIUSPlacedTextContent{90.951}{556.889}{418.313}{ALIUSC000000}{\ALIUSFontCormorantRegular}{13.920}{Nevertheless, as mentioned above, I think this literature is interesting and}
  \ALIUSPlacedTextContent{90.951}{573.928}{418.293}{ALIUSC000000}{\ALIUSFontCormorantRegular}{13.920}{important for various reasons (e.g., the phenomenological control). If maybe}
  \ALIUSPlacedTextContent{90.951}{590.969}{418.212}{ALIUSC000000}{\ALIUSFontCormorantRegular}{13.920}{not telling us much about bodily self-consciousness, animal studies might}
  \ALIUSPlacedTextContent{90.951}{607.768}{418.304}{ALIUSC000000}{\ALIUSFontCormorantRegular}{13.920}{certainly help to shed light on basic underlying mechanisms, especially on}
  \ALIUSPlacedTextContent{90.951}{624.808}{418.085}{ALIUSC000000}{\ALIUSFontCormorantRegular}{13.920}{multisensory weighing and integration. Next to the full body illusion in rats}
  \ALIUSPlacedTextContent{90.951}{641.849}{418.257}{ALIUSC000000}{\ALIUSFontCormorantRegular}{13.920}{that I mentioned above, I would really love to do longer term experiments in,}
  \ALIUSPlacedTextContent{90.951}{658.648}{418.320}{ALIUSC000000}{\ALIUSFontCormorantRegular}{13.920}{e.g., mice, if it is possible to do it in an ethically acceptable way, to better}
  \ALIUSPlacedTextContent{90.951}{675.688}{418.254}{ALIUSC000000}{\ALIUSFontCormorantRegular}{13.920}{understand the plasticity of such multisensory integration mechanisms,}
  \ALIUSPlacedTextContent{90.951}{692.729}{418.243}{ALIUSC000000}{\ALIUSFontCormorantRegular}{13.920}{which you cannot do in humans. What if a rat for example would, using}
  \ALIUSPlacedTextContent{90.951}{709.529}{418.216}{ALIUSC000000}{\ALIUSFontCormorantRegular}{13.920}{binaural headphones, from birth on always hear from another rat's}
  \ALIUSPlacedTextContent{90.951}{726.568}{418.201}{ALIUSC000000}{\ALIUSFontCormorantRegular}{13.920}{perspective (or from a stationary perspective in the room similar to}
  \ALIUSPlacedTextContent{90.951}{743.609}{418.240}{ALIUSC000000}{\ALIUSFontCormorantRegular}{13.920}{(Mizumoto \& Ishikawa, 2005)), while seeing and smelling from their own}
  \ALIUSPlacedTextContent{90.951}{793.143}{122.397}{ALIUSC767171}{\ALIUSFontCormorantLight}{12.000}{ALIUS Bulletin n°4 (2020)}
  \ALIUSPlacedTextContent{387.435}{793.143}{120.728}{ALIUSC767171}{\ALIUSFontCormorantLight}{12.000}{aliusresearch.org/bulletin}
\end{tikzpicture}
\null
\newpage

% Page 12
\thispagestyle{empty}
\noindent\begin{tikzpicture}[remember picture,overlay,x=1bp,y=1bp,shift={(current page.north west)}]
  \ALIUSPlacedTextContent{90.951}{36.329}{258.472}{ALIUSC767171}{\ALIUSFontCormorantRegular}{13.920}{Lenggenhager – Bodily boundaries and beyond}
  \ALIUSPlacedTextContent{493.911}{36.329}{15.213}{ALIUSC000000}{\ALIUSFontCormorantRegular}{13.920}{93}
  \ALIUSPlacedTextContent{90.951}{73.048}{418.258}{ALIUSC000000}{\ALIUSFontCormorantRegular}{13.920}{perspective? How would that influence the rat's perceived self-location, as for}
  \ALIUSPlacedTextContent{90.951}{89.848}{418.298}{ALIUSC000000}{\ALIUSFontCormorantRegular}{13.920}{example measured with a reaction to threat on various positions?}
  \ALIUSPlacedTextContent{90.951}{106.888}{418.359}{ALIUSC000000}{\ALIUSFontCormorantRegular}{13.920}{Investigating such long-term multisensory mismatching stimulation would}
  \ALIUSPlacedTextContent{90.951}{123.929}{418.107}{ALIUSC000000}{\ALIUSFontCormorantRegular}{13.920}{certainly tell us a lot about the plasticity of multisensory mechanisms}
  \ALIUSPlacedTextContent{90.951}{140.729}{145.242}{ALIUSC000000}{\ALIUSFontCormorantRegular}{13.920}{underlying the bodily self.}
  \ALIUSPlacedTextContent{90.951}{169.893}{414.986}{ALIUSC1F8135}{\ALIUSFontLatoLight}{12.357}{Several influential theories presuppose that basic, largely implicit and pre-}
  \ALIUSPlacedTextContent{90.951}{187.173}{418.468}{ALIUSC1F8135}{\ALIUSFontLatoLight}{12.357}{reflective bodily processes underlie self-consciousness, where the integration}
  \ALIUSPlacedTextContent{90.951}{204.453}{418.372}{ALIUSC1F8135}{\ALIUSFontLatoLight}{12.357}{of sensory and motor bodily signals with the self as an agent of intentional}
  \ALIUSPlacedTextContent{90.951}{221.733}{418.372}{ALIUSC1F8135}{\ALIUSFontLatoLight}{12.357}{object remains anchored in an embodied self (Blanke \& Metzinger, 2009; S.}
  \ALIUSPlacedTextContent{90.951}{239.013}{418.463}{ALIUSC1F8135}{\ALIUSFontLatoLight}{12.357}{Gallagher, 2005, 2013; Newen, 2018). On the contrary, opponents of this}
  \ALIUSPlacedTextContent{90.951}{256.293}{418.630}{ALIUSC1F8135}{\ALIUSFontLatoLight}{12.357}{claim do not consider bodily sensations or components of bodily awareness}
  \ALIUSPlacedTextContent{90.951}{273.333}{418.437}{ALIUSC1F8135}{\ALIUSFontLatoLight}{12.357}{(i.e., sense of ownership, self-location, and agency) a necessary prerequisite}
  \ALIUSPlacedTextContent{90.951}{290.613}{418.362}{ALIUSC1F8135}{\ALIUSFontLatoLight}{12.357}{for self-conscious experience, or even an essential requirement for}
  \ALIUSPlacedTextContent{90.951}{307.893}{376.230}{ALIUSC1F8135}{\ALIUSFontLatoLight}{12.357}{consciousness in general (Millière, 2020; Millière \& Metzinger, 2020).}
  \ALIUSPlacedTextContent{90.951}{337.173}{418.429}{ALIUSC1F8135}{\ALIUSFontLatoLight}{12.357}{Cases from clinical samples encompass unique aberrations from ordinary}
  \ALIUSPlacedTextContent{90.951}{354.453}{418.370}{ALIUSC1F8135}{\ALIUSFontLatoLight}{12.357}{bodily self-consciousness; for example, patients with asomatognosia suffer}
  \ALIUSPlacedTextContent{90.951}{371.733}{418.448}{ALIUSC1F8135}{\ALIUSFontLatoLight}{12.357}{from an unawareness of ownership of a part of their body (loss of ownership),}
  \ALIUSPlacedTextContent{90.951}{389.013}{418.392}{ALIUSC1F8135}{\ALIUSFontLatoLight}{12.357}{patients with tetraplegia experience paralysis that can result in the total loss}
  \ALIUSPlacedTextContent{90.951}{406.293}{418.395}{ALIUSC1F8135}{\ALIUSFontLatoLight}{12.357}{of movement and sensation in all four limbs and torso (loss of agency),}
  \ALIUSPlacedTextContent{90.951}{423.333}{418.288}{ALIUSC1F8135}{\ALIUSFontLatoLight}{12.357}{whereas out-of-body experiences from neurological origins include the feeling}
  \ALIUSPlacedTextContent{90.951}{440.613}{418.345}{ALIUSC1F8135}{\ALIUSFontLatoLight}{12.357}{of disembodiment from one's own body and viewing it from an elevated}
  \ALIUSPlacedTextContent{90.951}{457.893}{418.447}{ALIUSC1F8135}{\ALIUSFontLatoLight}{12.357}{visuospatial perspective (alterations of self-location). Individuals with body}
  \ALIUSPlacedTextContent{90.951}{475.173}{418.451}{ALIUSC1F8135}{\ALIUSFontLatoLight}{12.357}{integrity dysphoria may not feel like a limb belongs to them and desire}
  \ALIUSPlacedTextContent{90.951}{492.453}{418.262}{ALIUSC1F8135}{\ALIUSFontLatoLight}{12.357}{amputation or being paraplegic. Accumulating empirical evidence suggests}
  \ALIUSPlacedTextContent{90.951}{509.733}{418.380}{ALIUSC1F8135}{\ALIUSFontLatoLight}{12.357}{that transient absence of bodily experiences or weakly embodied states are}
  \ALIUSPlacedTextContent{90.951}{527.013}{418.375}{ALIUSC1F8135}{\ALIUSFontLatoLight}{12.357}{also encountered in non-clinical populations. Examples of such states include}
  \ALIUSPlacedTextContent{90.951}{544.293}{418.336}{ALIUSC1F8135}{\ALIUSFontLatoLight}{12.357}{feelings of disembodiment or a detachment between the body and mental}
  \ALIUSPlacedTextContent{90.951}{561.333}{418.335}{ALIUSC1F8135}{\ALIUSFontLatoLight}{12.357}{processes under the influence of psychedelic substances (Preller \&}
  \ALIUSPlacedTextContent{90.951}{578.613}{418.367}{ALIUSC1F8135}{\ALIUSFontLatoLight}{12.357}{Vollenweider, 2018; Timmermann et al., 2019; Vollenweider \& Kometer,}
  \ALIUSPlacedTextContent{90.951}{595.893}{418.442}{ALIUSC1F8135}{\ALIUSFontLatoLight}{12.357}{2010) or a loss of spatiotemporal awareness in dreams (Occhionero \&}
  \ALIUSPlacedTextContent{90.951}{613.173}{88.168}{ALIUSC1F8135}{\ALIUSFontLatoLight}{12.357}{Cicogna, 2011).}
  \ALIUSPlacedTextContent{90.951}{642.453}{418.526}{ALIUSC1F8135}{\ALIUSFontLatoLight}{12.357}{Do you think that such experiences contradict the claims that purport a}
  \ALIUSPlacedTextContent{90.951}{659.733}{418.450}{ALIUSC1F8135}{\ALIUSFontLatoLight}{12.357}{putative role of the body for self-consciousness? One might argue that even}
  \ALIUSPlacedTextContent{90.951}{677.013}{418.413}{ALIUSC1F8135}{\ALIUSFontLatoLight}{12.357}{in the presence of loss of agency (e.g., tetraplegia), alterations of self-location}
  \ALIUSPlacedTextContent{90.951}{694.293}{414.986}{ALIUSC1F8135}{\ALIUSFontLatoLight}{12.357}{(e.g., out-of-body experiences), or loss of self-location (reported in some drug-}
  \ALIUSPlacedTextContent{90.951}{711.333}{418.381}{ALIUSC1F8135}{\ALIUSFontLatoLight}{12.357}{induced states, in sensory deprivation and in deaf blindness, (Millière, 2019))}
  \ALIUSPlacedTextContent{90.951}{728.613}{418.377}{ALIUSC1F8135}{\ALIUSFontLatoLight}{12.357}{one could nevertheless preserve a connection to a physical body and}
  \ALIUSPlacedTextContent{90.951}{745.893}{418.517}{ALIUSC1F8135}{\ALIUSFontLatoLight}{12.357}{therefore to bodily self-consciousness. In the presence of "bodiless" dreams,}
  \ALIUSPlacedTextContent{90.951}{793.143}{122.397}{ALIUSC767171}{\ALIUSFontCormorantLight}{12.000}{ALIUS Bulletin n°4 (2020)}
  \ALIUSPlacedTextContent{387.435}{793.143}{120.728}{ALIUSC767171}{\ALIUSFontCormorantLight}{12.000}{aliusresearch.org/bulletin}
\end{tikzpicture}
\null
\newpage

% Page 13
\thispagestyle{empty}
\noindent\begin{tikzpicture}[remember picture,overlay,x=1bp,y=1bp,shift={(current page.north west)}]
  \ALIUSPlacedTextContent{90.951}{36.329}{258.472}{ALIUSC767171}{\ALIUSFontCormorantRegular}{13.920}{Lenggenhager – Bodily boundaries and beyond}
  \ALIUSPlacedTextContent{493.191}{36.329}{16.081}{ALIUSC000000}{\ALIUSFontCormorantRegular}{13.920}{94}
  \ALIUSPlacedTextContent{90.951}{73.173}{418.385}{ALIUSC1F8135}{\ALIUSFontLatoLight}{12.357}{however, one may retain a sense of agency, self-location, and first-person}
  \ALIUSPlacedTextContent{90.951}{90.453}{418.325}{ALIUSC1F8135}{\ALIUSFontLatoLight}{12.357}{perspective, despite lacking the experience of a body. Would it still make}
  \ALIUSPlacedTextContent{90.951}{107.733}{418.483}{ALIUSC1F8135}{\ALIUSFontLatoLight}{12.357}{sense to strap such states to bodily self-consciousness, or would it be more}
  \ALIUSPlacedTextContent{90.951}{125.013}{271.276}{ALIUSC1F8135}{\ALIUSFontLatoLight}{12.357}{appropriate to talk of perspectival consciousness?}
  \ALIUSPlacedTextContent{90.951}{154.169}{418.311}{ALIUSC000000}{\ALIUSFontCormorantRegular}{13.920}{I think that the mentioned examples do not contradict the claim of an}
  \ALIUSPlacedTextContent{90.951}{173.609}{418.233}{ALIUSC000000}{\ALIUSFontCormorantRegular}{13.920}{important role of the body in self-consciousness for several reasons. First of}
  \ALIUSPlacedTextContent{90.951}{193.049}{418.190}{ALIUSC000000}{\ALIUSFontCormorantRegular}{13.920}{all, I would argue that even if we find some rare instances (for example}
  \ALIUSPlacedTextContent{90.951}{212.488}{418.231}{ALIUSC000000}{\ALIUSFontCormorantRegular}{13.920}{pathological or drug-induced) in which self-consciousness might be possible}
  \ALIUSPlacedTextContent{90.951}{231.928}{418.330}{ALIUSC000000}{\ALIUSFontCormorantRegular}{13.920}{without momentary bodily experiences, it does not at all mean that bodily}
  \ALIUSPlacedTextContent{90.951}{251.609}{418.244}{ALIUSC000000}{\ALIUSFontCormorantRegular}{13.920}{self-consciousness is not generally important in self-consciousness. There is}
  \ALIUSPlacedTextContent{90.951}{271.048}{418.260}{ALIUSC000000}{\ALIUSFontCormorantRegular}{13.920}{the famous saying that the exception confirms the rule, and we know that the}
  \ALIUSPlacedTextContent{90.951}{290.488}{418.238}{ALIUSC000000}{\ALIUSFontCormorantRegular}{13.920}{brain is highly plastic and adaptive, thus it is difficult to generalize from very}
  \ALIUSPlacedTextContent{90.951}{309.928}{418.250}{ALIUSC000000}{\ALIUSFontCormorantRegular}{13.920}{specific cases to "normal" consciousness. Furthermore, all examples given are}
  \ALIUSPlacedTextContent{90.951}{329.608}{418.346}{ALIUSC000000}{\ALIUSFontCormorantRegular}{13.920}{specific states in individuals who have or at least had experienced normal}
  \ALIUSPlacedTextContent{90.951}{349.048}{37.812}{ALIUSC000000}{\ALIUSFontCormorantRegular}{13.920}{bodily}
  \ALIUSPlacedTextContent{140.886}{349.048}{80.430}{ALIUSC000000}{\ALIUSFontCormorantRegular}{13.920}{self-awareness}
  \ALIUSPlacedTextContent{233.435}{349.048}{61.915}{ALIUSC000000}{\ALIUSFontCormorantRegular}{13.920}{previously.}
  \ALIUSPlacedTextContent{307.477}{349.048}{22.193}{ALIUSC000000}{\ALIUSFontCormorantRegular}{13.920}{For}
  \ALIUSPlacedTextContent{341.795}{349.048}{51.223}{ALIUSC000000}{\ALIUSFontCormorantRegular}{13.920}{example,}
  \ALIUSPlacedTextContent{405.143}{349.048}{9.119}{ALIUSC000000}{\ALIUSFontCormorantRegular}{13.920}{a}
  \ALIUSPlacedTextContent{426.388}{349.048}{42.521}{ALIUSC000000}{\ALIUSFontCormorantRegular}{13.920}{patient}
  \ALIUSPlacedTextContent{481.035}{349.048}{28.207}{ALIUSC000000}{\ALIUSFontCormorantRegular}{13.920}{with}
  \ALIUSPlacedTextContent{90.951}{368.488}{418.304}{ALIUSC000000}{\ALIUSFontCormorantRegular}{13.920}{somatoparaphrenia might lose the sense of ownership for one hand but still}
  \ALIUSPlacedTextContent{90.951}{387.928}{418.291}{ALIUSC000000}{\ALIUSFontCormorantRegular}{13.920}{has a normal sense of ownership for the other hand. A tetraplegic might still}
  \ALIUSPlacedTextContent{90.951}{407.608}{418.262}{ALIUSC000000}{\ALIUSFontCormorantRegular}{13.920}{have a normal sense of agency over some body parts like the eye-lids. And a}
  \ALIUSPlacedTextContent{90.951}{427.048}{418.278}{ALIUSC000000}{\ALIUSFontCormorantRegular}{13.920}{dreaming person senses her body in waking state. These instances might be}
  \ALIUSPlacedTextContent{90.951}{446.488}{418.257}{ALIUSC000000}{\ALIUSFontCormorantRegular}{13.920}{an important prerequisite to be able to transiently perceive a "bodiless" state.}
  \ALIUSPlacedTextContent{90.951}{477.928}{418.235}{ALIUSC000000}{\ALIUSFontCormorantRegular}{13.920}{Thus, I would still see the lack of the sense of a body as a variation of an}
  \ALIUSPlacedTextContent{90.951}{497.608}{418.313}{ALIUSC000000}{\ALIUSFontCormorantRegular}{13.920}{altered body perception. After all, even if you wanted to call it perspectival}
  \ALIUSPlacedTextContent{90.951}{517.049}{418.171}{ALIUSC000000}{\ALIUSFontCormorantRegular}{13.920}{consciousness, a visual or auditory perspective is only experienced the way}
  \ALIUSPlacedTextContent{90.951}{536.489}{418.246}{ALIUSC000000}{\ALIUSFontCormorantRegular}{13.920}{we experience it, due to the way our body is physically shaped.}
  \ALIUSPlacedTextContent{119.311}{566.919}{14.160}{ALIUSC595959}{\ALIUSFontCormorantLight}{48.000}{?}
  \ALIUSPlacedTextContent{147.631}{580.027}{314.874}{ALIUSC000000}{\ALIUSFontLatoLight}{15.120}{I would still see the lack of the body as a variation}
  \ALIUSPlacedTextContent{472.591}{588.999}{14.160}{ALIUSC595959}{\ALIUSFontCormorantLight}{48.000}{?}
  \ALIUSPlacedTextContent{147.631}{598.027}{192.414}{ALIUSC000000}{\ALIUSFontLatoLight}{15.120}{of an altered body perception}
  \ALIUSPlacedTextContent{90.951}{629.973}{418.426}{ALIUSC1F8135}{\ALIUSFontLatoLight}{12.357}{One can certainly argue that persons with surgical or congenital amputations}
  \ALIUSPlacedTextContent{90.951}{647.253}{418.488}{ALIUSC1F8135}{\ALIUSFontLatoLight}{12.357}{retain a sense of bodily self-consciousness, despite the lack or loss of their}
  \ALIUSPlacedTextContent{90.951}{664.533}{418.439}{ALIUSC1F8135}{\ALIUSFontLatoLight}{12.357}{limbs. In the most severe cases, individuals may physically comprise no more}
  \ALIUSPlacedTextContent{90.951}{681.813}{418.370}{ALIUSC1F8135}{\ALIUSFontLatoLight}{12.357}{than their trunk and head. Yet, the fundamental sense of selfhood associated}
  \ALIUSPlacedTextContent{90.951}{699.093}{414.986}{ALIUSC1F8135}{\ALIUSFontLatoLight}{12.357}{with bodily self-consciousness is ostensibly experienced as a global whole-}
  \ALIUSPlacedTextContent{90.951}{716.133}{418.317}{ALIUSC1F8135}{\ALIUSFontLatoLight}{12.357}{body representation, rather than individual representation of separate body}
  \ALIUSPlacedTextContent{90.951}{733.413}{418.350}{ALIUSC1F8135}{\ALIUSFontLatoLight}{12.357}{parts (Lenggenhager et al., 2007). In theory – and not considering the}
  \ALIUSPlacedTextContent{90.951}{750.693}{418.394}{ALIUSC1F8135}{\ALIUSFontLatoLight}{12.357}{constraints of vital organs – how much of the body could we progressively}
  \ALIUSPlacedTextContent{90.951}{793.143}{122.397}{ALIUSC767171}{\ALIUSFontCormorantLight}{12.000}{ALIUS Bulletin n°4 (2020)}
  \ALIUSPlacedTextContent{387.435}{793.143}{120.728}{ALIUSC767171}{\ALIUSFontCormorantLight}{12.000}{aliusresearch.org/bulletin}
\end{tikzpicture}
\null
\newpage

% Page 14
\thispagestyle{empty}
\noindent\begin{tikzpicture}[remember picture,overlay,x=1bp,y=1bp,shift={(current page.north west)}]
  \ALIUSPlacedTextContent{90.951}{36.329}{258.472}{ALIUSC767171}{\ALIUSFontCormorantRegular}{13.920}{Lenggenhager – Bodily boundaries and beyond}
  \ALIUSPlacedTextContent{493.671}{36.329}{15.493}{ALIUSC000000}{\ALIUSFontCormorantRegular}{13.920}{95}
  \ALIUSPlacedTextContent{90.951}{73.173}{418.467}{ALIUSC1F8135}{\ALIUSFontLatoLight}{12.357}{remove before the sense of (bodily) self-consciousness is lost? Along these}
  \ALIUSPlacedTextContent{90.951}{90.453}{393.494}{ALIUSC1F8135}{\ALIUSFontLatoLight}{12.357}{lines, could prosthetics serve as a viable substitute for the physical body?}
  \ALIUSPlacedTextContent{90.951}{119.608}{418.127}{ALIUSC000000}{\ALIUSFontCormorantRegular}{13.920}{This is an intriguing question. Similar to the famous ship of Theseus example,}
  \ALIUSPlacedTextContent{90.951}{139.049}{418.249}{ALIUSC000000}{\ALIUSFontCormorantRegular}{13.920}{which questions whether a ship would still be the same if all parts were}
  \ALIUSPlacedTextContent{90.951}{158.489}{418.246}{ALIUSC000000}{\ALIUSFontCormorantRegular}{13.920}{successively replaced over time. This example would be even more interesting}
  \ALIUSPlacedTextContent{90.951}{177.928}{418.230}{ALIUSC000000}{\ALIUSFontCormorantRegular}{13.920}{for a human being than for a ship, what if you replaced one organ/body part}
  \ALIUSPlacedTextContent{90.951}{197.609}{418.159}{ALIUSC000000}{\ALIUSFontCormorantRegular}{13.920}{after the other in a human being, would she still be the same? While this}
  \ALIUSPlacedTextContent{90.951}{217.049}{414.974}{ALIUSC000000}{\ALIUSFontCormorantRegular}{13.920}{remains thus far a thought experiment, the question on how self-}
  \ALIUSPlacedTextContent{90.951}{236.488}{418.281}{ALIUSC000000}{\ALIUSFontCormorantRegular}{13.920}{consciousness is altered by, e.g., organ transplantation or prosthetics, is a very}
  \ALIUSPlacedTextContent{90.951}{255.929}{418.168}{ALIUSC000000}{\ALIUSFontCormorantRegular}{13.920}{interesting and important one, which should be carefully investigated. I think}
  \ALIUSPlacedTextContent{90.951}{275.608}{418.279}{ALIUSC000000}{\ALIUSFontCormorantRegular}{13.920}{the "problem" in your question might exactly be the "not considering the}
  \ALIUSPlacedTextContent{90.951}{295.048}{418.251}{ALIUSC000000}{\ALIUSFontCormorantRegular}{13.920}{constraints of vital organs", as both empirical and theoretical evidence}
  \ALIUSPlacedTextContent{90.951}{314.488}{418.274}{ALIUSC000000}{\ALIUSFontCormorantRegular}{13.920}{suggest that the trunk and head, where also the vital organs are located, is the}
  \ALIUSPlacedTextContent{90.951}{333.928}{418.377}{ALIUSC000000}{\ALIUSFontCormorantRegular}{13.920}{core of this fundamental selfhood, which is also where people typically}
  \ALIUSPlacedTextContent{90.951}{353.608}{418.261}{ALIUSC000000}{\ALIUSFontCormorantRegular}{13.920}{localize their "self" if forced to localized it in a single point (Alsmith \& Longo,}
  \ALIUSPlacedTextContent{90.951}{373.048}{418.273}{ALIUSC000000}{\ALIUSFontCormorantRegular}{13.920}{2014). This also fits the idea that interoceptive cues, mainly from the trunk,}
  \ALIUSPlacedTextContent{90.951}{392.488}{418.101}{ALIUSC000000}{\ALIUSFontCormorantRegular}{13.920}{might crucially underlie our sense of a bodily self (e.g., Park \& Blanke, 2019).}
  \ALIUSPlacedTextContent{90.951}{411.928}{418.243}{ALIUSC000000}{\ALIUSFontCormorantRegular}{13.920}{But in a pure thought experiment, my guess would be that as long as you have}
  \ALIUSPlacedTextContent{90.951}{431.608}{418.265}{ALIUSC000000}{\ALIUSFontCormorantRegular}{13.920}{some physical body (even a prosthetic one if it were integrated in your}
  \ALIUSPlacedTextContent{90.951}{451.048}{418.352}{ALIUSC000000}{\ALIUSFontCormorantRegular}{13.920}{sensorimotor loop) you would feel as physically embodied. And again, I}
  \ALIUSPlacedTextContent{90.951}{470.488}{418.146}{ALIUSC000000}{\ALIUSFontCormorantRegular}{13.920}{would argue that there is an important difference between whether you}
  \ALIUSPlacedTextContent{90.951}{489.928}{418.289}{ALIUSC000000}{\ALIUSFontCormorantRegular}{13.920}{previously had a body and lose/replace it as in your thought experiment, or}
  \ALIUSPlacedTextContent{90.951}{509.608}{294.885}{ALIUSC000000}{\ALIUSFontCormorantRegular}{13.920}{whether there would be no body from the beginning.}
  \ALIUSPlacedTextContent{90.951}{548.613}{418.215}{ALIUSC1F8135}{\ALIUSFontLatoLight}{12.357}{Several clinical trials using psychedelic substances have evinced impressive}
  \ALIUSPlacedTextContent{90.951}{565.893}{418.304}{ALIUSC1F8135}{\ALIUSFontLatoLight}{12.357}{and often persistent improvements in depression. During such sessions,}
  \ALIUSPlacedTextContent{90.951}{583.173}{418.316}{ALIUSC1F8135}{\ALIUSFontLatoLight}{12.357}{perceived disembodiment or detachment from the physical body constitute}
  \ALIUSPlacedTextContent{90.951}{600.453}{418.310}{ALIUSC1F8135}{\ALIUSFontLatoLight}{12.357}{common phenomenological experiences under the influence of psychedelic}
  \ALIUSPlacedTextContent{90.951}{617.733}{418.380}{ALIUSC1F8135}{\ALIUSFontLatoLight}{12.357}{substances (Belser et al., 2017; Watts et al., 2017). In a separate line of}
  \ALIUSPlacedTextContent{90.951}{634.773}{418.378}{ALIUSC1F8135}{\ALIUSFontLatoLight}{12.357}{research, an experimentally induced out-of-body illusion successfully}
  \ALIUSPlacedTextContent{90.951}{652.053}{418.440}{ALIUSC1F8135}{\ALIUSFontLatoLight}{12.357}{alleviated symptoms in a number of chronic pain conditions, including}
  \ALIUSPlacedTextContent{90.951}{669.333}{418.442}{ALIUSC1F8135}{\ALIUSFontLatoLight}{12.357}{fibromyalgia, endometriosis, chronic lower back pain, and spinal cord injury}
  \ALIUSPlacedTextContent{90.951}{686.613}{414.986}{ALIUSC1F8135}{\ALIUSFontLatoLight}{12.357}{(Pamment \& Aspell, 2017). Do you think that altered states of bodily self-}
  \ALIUSPlacedTextContent{90.951}{703.893}{418.562}{ALIUSC1F8135}{\ALIUSFontLatoLight}{12.357}{consciousness, specifically those related to disembodiment, are central to the}
  \ALIUSPlacedTextContent{90.951}{721.173}{418.474}{ALIUSC1F8135}{\ALIUSFontLatoLight}{12.357}{alleviation of symptoms in such conditions? If so, why and how would a}
  \ALIUSPlacedTextContent{90.951}{738.453}{418.400}{ALIUSC1F8135}{\ALIUSFontLatoLight}{12.357}{perceived "detachment" from our bodies operate beneficially for our mental}
  \ALIUSPlacedTextContent{90.951}{755.733}{108.273}{ALIUSC1F8135}{\ALIUSFontLatoLight}{12.357}{and physical health?}
  \ALIUSPlacedTextContent{90.951}{793.143}{122.397}{ALIUSC767171}{\ALIUSFontCormorantLight}{12.000}{ALIUS Bulletin n°4 (2020)}
  \ALIUSPlacedTextContent{387.435}{793.143}{120.728}{ALIUSC767171}{\ALIUSFontCormorantLight}{12.000}{aliusresearch.org/bulletin}
\end{tikzpicture}
\null
\newpage

% Page 15
\thispagestyle{empty}
\noindent\begin{tikzpicture}[remember picture,overlay,x=1bp,y=1bp,shift={(current page.north west)}]
  \ALIUSPlacedTextContent{90.951}{36.329}{258.472}{ALIUSC767171}{\ALIUSFontCormorantRegular}{13.920}{Lenggenhager – Bodily boundaries and beyond}
  \ALIUSPlacedTextContent{492.951}{36.329}{16.277}{ALIUSC000000}{\ALIUSFontCormorantRegular}{13.920}{96}
  \ALIUSPlacedTextContent{90.951}{73.048}{418.184}{ALIUSC000000}{\ALIUSFontCormorantRegular}{13.920}{While the mentioned examples are very promising there is still limited}
  \ALIUSPlacedTextContent{90.951}{92.488}{418.395}{ALIUSC000000}{\ALIUSFontCormorantRegular}{13.920}{empirical evidence. Replication studies are needed. However, I think that the}
  \ALIUSPlacedTextContent{90.951}{111.928}{418.377}{ALIUSC000000}{\ALIUSFontCormorantRegular}{13.920}{experience of disembodiment, as well as possibly related sensations, such as}
  \ALIUSPlacedTextContent{90.951}{131.609}{418.235}{ALIUSC000000}{\ALIUSFontCormorantRegular}{13.920}{the sense of lightness, might be beneficial. Again, I think it would be}
  \ALIUSPlacedTextContent{90.951}{151.049}{418.165}{ALIUSC000000}{\ALIUSFontCormorantRegular}{13.920}{important to better quantify different phenomenological aspects of such}
  \ALIUSPlacedTextContent{90.951}{170.488}{418.253}{ALIUSC000000}{\ALIUSFontCormorantRegular}{13.920}{altered sense of embodiment and to investigate in the different states (e.g., in}
  \ALIUSPlacedTextContent{90.951}{189.928}{418.267}{ALIUSC000000}{\ALIUSFontCormorantRegular}{13.920}{virtual reality or during psychedelic drugs) to better understand which}
  \ALIUSPlacedTextContent{90.951}{209.609}{418.196}{ALIUSC000000}{\ALIUSFontCormorantRegular}{13.920}{aspects of such experiences might be helpful. Generally, I am very excited by}
  \ALIUSPlacedTextContent{90.951}{229.049}{418.187}{ALIUSC000000}{\ALIUSFontCormorantRegular}{13.920}{the idea that the unusual body experience itself rather than, for example, the}
  \ALIUSPlacedTextContent{90.951}{248.489}{418.229}{ALIUSC000000}{\ALIUSFontCormorantRegular}{13.920}{pure neurochemical alterations due to psychedelic substances might cause}
  \ALIUSPlacedTextContent{90.951}{267.928}{418.285}{ALIUSC000000}{\ALIUSFontCormorantRegular}{13.920}{some of these changes. If you think of the hyperembodiment model of}
  \ALIUSPlacedTextContent{90.951}{287.608}{418.235}{ALIUSC000000}{\ALIUSFontCormorantRegular}{13.920}{Thomas Fuchs, for example (Fuchs \& Schlimme, 2009), in which depression}
  \ALIUSPlacedTextContent{90.951}{307.048}{418.129}{ALIUSC000000}{\ALIUSFontCormorantRegular}{13.920}{is suggested to be linked to a too rigid and strong sense of embodiment, it}
  \ALIUSPlacedTextContent{90.951}{326.488}{418.210}{ALIUSC000000}{\ALIUSFontCormorantRegular}{13.920}{could well be that the sense of transient disembodiment or lightness might}
  \ALIUSPlacedTextContent{90.951}{345.928}{418.136}{ALIUSC000000}{\ALIUSFontCormorantRegular}{13.920}{be helpful. Furthermore, while a constant sense of detachment from the body}
  \ALIUSPlacedTextContent{90.951}{365.608}{418.171}{ALIUSC000000}{\ALIUSFontCormorantRegular}{13.920}{might not be beneficial, I think that there is a lot of therapeutic potential in}
  \ALIUSPlacedTextContent{90.951}{385.048}{418.257}{ALIUSC000000}{\ALIUSFontCormorantRegular}{13.920}{using such tools to let people perceive how plastic the sense of the body}
  \ALIUSPlacedTextContent{90.951}{404.488}{418.338}{ALIUSC000000}{\ALIUSFontCormorantRegular}{13.920}{actually is. This could also be applied in a more educative way by showing}
  \ALIUSPlacedTextContent{90.951}{423.928}{421.533}{ALIUSC000000}{\ALIUSFontCormorantRegular}{13.920}{people that the sense of their own body might not be as rigid as they believe.}
  \ALIUSPlacedTextContent{108.511}{454.839}{14.160}{ALIUSC595959}{\ALIUSFontCormorantLight}{48.000}{?}
  \ALIUSPlacedTextContent{141.151}{464.107}{314.927}{ALIUSC000000}{\ALIUSFontLatoLight}{15.120}{While a constant sense of detachment from the}
  \ALIUSPlacedTextContent{141.151}{482.107}{314.913}{ALIUSC000000}{\ALIUSFontLatoLight}{15.120}{body might not be beneficial, I think that there is}
  \ALIUSPlacedTextContent{141.151}{500.107}{314.889}{ALIUSC000000}{\ALIUSFontLatoLight}{15.120}{a lot of therapeutic potential in using such tools}
  \ALIUSPlacedTextContent{141.151}{518.107}{314.911}{ALIUSC000000}{\ALIUSFontLatoLight}{15.120}{to let people perceive how plastic the sense of}
  \ALIUSPlacedTextContent{471.631}{528.039}{14.160}{ALIUSC595959}{\ALIUSFontCormorantLight}{48.000}{?}
  \ALIUSPlacedTextContent{141.151}{536.107}{124.744}{ALIUSC000000}{\ALIUSFontLatoLight}{15.120}{the body actually is}
  \ALIUSPlacedTextContent{90.951}{577.893}{418.426}{ALIUSC1F8135}{\ALIUSFontLatoLight}{12.357}{Virtual reality has undergone a transition from an expensive and arduous}
  \ALIUSPlacedTextContent{90.951}{595.173}{418.285}{ALIUSC1F8135}{\ALIUSFontLatoLight}{12.357}{device to a functional technology that is increasingly employed both in}
  \ALIUSPlacedTextContent{90.951}{612.453}{418.395}{ALIUSC1F8135}{\ALIUSFontLatoLight}{12.357}{empirical and home-use settings. While the immersive nature of virtual reality}
  \ALIUSPlacedTextContent{90.951}{629.733}{418.377}{ALIUSC1F8135}{\ALIUSFontLatoLight}{12.357}{has already demonstrated considerable success as an application for}
  \ALIUSPlacedTextContent{90.951}{647.013}{418.456}{ALIUSC1F8135}{\ALIUSFontLatoLight}{12.357}{behavioral health (e.g., exposure therapy for phobias or posttraumatic stress}
  \ALIUSPlacedTextContent{90.951}{664.293}{418.395}{ALIUSC1F8135}{\ALIUSFontLatoLight}{12.357}{disorder) (Riva et al., 2018), a particularly auspicious feature of virtual reality}
  \ALIUSPlacedTextContent{90.951}{681.573}{418.380}{ALIUSC1F8135}{\ALIUSFontLatoLight}{12.357}{involves the implementation of virtual avatars as bottom-up sensory}
  \ALIUSPlacedTextContent{90.951}{698.613}{418.388}{ALIUSC1F8135}{\ALIUSFontLatoLight}{12.357}{modulators of existing body representations (embodied virtual reality). Virtual}
  \ALIUSPlacedTextContent{90.951}{715.893}{418.514}{ALIUSC1F8135}{\ALIUSFontLatoLight}{12.357}{reality permits a degree of sensory control that would not be possible in}
  \ALIUSPlacedTextContent{90.951}{733.173}{418.372}{ALIUSC1F8135}{\ALIUSFontLatoLight}{12.357}{physical reality, and thus facilitates unique opportunities to update aberrant}
  \ALIUSPlacedTextContent{90.951}{750.453}{244.977}{ALIUSC1F8135}{\ALIUSFontLatoLight}{12.357}{representations of bodily self-consciousness.}
  \ALIUSPlacedTextContent{90.951}{793.143}{122.397}{ALIUSC767171}{\ALIUSFontCormorantLight}{12.000}{ALIUS Bulletin n°4 (2020)}
  \ALIUSPlacedTextContent{387.435}{793.143}{120.728}{ALIUSC767171}{\ALIUSFontCormorantLight}{12.000}{aliusresearch.org/bulletin}
\end{tikzpicture}
\null
\newpage

% Page 16
\thispagestyle{empty}
\noindent\begin{tikzpicture}[remember picture,overlay,x=1bp,y=1bp,shift={(current page.north west)}]
  \ALIUSPlacedTextContent{90.951}{36.329}{258.472}{ALIUSC767171}{\ALIUSFontCormorantRegular}{13.920}{Lenggenhager – Bodily boundaries and beyond}
  \ALIUSPlacedTextContent{493.431}{36.329}{15.773}{ALIUSC000000}{\ALIUSFontCormorantRegular}{13.920}{97}
  \ALIUSPlacedTextContent{90.951}{73.173}{418.347}{ALIUSC1F8135}{\ALIUSFontLatoLight}{12.357}{With a potentially increasing recognition of embodiment in the convergence}
  \ALIUSPlacedTextContent{90.951}{90.453}{418.395}{ALIUSC1F8135}{\ALIUSFontLatoLight}{12.357}{of phenomenology and pathology, do you consider embodied virtual reality a}
  \ALIUSPlacedTextContent{90.951}{107.733}{418.485}{ALIUSC1F8135}{\ALIUSFontLatoLight}{12.357}{promising therapeutic tool for disorders of bodily self-consciousness? Is it or}
  \ALIUSPlacedTextContent{90.951}{125.013}{418.451}{ALIUSC1F8135}{\ALIUSFontLatoLight}{12.357}{will it be possible to effectively "replace" our physical bodies with virtual}
  \ALIUSPlacedTextContent{90.951}{142.293}{418.437}{ALIUSC1F8135}{\ALIUSFontLatoLight}{12.357}{selves? Given adequate sensory information, can we embody almost anything}
  \ALIUSPlacedTextContent{90.951}{159.333}{418.506}{ALIUSC1F8135}{\ALIUSFontLatoLight}{12.357}{– e.g., an animal, a superhero, or an inanimate object such as a house? Along}
  \ALIUSPlacedTextContent{90.951}{176.613}{418.397}{ALIUSC1F8135}{\ALIUSFontLatoLight}{12.357}{these lines, could we further embody multiple selves, and if so, how might this}
  \ALIUSPlacedTextContent{90.951}{193.893}{418.329}{ALIUSC1F8135}{\ALIUSFontLatoLight}{12.357}{affect our sense of (bodily) self long term? Although the potential applications}
  \ALIUSPlacedTextContent{90.951}{211.173}{418.367}{ALIUSC1F8135}{\ALIUSFontLatoLight}{12.357}{for clinical disorders seem promising, what are some of the ethical concerns}
  \ALIUSPlacedTextContent{90.951}{228.453}{418.393}{ALIUSC1F8135}{\ALIUSFontLatoLight}{12.357}{and risks that need to be considered moving forward with virtual reality and}
  \ALIUSPlacedTextContent{90.951}{245.733}{315.963}{ALIUSC1F8135}{\ALIUSFontLatoLight}{12.357}{other "body surrogates", such as robotics and prosthetics?}
  \ALIUSPlacedTextContent{90.951}{274.888}{418.249}{ALIUSC000000}{\ALIUSFontCormorantRegular}{13.920}{Yes, I definitely consider embodied virtual reality a promising therapeutic}
  \ALIUSPlacedTextContent{90.951}{294.328}{418.384}{ALIUSC000000}{\ALIUSFontCormorantRegular}{13.920}{tool and potentially helpful for educational settings or training situations.}
  \ALIUSPlacedTextContent{90.951}{313.768}{418.257}{ALIUSC000000}{\ALIUSFontCormorantRegular}{13.920}{We have recently shown that people can even feel like they are strongly}
  \ALIUSPlacedTextContent{90.951}{333.208}{418.375}{ALIUSC000000}{\ALIUSFontCormorantRegular}{13.920}{embodying a grapefruit (Lesur, Aicher, et al., 2020). That's why I am}
  \ALIUSPlacedTextContent{90.951}{352.888}{418.341}{ALIUSC000000}{\ALIUSFontCormorantRegular}{13.920}{optimistic that healthy participants can basically embody anything, at least}
  \ALIUSPlacedTextContent{90.951}{372.328}{418.263}{ALIUSC000000}{\ALIUSFontCormorantRegular}{13.920}{as long as head-related visuo-motor coherency is given (Lesur et al., 2018b),}
  \ALIUSPlacedTextContent{90.951}{391.768}{418.330}{ALIUSC000000}{\ALIUSFontCormorantRegular}{13.920}{even if it might be an alief rather than an illusion. While I have never tried}
  \ALIUSPlacedTextContent{90.951}{411.208}{418.353}{ALIUSC000000}{\ALIUSFontCormorantRegular}{13.920}{to virtually embody several bodies and have a hard time imagining it,}
  \ALIUSPlacedTextContent{90.951}{430.888}{418.182}{ALIUSC000000}{\ALIUSFontCormorantRegular}{13.920}{previous literature suggests it is possible (e.g., Heydrich et al., 2013). I agree}
  \ALIUSPlacedTextContent{90.951}{450.328}{418.201}{ALIUSC000000}{\ALIUSFontCormorantRegular}{13.920}{with you that it is important to do more basic research before developing too}
  \ALIUSPlacedTextContent{90.951}{469.768}{418.107}{ALIUSC000000}{\ALIUSFontCormorantRegular}{13.920}{many therapeutic applications, especially when it comes to children, who are}
  \ALIUSPlacedTextContent{90.951}{489.208}{418.260}{ALIUSC000000}{\ALIUSFontCormorantRegular}{13.920}{presumably still developing their sense of a bodily self. I am surprised how}
  \ALIUSPlacedTextContent{90.951}{508.648}{418.281}{ALIUSC000000}{\ALIUSFontCormorantRegular}{13.920}{little is known about how embodied virtual reality exposure alters (bodily)}
  \ALIUSPlacedTextContent{90.951}{528.328}{418.214}{ALIUSC000000}{\ALIUSFontCormorantRegular}{13.920}{self-consciousness in children. It is one of my core objectives of my current}
  \ALIUSPlacedTextContent{90.951}{547.768}{418.270}{ALIUSC000000}{\ALIUSFontCormorantRegular}{13.920}{Swiss National Science Foundation-funded project to investigate these}
  \ALIUSPlacedTextContent{90.951}{567.208}{418.262}{ALIUSC000000}{\ALIUSFontCormorantRegular}{13.920}{mechanisms. It is important to carefully evaluate ethical considerations (e.g.,}
  \ALIUSPlacedTextContent{90.951}{586.648}{418.257}{ALIUSC000000}{\ALIUSFontCormorantRegular}{13.920}{Madary \& Metzinger, 2016), especially as potential therapeutic embodied}
  \ALIUSPlacedTextContent{90.951}{606.328}{418.255}{ALIUSC000000}{\ALIUSFontCormorantRegular}{13.920}{virtual reality, as well as embodied virtual reality for leisure (e.g., in games),}
  \ALIUSPlacedTextContent{90.951}{625.768}{418.229}{ALIUSC000000}{\ALIUSFontCormorantRegular}{13.920}{will use much longer exposure times than what we typically use in the}
  \ALIUSPlacedTextContent{90.951}{645.208}{66.200}{ALIUSC000000}{\ALIUSFontCormorantRegular}{13.920}{laboratory.}
  \ALIUSPlacedTextContent{90.951}{676.773}{418.395}{ALIUSC1F8135}{\ALIUSFontLatoLight}{12.357}{In the virtual reality literature, it is often claimed that good virtual reality}
  \ALIUSPlacedTextContent{90.951}{694.053}{418.386}{ALIUSC1F8135}{\ALIUSFontLatoLight}{12.357}{technology should induce a "sense of presence", defined as the sense of being}
  \ALIUSPlacedTextContent{90.951}{711.333}{414.986}{ALIUSC1F8135}{\ALIUSFontLatoLight}{12.357}{present within the virtual environment that one perceived through a head-}
  \ALIUSPlacedTextContent{90.951}{728.613}{400.654}{ALIUSC1F8135}{\ALIUSFontLatoLight}{12.357}{mounted display (e.g., Heeter, 1992; Held \& Durlach, 1992; Slater, 2009).}
  \ALIUSPlacedTextContent{90.951}{793.143}{122.397}{ALIUSC767171}{\ALIUSFontCormorantLight}{12.000}{ALIUS Bulletin n°4 (2020)}
  \ALIUSPlacedTextContent{387.435}{793.143}{120.728}{ALIUSC767171}{\ALIUSFontCormorantLight}{12.000}{aliusresearch.org/bulletin}
\end{tikzpicture}
\null
\newpage

% Page 17
\thispagestyle{empty}
\noindent\begin{tikzpicture}[remember picture,overlay,x=1bp,y=1bp,shift={(current page.north west)}]
  \ALIUSPlacedTextContent{90.951}{36.329}{258.472}{ALIUSC767171}{\ALIUSFontCormorantRegular}{13.920}{Lenggenhager – Bodily boundaries and beyond}
  \ALIUSPlacedTextContent{492.471}{36.329}{16.613}{ALIUSC000000}{\ALIUSFontCormorantRegular}{13.920}{98}
  \ALIUSPlacedTextContent{90.951}{73.173}{418.390}{ALIUSC1F8135}{\ALIUSFontLatoLight}{12.357}{There are a number of technical specifications of virtual reality systems that}
  \ALIUSPlacedTextContent{90.951}{90.453}{418.304}{ALIUSC1F8135}{\ALIUSFontLatoLight}{12.357}{appear to be strong mediators, if not requirements, for the induction of a}
  \ALIUSPlacedTextContent{90.951}{107.733}{418.429}{ALIUSC1F8135}{\ALIUSFontLatoLight}{12.357}{sense of presence. These include, among others: stereoscopy (providing}
  \ALIUSPlacedTextContent{90.951}{125.013}{418.426}{ALIUSC1F8135}{\ALIUSFontLatoLight}{12.357}{binocular depth cues), head tracking (providing action-contingent visual}
  \ALIUSPlacedTextContent{90.951}{142.293}{418.382}{ALIUSC1F8135}{\ALIUSFontLatoLight}{12.357}{feedback), hand/body tracking (giving the user a virtual body), high screen}
  \ALIUSPlacedTextContent{90.951}{159.333}{418.422}{ALIUSC1F8135}{\ALIUSFontLatoLight}{12.357}{refresh rate (enabling smooth visual motion), high screen resolution (providing}
  \ALIUSPlacedTextContent{90.951}{176.613}{418.484}{ALIUSC1F8135}{\ALIUSFontLatoLight}{12.357}{clear visual input), and wide field of view (providing peripheral visual input). A}
  \ALIUSPlacedTextContent{90.951}{193.893}{418.339}{ALIUSC1F8135}{\ALIUSFontLatoLight}{12.357}{meta-analysis of the relationship between the sense of presence (as measured}
  \ALIUSPlacedTextContent{90.951}{211.173}{418.395}{ALIUSC1F8135}{\ALIUSFontLatoLight}{12.357}{by various questionnaires) and different technical specifications virtual reality}
  \ALIUSPlacedTextContent{90.951}{228.453}{418.345}{ALIUSC1F8135}{\ALIUSFontLatoLight}{12.357}{systems in 83 studies found that head tracking was one of the most important}
  \ALIUSPlacedTextContent{90.951}{245.733}{384.506}{ALIUSC1F8135}{\ALIUSFontLatoLight}{12.357}{features to induce a sense of presence (Cummings \& Bailenson, 2016).}
  \ALIUSPlacedTextContent{90.951}{275.013}{418.399}{ALIUSC1F8135}{\ALIUSFontLatoLight}{12.357}{When these technical specifications are not met, the experience of virtual}
  \ALIUSPlacedTextContent{90.951}{292.293}{418.275}{ALIUSC1F8135}{\ALIUSFontLatoLight}{12.357}{reality users can become quite uncomfortable. In particular, the lack of head}
  \ALIUSPlacedTextContent{90.951}{309.333}{418.392}{ALIUSC1F8135}{\ALIUSFontLatoLight}{12.357}{tracking (or head tracking with high latency) causes a mismatch between visual}
  \ALIUSPlacedTextContent{90.951}{326.613}{418.384}{ALIUSC1F8135}{\ALIUSFontLatoLight}{12.357}{input and proprioceptive/vestibular cues about self-motion, since the}
  \ALIUSPlacedTextContent{90.951}{343.893}{418.467}{ALIUSC1F8135}{\ALIUSFontLatoLight}{12.357}{viewpoint rendered by the computer is not sensitive to the user's head}
  \ALIUSPlacedTextContent{90.951}{361.173}{418.372}{ALIUSC1F8135}{\ALIUSFontLatoLight}{12.357}{movements. Such mismatch is often associated with an unpleasant}
  \ALIUSPlacedTextContent{90.951}{378.453}{418.348}{ALIUSC1F8135}{\ALIUSFontLatoLight}{12.357}{combination of symptoms known as 'cybersickness' – a special case of motion}
  \ALIUSPlacedTextContent{90.951}{395.733}{418.337}{ALIUSC1F8135}{\ALIUSFontLatoLight}{12.357}{sickness (Gallagher \& Ferrè, 2018). There is some evidence that cybersickness}
  \ALIUSPlacedTextContent{90.951}{413.013}{418.458}{ALIUSC1F8135}{\ALIUSFontLatoLight}{12.357}{is negatively correlated with scores on various presence questionnaires}
  \ALIUSPlacedTextContent{90.951}{430.293}{418.338}{ALIUSC1F8135}{\ALIUSFontLatoLight}{12.357}{(Weech et al., 2019). This raises the following question: is the sense of}
  \ALIUSPlacedTextContent{90.951}{447.333}{418.384}{ALIUSC1F8135}{\ALIUSFontLatoLight}{12.357}{presence in virtual reality really a positive experience that one has – in addition}
  \ALIUSPlacedTextContent{90.951}{464.613}{418.385}{ALIUSC1F8135}{\ALIUSFontLatoLight}{12.357}{to whatever else one might experience while using virtual reality –, or is it}
  \ALIUSPlacedTextContent{90.951}{481.893}{418.315}{ALIUSC1F8135}{\ALIUSFontLatoLight}{12.357}{simply the absence of the discomfort and abnormal sensations, such as}
  \ALIUSPlacedTextContent{90.951}{499.173}{418.414}{ALIUSC1F8135}{\ALIUSFontLatoLight}{12.357}{cybersickness, associated with less sophisticated virtual reality systems? Do}
  \ALIUSPlacedTextContent{90.951}{516.453}{418.520}{ALIUSC1F8135}{\ALIUSFontLatoLight}{12.357}{you think one of these two hypotheses is more plausible than the other, on}
  \ALIUSPlacedTextContent{90.951}{533.733}{418.401}{ALIUSC1F8135}{\ALIUSFontLatoLight}{12.357}{the basis of evidence provided by your own research or virtual reality research}
  \ALIUSPlacedTextContent{90.951}{551.013}{48.706}{ALIUSC1F8135}{\ALIUSFontLatoLight}{12.357}{at large?}
  \ALIUSPlacedTextContent{90.951}{580.169}{418.279}{ALIUSC000000}{\ALIUSFontCormorantRegular}{13.920}{I think it depends on the content of the virtual reality. If you imagine a virtual}
  \ALIUSPlacedTextContent{90.951}{599.609}{418.266}{ALIUSC000000}{\ALIUSFontCormorantRegular}{13.920}{reality that exactly simulates reality, I would clearly expect that a strong sense}
  \ALIUSPlacedTextContent{90.951}{619.048}{418.345}{ALIUSC000000}{\ALIUSFontCormorantRegular}{13.920}{of presence is just the absence of any discomfort and not per se a pleasant}
  \ALIUSPlacedTextContent{90.951}{638.489}{418.219}{ALIUSC000000}{\ALIUSFontCormorantRegular}{13.920}{sensation. However, if you take virtual reality as a tool to let participants feel}
  \ALIUSPlacedTextContent{90.951}{657.928}{418.288}{ALIUSC000000}{\ALIUSFontCormorantRegular}{13.920}{the sense of presence in a pleasant and surprising world that they otherwise}
  \ALIUSPlacedTextContent{90.951}{677.609}{418.349}{ALIUSC000000}{\ALIUSFontCormorantRegular}{13.920}{could not be present, apart from maybe in dreams or during a psychedelic}
  \ALIUSPlacedTextContent{90.951}{697.048}{418.275}{ALIUSC000000}{\ALIUSFontCormorantRegular}{13.920}{experience or similar, the mere sense of presence might be associated with a}
  \ALIUSPlacedTextContent{90.951}{716.489}{418.175}{ALIUSC000000}{\ALIUSFontCormorantRegular}{13.920}{positive experience. But of course, depending on the content of the}
  \ALIUSPlacedTextContent{90.951}{735.929}{324.578}{ALIUSC000000}{\ALIUSFontCormorantRegular}{13.920}{environment this could also rapidly turn into a nightmare.}
  \ALIUSPlacedTextContent{90.951}{793.143}{122.397}{ALIUSC767171}{\ALIUSFontCormorantLight}{12.000}{ALIUS Bulletin n°4 (2020)}
  \ALIUSPlacedTextContent{387.435}{793.143}{120.728}{ALIUSC767171}{\ALIUSFontCormorantLight}{12.000}{aliusresearch.org/bulletin}
\end{tikzpicture}
\null
\newpage

% Page 18
\thispagestyle{empty}
\noindent\begin{tikzpicture}[remember picture,overlay,x=1bp,y=1bp,shift={(current page.north west)}]
  \ALIUSPlacedTextContent{90.951}{36.329}{258.472}{ALIUSC767171}{\ALIUSFontCormorantRegular}{13.920}{Lenggenhager – Bodily boundaries and beyond}
  \ALIUSPlacedTextContent{492.951}{36.329}{16.277}{ALIUSC000000}{\ALIUSFontCormorantRegular}{13.920}{99}
  \ALIUSPlacedTextContent{90.951}{73.173}{418.322}{ALIUSC1F8135}{\ALIUSFontLatoLight}{12.357}{If ethics, technology, and finances were no concern, what would be your}
  \ALIUSPlacedTextContent{90.951}{90.453}{125.273}{ALIUSC1F8135}{\ALIUSFontLatoLight}{12.357}{dream research study?}
  \ALIUSPlacedTextContent{90.951}{119.608}{418.199}{ALIUSC000000}{\ALIUSFontCormorantRegular}{13.920}{Hmm, as you know I am a great fan of the virtual body swapping, as for}
  \ALIUSPlacedTextContent{90.951}{139.049}{418.315}{ALIUSC000000}{\ALIUSFontCormorantRegular}{13.920}{example in the Machine To Be Another (Oliveira et al., 2016) and I love trying}
  \ALIUSPlacedTextContent{90.951}{158.489}{418.303}{ALIUSC000000}{\ALIUSFontCormorantRegular}{13.920}{out all experiments as a participant before I start an experiment. My dream}
  \ALIUSPlacedTextContent{90.951}{177.928}{418.260}{ALIUSC000000}{\ALIUSFontCormorantRegular}{13.920}{study would be to swap bodies in reality rather than in virtual reality. Maybe}
  \ALIUSPlacedTextContent{90.951}{197.609}{418.176}{ALIUSC000000}{\ALIUSFontCormorantRegular}{13.920}{related to what was mentioned before, to create an experiment in which you}
  \ALIUSPlacedTextContent{90.951}{217.049}{418.207}{ALIUSC000000}{\ALIUSFontCormorantRegular}{13.920}{could progressively swap out body parts and organs of the other person and}
  \ALIUSPlacedTextContent{90.951}{236.488}{418.271}{ALIUSC000000}{\ALIUSFontCormorantRegular}{13.920}{test how it changes affective and cognitive processes, especially the sense of a}
  \ALIUSPlacedTextContent{90.951}{255.929}{418.272}{ALIUSC000000}{\ALIUSFontCormorantRegular}{13.920}{bodily self and self-consciousness. If this answer is a bit too science-fictional,}
  \ALIUSPlacedTextContent{90.951}{275.608}{418.310}{ALIUSC000000}{\ALIUSFontCormorantRegular}{13.920}{I would dream to do a study in humans as I suggested in animals above,}
  \ALIUSPlacedTextContent{90.951}{295.048}{418.251}{ALIUSC000000}{\ALIUSFontCormorantRegular}{13.920}{letting them grow up with permanently altered multisensory perspectives or}
  \ALIUSPlacedTextContent{90.951}{314.488}{418.276}{ALIUSC000000}{\ALIUSFontCormorantRegular}{13.920}{multisensory contingencies. If that sounds too scary, I would at least love to}
  \ALIUSPlacedTextContent{90.951}{333.928}{418.246}{ALIUSC000000}{\ALIUSFontCormorantRegular}{13.920}{do more long-term exposure experiments. But even those need a lot of money}
  \ALIUSPlacedTextContent{90.951}{353.608}{360.095}{ALIUSC000000}{\ALIUSFontCormorantRegular}{13.920}{and facilities. And there might of course still be ethical concerns.}
  \ALIUSPlacedTextContent{90.951}{385.173}{418.356}{ALIUSC1F8135}{\ALIUSFontLatoLight}{12.357}{With the rapid advancement of technology, what types of questions do you}
  \ALIUSPlacedTextContent{90.951}{402.453}{418.478}{ALIUSC1F8135}{\ALIUSFontLatoLight}{12.357}{anticipate the field of bodily self-consciousness research will be attempting to}
  \ALIUSPlacedTextContent{90.951}{419.733}{110.549}{ALIUSC1F8135}{\ALIUSFontLatoLight}{12.357}{answer in 20 years?}
  \ALIUSPlacedTextContent{90.951}{448.648}{418.154}{ALIUSC000000}{\ALIUSFontCormorantRegular}{13.920}{The question on how digitalization and digital interactions change bodily}
  \ALIUSPlacedTextContent{90.951}{468.328}{418.270}{ALIUSC000000}{\ALIUSFontCormorantRegular}{13.920}{self-consciousness will be increasingly relevant. The fact that we interact}
  \ALIUSPlacedTextContent{90.951}{487.768}{418.163}{ALIUSC000000}{\ALIUSFontCormorantRegular}{13.920}{more and more in quasi disembodied states through various digital media}
  \ALIUSPlacedTextContent{90.951}{507.208}{418.065}{ALIUSC000000}{\ALIUSFontCormorantRegular}{13.920}{might shape the way we perceive ourselves and others. The fact that we might}
  \ALIUSPlacedTextContent{90.951}{526.648}{418.260}{ALIUSC000000}{\ALIUSFontCormorantRegular}{13.920}{increasingly have the possibility to digitally change and represent certain}
  \ALIUSPlacedTextContent{90.951}{546.328}{418.283}{ALIUSC000000}{\ALIUSFontCormorantRegular}{13.920}{aspects of our body and self in much more flexible ways (e.g., using facial}
  \ALIUSPlacedTextContent{90.951}{565.768}{414.974}{ALIUSC000000}{\ALIUSFontCormorantRegular}{13.920}{filters, self-chosen avatars or holograms) might change our notion of self-}
  \ALIUSPlacedTextContent{90.951}{585.208}{259.731}{ALIUSC000000}{\ALIUSFontCormorantRegular}{13.920}{consciousness and reveal new ethical questions.}
  \ALIUSPlacedTextContent{114.271}{631.959}{14.160}{ALIUSC595959}{\ALIUSFontCormorantLight}{48.000}{?}
  \ALIUSPlacedTextContent{145.951}{640.987}{314.234}{ALIUSC000000}{\ALIUSFontLatoLight}{15.120}{The fact that we interact more and more in quasi-}
  \ALIUSPlacedTextContent{145.951}{658.987}{318.061}{ALIUSC000000}{\ALIUSFontLatoLight}{15.120}{disembodied states through various digital media}
  \ALIUSPlacedTextContent{145.951}{676.987}{317.996}{ALIUSC000000}{\ALIUSFontLatoLight}{15.120}{might shape the way we perceive ourselves and}
  \ALIUSPlacedTextContent{474.271}{688.359}{14.160}{ALIUSC595959}{\ALIUSFontCormorantLight}{48.000}{?}
  \ALIUSPlacedTextContent{145.951}{694.987}{44.825}{ALIUSC000000}{\ALIUSFontLatoLight}{15.120}{others.}
  \ALIUSPlacedTextContent{90.951}{793.143}{122.397}{ALIUSC767171}{\ALIUSFontCormorantLight}{12.000}{ALIUS Bulletin n°4 (2020)}
  \ALIUSPlacedTextContent{387.435}{793.143}{120.728}{ALIUSC767171}{\ALIUSFontCormorantLight}{12.000}{aliusresearch.org/bulletin}
\end{tikzpicture}
\null
\newpage

% Page 19
\thispagestyle{empty}
\noindent\begin{tikzpicture}[remember picture,overlay,x=1bp,y=1bp,shift={(current page.north west)}]
  \fill[ALIUSC0000FF] (126.951bp,-164.391bp) rectangle ++(213.600bp,-0.720bp);
  \fill[ALIUSC0000FF] (126.951bp,-233.511bp) rectangle ++(190.800bp,-0.720bp);
  \fill[ALIUSC0000FF] (126.951bp,-302.391bp) rectangle ++(196.080bp,-0.720bp);
  \fill[ALIUSC0000FF] (126.951bp,-387.111bp) rectangle ++(206.640bp,-0.720bp);
  \fill[ALIUSC0000FF] (214.551bp,-440.391bp) rectangle ++(192.240bp,-0.720bp);
  \fill[ALIUSC0000FF] (126.951bp,-493.671bp) rectangle ++(156.960bp,-0.720bp);
  \fill[ALIUSC0000FF] (126.951bp,-548.391bp) rectangle ++(187.440bp,-0.720bp);
  \fill[ALIUSC0000FF] (126.951bp,-601.431bp) rectangle ++(199.440bp,-0.720bp);
  \fill[ALIUSC0000FF] (126.951bp,-686.151bp) rectangle ++(248.400bp,-0.720bp);
  \fill[ALIUSC0000FF] (126.951bp,-739.431bp) rectangle ++(156.240bp,-0.720bp);
  \ALIUSPlacedTextContent{90.951}{36.329}{258.472}{ALIUSC767171}{\ALIUSFontCormorantRegular}{13.920}{Lenggenhager – Bodily boundaries and beyond}
  \ALIUSPlacedTextContent{487.911}{36.329}{21.261}{ALIUSC000000}{\ALIUSFontCormorantRegular}{13.920}{100}
  \ALIUSPlacedTextContent{264.617}{73.132}{71.454}{ALIUSC000000}{\ALIUSFontLatoLight}{13.920}{References}
  \ALIUSPlacedTextContent{90.951}{120.016}{418.086}{ALIUSC000000}{\ALIUSFontCormorantRegular}{12.960}{Alsmith, A., \& Longo, M. R. (2014). Where exactly am I? Self-location judgements}
  \ALIUSPlacedTextContent{126.951}{135.616}{178.827}{ALIUSC000000}{\ALIUSFontCormorantRegular}{12.960}{distribute between head and torso.}
  \ALIUSPlacedTextContent{305.800}{135.616}{148.885}{ALIUSC000000}{\ALIUSFontCormorantItalic}{12.960}{Consciousness and Cognition, 24}
  \ALIUSPlacedTextContent{454.712}{135.616}{49.267}{ALIUSC000000}{\ALIUSFontCormorantRegular}{12.960}{, 70– 74.}
  \ALIUSPlacedTextContent{126.951}{151.456}{213.552}{ALIUSC0000FF}{\ALIUSFontCormorantRegular}{12.960}{https://doi.org/10.1016/j.concog.2013.12.005}
  \ALIUSPlacedTextContent{90.951}{173.296}{417.940}{ALIUSC000000}{\ALIUSFontCormorantRegular}{12.960}{Armel, K. C., \& Ramachandran, V. S. (2003). Projecting sensations to external}
  \ALIUSPlacedTextContent{126.951}{188.896}{262.511}{ALIUSC000000}{\ALIUSFontCormorantRegular}{12.960}{objects: Evidence from skin conductance response.}
  \ALIUSPlacedTextContent{390.398}{188.896}{118.597}{ALIUSC000000}{\ALIUSFontCormorantItalic}{12.960}{Proceedings of the Royal}
  \ALIUSPlacedTextContent{126.951}{204.736}{203.920}{ALIUSC000000}{\ALIUSFontCormorantItalic}{12.960}{Society of London B: Biological Sciences, 270}
  \ALIUSPlacedTextContent{330.890}{204.736}{93.354}{ALIUSC000000}{\ALIUSFontCormorantRegular}{12.960}{(1523), 1499–1506.}
  \ALIUSPlacedTextContent{126.951}{220.576}{190.661}{ALIUSC0000FF}{\ALIUSFontCormorantRegular}{12.960}{https://doi.org/10.1098/rspb.2003.2364}
  \ALIUSPlacedTextContent{90.951}{242.176}{418.227}{ALIUSC000000}{\ALIUSFontCormorantRegular}{12.960}{Barnsley, N., McAuley, J. H., Mohan, R., Dey, A., Thomas, P., \& Moseley, G. L.}
  \ALIUSPlacedTextContent{126.951}{258.016}{382.003}{ALIUSC000000}{\ALIUSFontCormorantRegular}{12.960}{(2011). The rubber hand illusion increases histamine reactivity in the real}
  \ALIUSPlacedTextContent{126.951}{273.616}{25.689}{ALIUSC000000}{\ALIUSFontCormorantRegular}{12.960}{arm.}
  \ALIUSPlacedTextContent{152.651}{273.616}{88.429}{ALIUSC000000}{\ALIUSFontCormorantItalic}{12.960}{Current Biology, 21}
  \ALIUSPlacedTextContent{241.088}{273.616}{92.249}{ALIUSC000000}{\ALIUSFontCormorantRegular}{12.960}{(23), R945–R946.}
  \ALIUSPlacedTextContent{126.951}{289.456}{196.033}{ALIUSC0000FF}{\ALIUSFontCormorantRegular}{12.960}{https://doi.org/10.1016/j.cub.2011.10.039}
  \ALIUSPlacedTextContent{90.951}{311.056}{418.028}{ALIUSC000000}{\ALIUSFontCormorantRegular}{12.960}{Belser, A. B., Agin-Liebes, G., Swift, T. C., Terrana, S., Devenot, N., Friedman, H.}
  \ALIUSPlacedTextContent{126.951}{326.896}{378.987}{ALIUSC000000}{\ALIUSFontCormorantRegular}{12.960}{L., Guss, J., Bossis, A., \& Ross, S. (2017). Patient Experiences of Psilocybin-}
  \ALIUSPlacedTextContent{126.951}{342.736}{381.991}{ALIUSC000000}{\ALIUSFontCormorantRegular}{12.960}{Assisted Psychotherapy: An Interpretative Phenomenological Analysis.}
  \ALIUSPlacedTextContent{126.951}{358.336}{172.038}{ALIUSC000000}{\ALIUSFontCormorantItalic}{12.960}{Journal of Humanistic Psychology, 57}
  \ALIUSPlacedTextContent{299.066}{358.336}{69.162}{ALIUSC000000}{\ALIUSFontCormorantRegular}{12.960}{(4), 354–388.}
  \ALIUSPlacedTextContent{126.951}{374.176}{206.563}{ALIUSC0000FF}{\ALIUSFontCormorantRegular}{12.960}{https://doi.org/10.1177/0022167817706884}
  \ALIUSPlacedTextContent{90.951}{395.776}{418.059}{ALIUSC000000}{\ALIUSFontCormorantRegular}{12.960}{Bergouignan, L., Nyberg, L., \& Ehrsson, H. H. (2014). Out-of-body–induced}
  \ALIUSPlacedTextContent{126.951}{411.616}{118.493}{ALIUSC000000}{\ALIUSFontCormorantRegular}{12.960}{hippocampal amnesia.}
  \ALIUSPlacedTextContent{248.822}{411.616}{257.211}{ALIUSC000000}{\ALIUSFontCormorantItalic}{12.960}{Proceedings of the National Academy of Sciences, 1}
  \ALIUSPlacedTextContent{126.951}{427.456}{7.554}{ALIUSC000000}{\ALIUSFontCormorantItalic}{12.960}{11}
  \ALIUSPlacedTextContent{134.516}{427.456}{80.055}{ALIUSC000000}{\ALIUSFontCormorantRegular}{12.960}{(12), 4421–4426.}
  \ALIUSPlacedTextContent{214.581}{427.456}{192.093}{ALIUSC0000FF}{\ALIUSFontCormorantRegular}{12.960}{https://doi.org/10.1073/pnas.1318801111}
  \ALIUSPlacedTextContent{90.951}{449.056}{418.023}{ALIUSC000000}{\ALIUSFontCormorantRegular}{12.960}{Blanke, O. (2012). Multisensory brain mechanisms of bodily self-consciousness.}
  \ALIUSPlacedTextContent{126.951}{464.896}{234.402}{ALIUSC000000}{\ALIUSFontCormorantRegular}{12.960}{Nature Reviews Neuroscience, 13(8), 556–571.}
  \ALIUSPlacedTextContent{126.951}{480.736}{157.010}{ALIUSC0000FF}{\ALIUSFontCormorantRegular}{12.960}{https://doi.org/10.1038/nrn3292}
  \ALIUSPlacedTextContent{383.859}{502.239}{4.303}{ALIUSC000000}{\ALIUSFontCambria}{12.960}{-}
  \ALIUSPlacedTextContent{398.367}{502.239}{4.303}{ALIUSC000000}{\ALIUSFontCambria}{12.960}{-}
  \ALIUSPlacedTextContent{90.951}{502.576}{292.892}{ALIUSC000000}{\ALIUSFontCormorantRegular}{12.960}{Blanke, O., Landis, T., Spinelli, L., \& Seeck, M. (2004). Out}
  \ALIUSPlacedTextContent{388.176}{502.576}{10.180}{ALIUSC000000}{\ALIUSFontCormorantRegular}{12.960}{of}
  \ALIUSPlacedTextContent{402.684}{502.576}{106.316}{ALIUSC000000}{\ALIUSFontCormorantRegular}{12.960}{body experience and}
  \ALIUSPlacedTextContent{126.951}{519.376}{168.595}{ALIUSC000000}{\ALIUSFontCormorantRegular}{12.960}{autoscopy of neurological origin.}
  \ALIUSPlacedTextContent{295.621}{519.376}{45.547}{ALIUSC000000}{\ALIUSFontCormorantItalic}{12.960}{Brain, 127}
  \ALIUSPlacedTextContent{341.198}{519.376}{35.307}{ALIUSC000000}{\ALIUSFontCormorantRegular}{12.960}{(2), 243}
  \ALIUSPlacedTextContent{383.031}{519.376}{25.535}{ALIUSC000000}{\ALIUSFontCormorantRegular}{12.960}{258.}
  \ALIUSPlacedTextContent{376.531}{520.218}{6.480}{ALIUSC000000}{\ALIUSFontCormorantRegular}{12.960}{–}
  \ALIUSPlacedTextContent{126.951}{535.456}{187.395}{ALIUSC0000FF}{\ALIUSFontCormorantRegular}{12.960}{https://doi.org/10.1093/brain/awh040}
  \ALIUSPlacedTextContent{90.951}{557.056}{418.035}{ALIUSC000000}{\ALIUSFontCormorantRegular}{12.960}{Blanke, O., \& Metzinger, T. (2009). Full-body illusions and minimal phenomenal}
  \ALIUSPlacedTextContent{126.951}{572.656}{48.237}{ALIUSC000000}{\ALIUSFontCormorantRegular}{12.960}{selfhood.}
  \ALIUSPlacedTextContent{175.166}{572.656}{145.498}{ALIUSC000000}{\ALIUSFontCormorantItalic}{12.960}{Trends in Cognitive Sciences, 13}
  \ALIUSPlacedTextContent{320.711}{572.656}{48.558}{ALIUSC000000}{\ALIUSFontCormorantRegular}{12.960}{(1), 7–13.}
  \ALIUSPlacedTextContent{126.951}{588.496}{199.409}{ALIUSC0000FF}{\ALIUSFontCormorantRegular}{12.960}{https://doi.org/10.1016/j.tics.2008.10.003}
  \ALIUSPlacedTextContent{90.951}{610.336}{418.080}{ALIUSC000000}{\ALIUSFontCormorantRegular}{12.960}{Blanke, O., Mohr, C., Michel, C. M., Pascual-Leone, A., Brugger, P., Seeck, M.,}
  \ALIUSPlacedTextContent{126.951}{625.935}{382.073}{ALIUSC000000}{\ALIUSFontCormorantRegular}{12.960}{Landis, T., \& Thut, G. (2005). Linking Out-of-Body Experience and Self}
  \ALIUSPlacedTextContent{126.951}{641.776}{382.071}{ALIUSC000000}{\ALIUSFontCormorantRegular}{12.960}{Processing to Mental Own-Body Imagery at the Temporoparietal Junction.}
  \ALIUSPlacedTextContent{126.951}{657.376}{123.812}{ALIUSC000000}{\ALIUSFontCormorantItalic}{12.960}{Journal of Neuroscience, 25}
  \ALIUSPlacedTextContent{250.825}{657.376}{67.394}{ALIUSC000000}{\ALIUSFontCormorantRegular}{12.960}{(3), 550–557.}
  \ALIUSPlacedTextContent{126.951}{673.216}{248.478}{ALIUSC0000FF}{\ALIUSFontCormorantRegular}{12.960}{https://doi.org/10.1523/JNEUROSCI.2612-04.2005}
  \ALIUSPlacedTextContent{90.951}{695.056}{414.987}{ALIUSC000000}{\ALIUSFontCormorantRegular}{12.960}{Blanke, O., Ortigue, S., Landis, T., \& Seeck, M. (2002). Stimulating illusory own-}
  \ALIUSPlacedTextContent{126.951}{710.656}{93.055}{ALIUSC000000}{\ALIUSFontCormorantRegular}{12.960}{body perceptions.}
  \ALIUSPlacedTextContent{220.015}{710.656}{53.538}{ALIUSC000000}{\ALIUSFontCormorantItalic}{12.960}{Nature, 419}
  \ALIUSPlacedTextContent{273.574}{710.656}{87.765}{ALIUSC000000}{\ALIUSFontCormorantRegular}{12.960}{(6904), 269–270.}
  \ALIUSPlacedTextContent{126.951}{726.496}{156.259}{ALIUSC0000FF}{\ALIUSFontCormorantRegular}{12.960}{https://doi.org/10.1038/419269a}
  \ALIUSPlacedTextContent{90.951}{793.143}{122.397}{ALIUSC767171}{\ALIUSFontCormorantLight}{12.000}{ALIUS Bulletin n°4 (2020)}
  \ALIUSPlacedTextContent{387.435}{793.143}{120.728}{ALIUSC767171}{\ALIUSFontCormorantLight}{12.000}{aliusresearch.org/bulletin}
\end{tikzpicture}
\null
\newpage

% Page 20
\thispagestyle{empty}
\noindent\begin{tikzpicture}[remember picture,overlay,x=1bp,y=1bp,shift={(current page.north west)}]
  \fill[ALIUSC0000FF] (224.871bp,-101.751bp) rectangle ++(145.680bp,-0.720bp);
  \fill[ALIUSC0000FF] (149.991bp,-154.791bp) rectangle ++(232.080bp,-0.720bp);
  \fill[ALIUSC0000FF] (194.631bp,-208.071bp) rectangle ++(199.920bp,-0.720bp);
  \fill[ALIUSC0000FF] (257.511bp,-261.351bp) rectangle ++(224.160bp,-0.720bp);
  \fill[ALIUSC0000FF] (183.351bp,-336.231bp) rectangle ++(213.600bp,-0.720bp);
  \fill[ALIUSC0000FF] (126.951bp,-389.511bp) rectangle ++(234.000bp,-0.720bp);
  \fill[ALIUSC0000FF] (126.951bp,-442.791bp) rectangle ++(206.640bp,-0.720bp);
  \fill[ALIUSC0000FF] (126.951bp,-480.231bp) rectangle ++(373.200bp,-0.720bp);
  \fill[ALIUSC0000FF] (126.951bp,-496.071bp) rectangle ++(88.320bp,-0.720bp);
  \fill[ALIUSC0000FF] (126.951bp,-533.511bp) rectangle ++(206.640bp,-0.720bp);
  \fill[ALIUSC0000FF] (126.951bp,-602.391bp) rectangle ++(239.520bp,-0.720bp);
  \fill[ALIUSC0000FF] (126.951bp,-655.671bp) rectangle ++(194.400bp,-0.720bp);
  \fill[ALIUSC0000FF] (309.111bp,-693.111bp) rectangle ++(193.440bp,-0.720bp);
  \fill[ALIUSC0000FF] (126.951bp,-762.231bp) rectangle ++(200.640bp,-0.720bp);
  \ALIUSPlacedTextContent{90.951}{36.329}{258.472}{ALIUSC767171}{\ALIUSFontCormorantRegular}{13.920}{Lenggenhager – Bodily boundaries and beyond}
  \ALIUSPlacedTextContent{490.071}{36.329}{19.231}{ALIUSC000000}{\ALIUSFontCormorantRegular}{13.920}{101}
  \ALIUSPlacedTextContent{90.951}{72.976}{378.659}{ALIUSC000000}{\ALIUSFontCormorantRegular}{12.960}{Botvinick, M., \& Cohen, J. (1998). Rubber hands 'feel' touch that eyes see.}
  \ALIUSPlacedTextContent{470.175}{72.976}{38.813}{ALIUSC000000}{\ALIUSFontCormorantItalic}{12.960}{Nature,}
  \ALIUSPlacedTextContent{126.951}{88.816}{13.673}{ALIUSC000000}{\ALIUSFontCormorantItalic}{12.960}{391}
  \ALIUSPlacedTextContent{140.639}{88.816}{84.302}{ALIUSC000000}{\ALIUSFontCormorantRegular}{12.960}{(6669), 756–756.}
  \ALIUSPlacedTextContent{224.956}{88.816}{145.450}{ALIUSC0000FF}{\ALIUSFontCormorantRegular}{12.960}{https://doi.org/10.1038/35784}
  \ALIUSPlacedTextContent{90.951}{110.416}{418.052}{ALIUSC000000}{\ALIUSFontCormorantRegular}{12.960}{Brugger, P., Blanke, O., Regard, M., Bradford, D. T., \& Landis, T. (2006). Polyopic}
  \ALIUSPlacedTextContent{126.951}{126.016}{288.207}{ALIUSC000000}{\ALIUSFontCormorantRegular}{12.960}{Heautoscopy: Case Report and Review of the Literature.}
  \ALIUSPlacedTextContent{414.972}{126.016}{46.970}{ALIUSC000000}{\ALIUSFontCormorantItalic}{12.960}{Cortex, 42}
  \ALIUSPlacedTextContent{461.982}{126.016}{43.948}{ALIUSC000000}{\ALIUSFontCormorantRegular}{12.960}{(5), 666–}
  \ALIUSPlacedTextContent{126.951}{141.856}{23.077}{ALIUSC000000}{\ALIUSFontCormorantRegular}{12.960}{674.}
  \ALIUSPlacedTextContent{150.038}{141.856}{231.961}{ALIUSC0000FF}{\ALIUSFontCormorantRegular}{12.960}{https://doi.org/10.1016/S0010-9452(08)70403-9}
  \ALIUSPlacedTextContent{90.951}{163.696}{418.050}{ALIUSC000000}{\ALIUSFontCormorantRegular}{12.960}{Brugger, P., Lenggenhager, B., \& Giummarra, M. J. (2013). Xenomelia: A Social}
  \ALIUSPlacedTextContent{126.951}{179.296}{315.792}{ALIUSC000000}{\ALIUSFontCormorantRegular}{12.960}{Neuroscience View of Altered Bodily Self-Consciousness.}
  \ALIUSPlacedTextContent{447.334}{179.296}{61.625}{ALIUSC000000}{\ALIUSFontCormorantItalic}{12.960}{Frontiers in}
  \ALIUSPlacedTextContent{126.951}{195.136}{62.052}{ALIUSC000000}{\ALIUSFontCormorantItalic}{12.960}{Psychology, 4}
  \ALIUSPlacedTextContent{189.010}{195.136}{5.594}{ALIUSC000000}{\ALIUSFontCormorantRegular}{12.960}{.}
  \ALIUSPlacedTextContent{194.613}{195.136}{199.973}{ALIUSC0000FF}{\ALIUSFontCormorantRegular}{12.960}{https://doi.org/10.3389/fpsyg.2013.00204}
  \ALIUSPlacedTextContent{90.951}{216.736}{414.987}{ALIUSC000000}{\ALIUSFontCormorantRegular}{12.960}{Cummings, J. J., \& Bailenson, J. N. (2016). How Immersive Is Enough? A Meta-}
  \ALIUSPlacedTextContent{126.951}{232.576}{348.842}{ALIUSC000000}{\ALIUSFontCormorantRegular}{12.960}{Analysis of the Effect of Immersive Technology on User Presence.}
  \ALIUSPlacedTextContent{477.351}{232.576}{31.635}{ALIUSC000000}{\ALIUSFontCormorantItalic}{12.960}{Media}
  \ALIUSPlacedTextContent{126.951}{248.416}{65.791}{ALIUSC000000}{\ALIUSFontCormorantItalic}{12.960}{Psychology, 19}
  \ALIUSPlacedTextContent{192.755}{248.416}{64.796}{ALIUSC000000}{\ALIUSFontCormorantRegular}{12.960}{(2), 272–309.}
  \ALIUSPlacedTextContent{257.571}{248.416}{224.007}{ALIUSC0000FF}{\ALIUSFontCormorantRegular}{12.960}{https://doi.org/10.1080/15213269.2015.1015740}
  \ALIUSPlacedTextContent{90.951}{270.256}{418.033}{ALIUSC000000}{\ALIUSFontCormorantRegular}{12.960}{Dieguez, S. (2018). The Illusion Illusion. Constructivist Foundations, 14(1), 108–110.}
  \ALIUSPlacedTextContent{90.951}{291.856}{417.945}{ALIUSC000000}{\ALIUSFontCormorantRegular}{12.960}{Fiorio, M., Modenese, M., \& Cesari, P. (2020). The rubber hand illusion in hypnosis}
  \ALIUSPlacedTextContent{126.951}{307.456}{293.362}{ALIUSC000000}{\ALIUSFontCormorantRegular}{12.960}{provides new insights into the sense of body ownership.}
  \ALIUSPlacedTextContent{421.440}{307.456}{87.578}{ALIUSC000000}{\ALIUSFontCormorantItalic}{12.960}{Scientific Reports,}
  \ALIUSPlacedTextContent{126.951}{323.296}{9.460}{ALIUSC000000}{\ALIUSFontCormorantItalic}{12.960}{10}
  \ALIUSPlacedTextContent{136.427}{323.296}{47.006}{ALIUSC000000}{\ALIUSFontCormorantRegular}{12.960}{(1), 5706.}
  \ALIUSPlacedTextContent{183.447}{323.296}{213.580}{ALIUSC0000FF}{\ALIUSFontCormorantRegular}{12.960}{https://doi.org/10.1038/s41598-020-62745-x}
  \ALIUSPlacedTextContent{90.951}{345.136}{418.025}{ALIUSC000000}{\ALIUSFontCormorantRegular}{12.960}{Fuchs, T., \& Schlimme, J. E. (2009). Embodiment and psychopathology: A}
  \ALIUSPlacedTextContent{126.951}{360.736}{158.519}{ALIUSC000000}{\ALIUSFontCormorantRegular}{12.960}{phenomenological perspective.}
  \ALIUSPlacedTextContent{285.455}{360.736}{157.397}{ALIUSC000000}{\ALIUSFontCormorantItalic}{12.960}{Current Opinion in Psychiatry, 22}
  \ALIUSPlacedTextContent{442.870}{360.736}{65.882}{ALIUSC000000}{\ALIUSFontCormorantRegular}{12.960}{(6), 570– 575.}
  \ALIUSPlacedTextContent{126.951}{376.576}{233.947}{ALIUSC0000FF}{\ALIUSFontCormorantRegular}{12.960}{https://doi.org/10.1097/YCO.0b013e3283318e5c}
  \ALIUSPlacedTextContent{90.951}{398.176}{417.997}{ALIUSC000000}{\ALIUSFontCormorantRegular}{12.960}{Gallagher, M., \& Ferrè, E. R. (2018). Cybersickness: A Multisensory Integration}
  \ALIUSPlacedTextContent{126.951}{414.016}{63.276}{ALIUSC000000}{\ALIUSFontCormorantRegular}{12.960}{Perspective.}
  \ALIUSPlacedTextContent{190.220}{414.016}{117.765}{ALIUSC000000}{\ALIUSFontCormorantItalic}{12.960}{Multisensory Research, 31}
  \ALIUSPlacedTextContent{307.997}{414.016}{69.578}{ALIUSC000000}{\ALIUSFontCormorantRegular}{12.960}{(7), 645–674.}
  \ALIUSPlacedTextContent{126.951}{429.856}{206.673}{ALIUSC0000FF}{\ALIUSFontCormorantRegular}{12.960}{https://doi.org/10.1163/22134808-20181293}
  \ALIUSPlacedTextContent{90.951}{451.456}{293.132}{ALIUSC000000}{\ALIUSFontCormorantRegular}{12.960}{Gallagher, S. (2005). How the Body Shapes the Mind.}
  \ALIUSPlacedTextContent{387.007}{451.456}{116.355}{ALIUSC000000}{\ALIUSFontCormorantItalic}{12.960}{Oxford University Press}
  \ALIUSPlacedTextContent{503.390}{451.456}{5.594}{ALIUSC000000}{\ALIUSFontCormorantRegular}{12.960}{.}
  \ALIUSPlacedTextContent{126.951}{467.296}{373.125}{ALIUSC0000FF}{\ALIUSFontCormorantRegular}{12.960}{http://www.oxfordscholarship.com/view/10.1093/0199271941.001.0001/acpr}
  \ALIUSPlacedTextContent{126.951}{483.136}{88.351}{ALIUSC0000FF}{\ALIUSFontCormorantRegular}{12.960}{of-9780199271948}
  \ALIUSPlacedTextContent{90.951}{504.736}{240.490}{ALIUSC000000}{\ALIUSFontCormorantRegular}{12.960}{Gallagher, S. (2013). A Pattern Theory of Self.}
  \ALIUSPlacedTextContent{332.682}{504.736}{170.716}{ALIUSC000000}{\ALIUSFontCormorantItalic}{12.960}{Frontiers in Human Neuroscience, 7}
  \ALIUSPlacedTextContent{503.390}{504.736}{5.594}{ALIUSC000000}{\ALIUSFontCormorantRegular}{12.960}{.}
  \ALIUSPlacedTextContent{126.951}{520.576}{206.472}{ALIUSC0000FF}{\ALIUSFontCormorantRegular}{12.960}{https://doi.org/10.3389/fnhum.2013.00443}
  \ALIUSPlacedTextContent{90.951}{542.176}{417.971}{ALIUSC000000}{\ALIUSFontCormorantRegular}{12.960}{Gentile, G., Björnsdotter, M., Petkova, V. I., Abdulkarim, Z., \& Ehrsson, H. H.}
  \ALIUSPlacedTextContent{126.951}{558.016}{381.954}{ALIUSC000000}{\ALIUSFontCormorantRegular}{12.960}{(2015). Patterns of neural activity in the human ventral premotor cortex}
  \ALIUSPlacedTextContent{126.951}{573.616}{238.065}{ALIUSC000000}{\ALIUSFontCormorantRegular}{12.960}{reflect a whole-body multisensory percept.}
  \ALIUSPlacedTextContent{369.936}{573.616}{82.141}{ALIUSC000000}{\ALIUSFontCormorantItalic}{12.960}{NeuroImage, 109}
  \ALIUSPlacedTextContent{452.107}{573.616}{56.866}{ALIUSC000000}{\ALIUSFontCormorantRegular}{12.960}{, 328–340.}
  \ALIUSPlacedTextContent{126.951}{589.456}{239.520}{ALIUSC0000FF}{\ALIUSFontCormorantRegular}{12.960}{https://doi.org/10.1016/j.neuroimage.2015.01.008}
  \ALIUSPlacedTextContent{90.951}{611.056}{418.064}{ALIUSC000000}{\ALIUSFontCormorantRegular}{12.960}{Heeter, C. (1992). Being There: The Subjective Experience of Presence. Presence:}
  \ALIUSPlacedTextContent{126.951}{626.896}{73.315}{ALIUSC000000}{\ALIUSFontCormorantRegular}{12.960}{Teleoperators}
  \ALIUSPlacedTextContent{220.794}{626.896}{21.864}{ALIUSC000000}{\ALIUSFontCormorantRegular}{12.960}{and}
  \ALIUSPlacedTextContent{263.186}{626.896}{39.567}{ALIUSC000000}{\ALIUSFontCormorantRegular}{12.960}{Virtual}
  \ALIUSPlacedTextContent{323.282}{626.896}{77.229}{ALIUSC000000}{\ALIUSFontCormorantRegular}{12.960}{Environments,}
  \ALIUSPlacedTextContent{421.039}{626.896}{23.497}{ALIUSC000000}{\ALIUSFontCormorantRegular}{12.960}{1(2),}
  \ALIUSPlacedTextContent{465.064}{626.896}{43.915}{ALIUSC000000}{\ALIUSFontCormorantRegular}{12.960}{262–271.}
  \ALIUSPlacedTextContent{126.951}{642.736}{194.219}{ALIUSC0000FF}{\ALIUSFontCormorantRegular}{12.960}{https://doi.org/10.1162/pres.1992.1.2.262}
  \ALIUSPlacedTextContent{90.951}{664.336}{417.983}{ALIUSC000000}{\ALIUSFontCormorantRegular}{12.960}{Held, R. M., \& Durlach, N. I. (1992). Telepresence. Presence: Teleoperators and}
  \ALIUSPlacedTextContent{126.951}{680.176}{182.119}{ALIUSC000000}{\ALIUSFontCormorantRegular}{12.960}{Virtual Environments, 1(1), 109–112.}
  \ALIUSPlacedTextContent{309.077}{680.176}{193.457}{ALIUSC0000FF}{\ALIUSFontCormorantRegular}{12.960}{https://doi.org/10.1162/pres.1992.1.1.109}
  \ALIUSPlacedTextContent{90.951}{701.776}{417.975}{ALIUSC000000}{\ALIUSFontCormorantRegular}{12.960}{Heydrich, L., Dodds, T., Aspell, J., Herbelin, B., Buelthoff, H., Mohler, B., \& Blanke,}
  \ALIUSPlacedTextContent{126.951}{717.616}{382.057}{ALIUSC000000}{\ALIUSFontCormorantRegular}{12.960}{O. (2013). Visual capture and the experience of having two bodies – Evidence}
  \ALIUSPlacedTextContent{126.951}{733.456}{245.014}{ALIUSC000000}{\ALIUSFontCormorantRegular}{12.960}{from two different virtual reality techniques.}
  \ALIUSPlacedTextContent{375.063}{733.456}{128.318}{ALIUSC000000}{\ALIUSFontCormorantItalic}{12.960}{Frontiers in Psychology, 4}
  \ALIUSPlacedTextContent{503.390}{733.456}{5.594}{ALIUSC000000}{\ALIUSFontCormorantRegular}{12.960}{.}
  \ALIUSPlacedTextContent{126.951}{749.296}{200.575}{ALIUSC0000FF}{\ALIUSFontCormorantRegular}{12.960}{https://doi.org/10.3389/fpsyg.2013.00946}
  \ALIUSPlacedTextContent{90.951}{793.143}{122.397}{ALIUSC767171}{\ALIUSFontCormorantLight}{12.000}{ALIUS Bulletin n°4 (2020)}
  \ALIUSPlacedTextContent{387.435}{793.143}{120.728}{ALIUSC767171}{\ALIUSFontCormorantLight}{12.000}{aliusresearch.org/bulletin}
\end{tikzpicture}
\null
\newpage

% Page 21
\thispagestyle{empty}
\noindent\begin{tikzpicture}[remember picture,overlay,x=1bp,y=1bp,shift={(current page.north west)}]
  \fill[ALIUSC0000FF] (242.391bp,-133.191bp) rectangle ++(215.760bp,-0.720bp);
  \fill[ALIUSC0000FF] (126.951bp,-186.471bp) rectangle ++(187.680bp,-0.720bp);
  \fill[ALIUSC0000FF] (126.951bp,-239.511bp) rectangle ++(190.320bp,-0.720bp);
  \fill[ALIUSC0000FF] (264.711bp,-292.791bp) rectangle ++(200.160bp,-0.720bp);
  \fill[ALIUSC0000FF] (269.271bp,-330.231bp) rectangle ++(174.240bp,-0.720bp);
  \fill[ALIUSC0000FF] (170.151bp,-399.351bp) rectangle ++(178.320bp,-0.720bp);
  \fill[ALIUSC0000FF] (126.951bp,-468.231bp) rectangle ++(200.640bp,-0.720bp);
  \fill[ALIUSC0000FF] (126.951bp,-537.111bp) rectangle ++(222.000bp,-0.720bp);
  \fill[ALIUSC0000FF] (126.951bp,-606.231bp) rectangle ++(206.400bp,-0.720bp);
  \fill[ALIUSC0000FF] (161.031bp,-681.111bp) rectangle ++(211.920bp,-0.720bp);
  \fill[ALIUSC0000FF] (126.951bp,-734.391bp) rectangle ++(211.200bp,-0.720bp);
  \ALIUSPlacedTextContent{90.951}{36.329}{258.472}{ALIUSC767171}{\ALIUSFontCormorantRegular}{13.920}{Lenggenhager – Bodily boundaries and beyond}
  \ALIUSPlacedTextContent{488.871}{36.329}{20.211}{ALIUSC000000}{\ALIUSFontCormorantRegular}{13.920}{102}
  \ALIUSPlacedTextContent{90.951}{72.976}{418.271}{ALIUSC000000}{\ALIUSFontCormorantRegular}{12.960}{Ionta, S., Heydrich, L., Lenggenhager, B., Mouthon, M., Fornari, E., Chapuis, D.,}
  \ALIUSPlacedTextContent{126.951}{88.576}{378.987}{ALIUSC000000}{\ALIUSFontCormorantRegular}{12.960}{Gassert, R., \& Blanke, O. (2011). Multisensory Mechanisms in Temporo-}
  \ALIUSPlacedTextContent{126.951}{104.416}{382.056}{ALIUSC000000}{\ALIUSFontCormorantRegular}{12.960}{Parietal Cortex Support Self-Location and First-Person Perspective.}
  \ALIUSPlacedTextContent{126.951}{120.256}{51.257}{ALIUSC000000}{\ALIUSFontCormorantItalic}{12.960}{Neuron, 70}
  \ALIUSPlacedTextContent{178.234}{120.256}{64.122}{ALIUSC000000}{\ALIUSFontCormorantRegular}{12.960}{(2), 363–374.}
  \ALIUSPlacedTextContent{242.375}{120.256}{215.661}{ALIUSC0000FF}{\ALIUSFontCormorantRegular}{12.960}{https://doi.org/10.1016/j.neuron.2011.03.009}
  \ALIUSPlacedTextContent{90.951}{141.856}{418.046}{ALIUSC000000}{\ALIUSFontCormorantRegular}{12.960}{Lackner, J. R. (1988). Some proprioceptive influences on the perceptual}
  \ALIUSPlacedTextContent{126.951}{157.696}{259.453}{ALIUSC000000}{\ALIUSFontCormorantRegular}{12.960}{representation of body shape and orientation.}
  \ALIUSPlacedTextContent{391.363}{157.696}{48.372}{ALIUSC000000}{\ALIUSFontCormorantItalic}{12.960}{Brain, 111}
  \ALIUSPlacedTextContent{439.758}{157.696}{69.215}{ALIUSC000000}{\ALIUSFontCormorantRegular}{12.960}{(2), 281–297.}
  \ALIUSPlacedTextContent{126.951}{173.536}{187.648}{ALIUSC0000FF}{\ALIUSFontCormorantRegular}{12.960}{https://doi.org/10.1093/brain/111.2.281}
  \ALIUSPlacedTextContent{90.951}{195.136}{418.070}{ALIUSC000000}{\ALIUSFontCormorantRegular}{12.960}{Lenggenhager, B., Tadi, T., Metzinger, T., \& Blanke, O. (2007). Video Ergo Sum:}
  \ALIUSPlacedTextContent{126.951}{210.736}{220.179}{ALIUSC000000}{\ALIUSFontCormorantRegular}{12.960}{Manipulating Bodily Self-Consciousness.}
  \ALIUSPlacedTextContent{352.548}{210.736}{57.742}{ALIUSC000000}{\ALIUSFontCormorantItalic}{12.960}{Science, 317}
  \ALIUSPlacedTextContent{410.289}{210.736}{98.695}{ALIUSC000000}{\ALIUSFontCormorantRegular}{12.960}{(5841), 1096–1099.}
  \ALIUSPlacedTextContent{126.951}{226.576}{190.266}{ALIUSC0000FF}{\ALIUSFontCormorantRegular}{12.960}{https://doi.org/10.1126/science.1143439}
  \ALIUSPlacedTextContent{90.951}{248.416}{417.902}{ALIUSC000000}{\ALIUSFontCormorantRegular}{12.960}{Lesur, M. R., Bolt, E., \& Lenggenhager, B. (2020). The monologue of the double:}
  \ALIUSPlacedTextContent{126.951}{264.016}{382.082}{ALIUSC000000}{\ALIUSFontCormorantRegular}{12.960}{Allocentric reduplication of the own voice alters bodily self perception.}
  \ALIUSPlacedTextContent{126.951}{279.856}{137.738}{ALIUSC000000}{\ALIUSFontCormorantRegular}{12.960}{BioRxiv, 2020.08.11.246397.}
  \ALIUSPlacedTextContent{264.695}{279.856}{200.088}{ALIUSC0000FF}{\ALIUSFontCormorantRegular}{12.960}{https://doi.org/10.1101/2020.08.11.246397}
  \ALIUSPlacedTextContent{90.951}{301.456}{418.017}{ALIUSC000000}{\ALIUSFontCormorantRegular}{12.960}{Lush, P. (2020). Demand Characteristics Confound the Rubber Hand Illusion.}
  \ALIUSPlacedTextContent{126.951}{317.296}{107.887}{ALIUSC000000}{\ALIUSFontCormorantItalic}{12.960}{Collabra: Psychology, 6}
  \ALIUSPlacedTextContent{234.873}{317.296}{34.321}{ALIUSC000000}{\ALIUSFontCormorantRegular}{12.960}{(1), 22.}
  \ALIUSPlacedTextContent{269.205}{317.296}{174.228}{ALIUSC0000FF}{\ALIUSFontCormorantRegular}{12.960}{https://doi.org/10.1525/collabra.325}
  \ALIUSPlacedTextContent{90.951}{339.136}{418.102}{ALIUSC000000}{\ALIUSFontCormorantRegular}{12.960}{Lush, P., Botan, V., Scott, R. B., Seth, A., Ward, J., \& Dienes, Z. (2019).}
  \ALIUSPlacedTextContent{126.951}{354.736}{382.142}{ALIUSC000000}{\ALIUSFontCormorantRegular}{12.960}{Phenomenological control: Response to imaginative suggestion predicts}
  \ALIUSPlacedTextContent{126.951}{370.576}{382.038}{ALIUSC000000}{\ALIUSFontCormorantRegular}{12.960}{measures of mirror touch synaesthesia, vicarious pain and the rubber hand}
  \ALIUSPlacedTextContent{126.951}{386.416}{43.159}{ALIUSC000000}{\ALIUSFontCormorantRegular}{12.960}{illusion.}
  \ALIUSPlacedTextContent{170.109}{386.416}{178.213}{ALIUSC0000FF}{\ALIUSFontCormorantRegular}{12.960}{https://doi.org/10.31234/osf.io/82jav}
  \ALIUSPlacedTextContent{90.951}{408.016}{417.963}{ALIUSC000000}{\ALIUSFontCormorantRegular}{12.960}{Madary, M., \& Metzinger, T. K. (2016). Real virtuality: a code of ethical conduct.}
  \ALIUSPlacedTextContent{126.951}{423.856}{378.987}{ALIUSC000000}{\ALIUSFontCormorantRegular}{12.960}{Recommendations for good scientific practice and the consumers of VR-}
  \ALIUSPlacedTextContent{126.951}{439.456}{238.587}{ALIUSC000000}{\ALIUSFontCormorantRegular}{12.960}{technology. Frontiers in Robotics and AI, 3, 3.}
  \ALIUSPlacedTextContent{126.951}{455.296}{200.575}{ALIUSC0000FF}{\ALIUSFontCormorantRegular}{12.960}{https://doi.org/10.3389/frobt.2016.00003}
  \ALIUSPlacedTextContent{90.951}{476.896}{418.055}{ALIUSC000000}{\ALIUSFontCormorantRegular}{12.960}{Marotta, A., Tinazzi, M., Cavedini, C., Zampini, M., \& Fiorio, M. (2016). Individual}
  \ALIUSPlacedTextContent{126.951}{492.736}{381.908}{ALIUSC000000}{\ALIUSFontCormorantRegular}{12.960}{Differences in the Rubber Hand Illusion Are Related to Sensory}
  \ALIUSPlacedTextContent{126.951}{508.336}{74.695}{ALIUSC000000}{\ALIUSFontCormorantRegular}{12.960}{Suggestibility.}
  \ALIUSPlacedTextContent{201.685}{508.336}{67.619}{ALIUSC000000}{\ALIUSFontCormorantItalic}{12.960}{PLOS ONE, 11}
  \ALIUSPlacedTextContent{269.310}{508.336}{78.626}{ALIUSC000000}{\ALIUSFontCormorantRegular}{12.960}{(12), e0168489.}
  \ALIUSPlacedTextContent{126.951}{524.176}{222.024}{ALIUSC0000FF}{\ALIUSFontCormorantRegular}{12.960}{https://doi.org/10.1371/journal.pone.0168489}
  \ALIUSPlacedTextContent{90.951}{546.016}{414.987}{ALIUSC000000}{\ALIUSFontCormorantRegular}{12.960}{Maselli, A. (2015). Allocentric and egocentric manipulations of the sense of self-}
  \ALIUSPlacedTextContent{126.951}{561.616}{382.019}{ALIUSC000000}{\ALIUSFontCormorantRegular}{12.960}{location in full-body illusions and their relation with the sense of body}
  \ALIUSPlacedTextContent{126.951}{577.456}{58.372}{ALIUSC000000}{\ALIUSFontCormorantRegular}{12.960}{ownership.}
  \ALIUSPlacedTextContent{185.332}{577.456}{109.401}{ALIUSC000000}{\ALIUSFontCormorantItalic}{12.960}{Cognitive Processing, 16}
  \ALIUSPlacedTextContent{294.777}{577.456}{65.509}{ALIUSC000000}{\ALIUSFontCormorantRegular}{12.960}{(1), 309–312.}
  \ALIUSPlacedTextContent{126.951}{593.296}{206.328}{ALIUSC0000FF}{\ALIUSFontCormorantRegular}{12.960}{https://doi.org/10.1007/s10339-015-0667-z}
  \ALIUSPlacedTextContent{90.951}{614.896}{353.667}{ALIUSC000000}{\ALIUSFontCormorantRegular}{12.960}{Millière, R. (2019). Are there degrees of self-consciousness?}
  \ALIUSPlacedTextContent{452.456}{614.896}{56.506}{ALIUSC000000}{\ALIUSFontCormorantItalic}{12.960}{Journal of}
  \ALIUSPlacedTextContent{126.951}{630.736}{116.370}{ALIUSC000000}{\ALIUSFontCormorantItalic}{12.960}{Consciousness Studies, 26}
  \ALIUSPlacedTextContent{243.337}{630.736}{76.481}{ALIUSC000000}{\ALIUSFontCormorantRegular}{12.960}{(3–4), 252–276.}
  \ALIUSPlacedTextContent{90.951}{652.336}{248.905}{ALIUSC000000}{\ALIUSFontCormorantRegular}{12.960}{Millière, R. (2020). The varieties of selflessness.}
  \ALIUSPlacedTextContent{341.548}{652.336}{167.489}{ALIUSC000000}{\ALIUSFontCormorantItalic}{12.960}{Philosophy and the Mind Sciences,}
  \ALIUSPlacedTextContent{126.951}{668.176}{3.771}{ALIUSC000000}{\ALIUSFontCormorantItalic}{12.960}{1}
  \ALIUSPlacedTextContent{130.733}{668.176}{30.240}{ALIUSC000000}{\ALIUSFontCormorantRegular}{12.960}{(I), 8.}
  \ALIUSPlacedTextContent{160.971}{668.176}{211.909}{ALIUSC0000FF}{\ALIUSFontCormorantRegular}{12.960}{https://doi.org/10.33735/phimisci.2020.I.48}
  \ALIUSPlacedTextContent{90.951}{689.776}{418.023}{ALIUSC000000}{\ALIUSFontCormorantRegular}{12.960}{Millière, R., \& Metzinger, T. (2020). Radical disruptions of self-consciousness.}
  \ALIUSPlacedTextContent{126.951}{705.616}{164.786}{ALIUSC000000}{\ALIUSFontCormorantItalic}{12.960}{Philosophy and the Mind Sciences, 1}
  \ALIUSPlacedTextContent{291.735}{705.616}{42.240}{ALIUSC000000}{\ALIUSFontCormorantRegular}{12.960}{(I), 1–1.}
  \ALIUSPlacedTextContent{126.951}{721.456}{211.209}{ALIUSC0000FF}{\ALIUSFontCormorantRegular}{12.960}{https://doi.org/10.33735/phimisci.2020.I.50}
  \ALIUSPlacedTextContent{90.951}{793.143}{122.397}{ALIUSC767171}{\ALIUSFontCormorantLight}{12.000}{ALIUS Bulletin n°4 (2020)}
  \ALIUSPlacedTextContent{387.435}{793.143}{120.728}{ALIUSC767171}{\ALIUSFontCormorantLight}{12.000}{aliusresearch.org/bulletin}
\end{tikzpicture}
\null
\newpage

% Page 22
\thispagestyle{empty}
\noindent\begin{tikzpicture}[remember picture,overlay,x=1bp,y=1bp,shift={(current page.north west)}]
  \fill[ALIUSC0000FF] (126.951bp,-208.071bp) rectangle ++(213.600bp,-0.720bp);
  \fill[ALIUSC0000FF] (126.951bp,-261.351bp) rectangle ++(199.920bp,-0.720bp);
  \fill[ALIUSC0000FF] (126.951bp,-314.631bp) rectangle ++(227.040bp,-0.720bp);
  \fill[ALIUSC0000FF] (152.631bp,-367.911bp) rectangle ++(214.320bp,-0.720bp);
  \fill[ALIUSC0000FF] (323.991bp,-436.791bp) rectangle ++(177.360bp,-0.720bp);
  \fill[ALIUSC0000FF] (126.951bp,-490.071bp) rectangle ++(152.160bp,-0.720bp);
  \fill[ALIUSC0000FF] (126.951bp,-543.111bp) rectangle ++(198.480bp,-0.720bp);
  \fill[ALIUSC0000FF] (277.911bp,-596.391bp) rectangle ++(170.880bp,-0.720bp);
  \fill[ALIUSC0000FF] (126.951bp,-665.511bp) rectangle ++(209.040bp,-0.720bp);
  \fill[ALIUSC0000FF] (296.151bp,-734.391bp) rectangle ++(190.080bp,-0.720bp);
  \ALIUSPlacedTextContent{90.951}{36.329}{258.472}{ALIUSC767171}{\ALIUSFontCormorantRegular}{13.920}{Lenggenhager – Bodily boundaries and beyond}
  \ALIUSPlacedTextContent{489.111}{36.329}{20.029}{ALIUSC000000}{\ALIUSFontCormorantRegular}{13.920}{103}
  \ALIUSPlacedTextContent{90.951}{72.976}{417.963}{ALIUSC000000}{\ALIUSFontCormorantRegular}{12.960}{Mizumoto, M., \& Ishikawa, M. (2005). Immunity to Error through}
  \ALIUSPlacedTextContent{126.951}{88.576}{316.403}{ALIUSC000000}{\ALIUSFontCormorantRegular}{12.960}{Misidentification and the Bodily Illusion Experiment.}
  \ALIUSPlacedTextContent{451.845}{88.576}{57.115}{ALIUSC000000}{\ALIUSFontCormorantItalic}{12.960}{Journal of}
  \ALIUSPlacedTextContent{126.951}{104.416}{114.695}{ALIUSC000000}{\ALIUSFontCormorantItalic}{12.960}{Consciousness Studies, 12}
  \ALIUSPlacedTextContent{241.660}{104.416}{47.245}{ALIUSC000000}{\ALIUSFontCormorantRegular}{12.960}{(7), 3–19.}
  \ALIUSPlacedTextContent{90.951}{126.016}{358.738}{ALIUSC000000}{\ALIUSFontCormorantRegular}{12.960}{Müller-Lyer, F. (1889). Optische Urteilstäuschungen, du Bois Arch.}
  \ALIUSPlacedTextContent{452.283}{126.016}{56.700}{ALIUSC000000}{\ALIUSFontCormorantItalic}{12.960}{Archiv Für}
  \ALIUSPlacedTextContent{126.951}{141.856}{103.136}{ALIUSC000000}{\ALIUSFontCormorantItalic}{12.960}{Physiologie Suppl., 263}
  \ALIUSPlacedTextContent{230.089}{141.856}{51.548}{ALIUSC000000}{\ALIUSFontCormorantRegular}{12.960}{, 263–270.}
  \ALIUSPlacedTextContent{90.951}{163.696}{418.031}{ALIUSC000000}{\ALIUSFontCormorantRegular}{12.960}{Nakul, E., Orlando-Dessaints, N., Lenggenhager, B., \& Lopez, C. (2020). Measuring}
  \ALIUSPlacedTextContent{126.951}{179.296}{221.737}{ALIUSC000000}{\ALIUSFontCormorantRegular}{12.960}{perceived self-location in virtual reality.}
  \ALIUSPlacedTextContent{352.821}{179.296}{104.294}{ALIUSC000000}{\ALIUSFontCormorantItalic}{12.960}{Scientific Reports, 10}
  \ALIUSPlacedTextContent{457.100}{179.296}{51.892}{ALIUSC000000}{\ALIUSFontCormorantRegular}{12.960}{(1), 6802.}
  \ALIUSPlacedTextContent{126.951}{195.136}{213.567}{ALIUSC0000FF}{\ALIUSFontCormorantRegular}{12.960}{https://doi.org/10.1038/s41598-020-63643-y}
  \ALIUSPlacedTextContent{90.951}{216.736}{417.978}{ALIUSC000000}{\ALIUSFontCormorantRegular}{12.960}{Newen, A. (2018). The Embodied Self, the Pattern Theory of Self, and the}
  \ALIUSPlacedTextContent{126.951}{232.576}{87.975}{ALIUSC000000}{\ALIUSFontCormorantRegular}{12.960}{Predictive Mind.}
  \ALIUSPlacedTextContent{214.919}{232.576}{119.096}{ALIUSC000000}{\ALIUSFontCormorantItalic}{12.960}{Frontiers in Psychology, 9}
  \ALIUSPlacedTextContent{334.048}{232.576}{8.635}{ALIUSC000000}{\ALIUSFontCormorantRegular}{12.960}{.}
  \ALIUSPlacedTextContent{126.951}{248.416}{199.953}{ALIUSC0000FF}{\ALIUSFontCormorantRegular}{12.960}{https://doi.org/10.3389/fpsyg.2018.02270}
  \ALIUSPlacedTextContent{90.951}{270.016}{418.087}{ALIUSC000000}{\ALIUSFontCormorantRegular}{12.960}{Noel, J.-P., Pfeiffer, C., Blanke, O., \& Serino, A. (2015). Peripersonal space as the}
  \ALIUSPlacedTextContent{126.951}{285.856}{120.398}{ALIUSC000000}{\ALIUSFontCormorantRegular}{12.960}{space of the bodily self.}
  \ALIUSPlacedTextContent{247.314}{285.856}{65.946}{ALIUSC000000}{\ALIUSFontCormorantItalic}{12.960}{Cognition, 144}
  \ALIUSPlacedTextContent{313.287}{285.856}{44.060}{ALIUSC000000}{\ALIUSFontCormorantRegular}{12.960}{, 49–57.}
  \ALIUSPlacedTextContent{126.951}{301.696}{226.865}{ALIUSC0000FF}{\ALIUSFontCormorantRegular}{12.960}{https://doi.org/10.1016/j.cognition.2015.07.012}
  \ALIUSPlacedTextContent{90.951}{323.296}{418.019}{ALIUSC000000}{\ALIUSFontCormorantRegular}{12.960}{Occhionero, M., \& Cicogna, P. C. (2011). Autoscopic phenomena and one's own}
  \ALIUSPlacedTextContent{126.951}{339.136}{168.001}{ALIUSC000000}{\ALIUSFontCormorantRegular}{12.960}{body representation in dreams.}
  \ALIUSPlacedTextContent{297.497}{339.136}{156.622}{ALIUSC000000}{\ALIUSFontCormorantItalic}{12.960}{Consciousness and Cognition, 20}
  \ALIUSPlacedTextContent{454.152}{339.136}{51.778}{ALIUSC000000}{\ALIUSFontCormorantRegular}{12.960}{(4), 1009–}
  \ALIUSPlacedTextContent{126.951}{354.976}{25.752}{ALIUSC000000}{\ALIUSFontCormorantRegular}{12.960}{1015.}
  \ALIUSPlacedTextContent{152.703}{354.976}{214.261}{ALIUSC0000FF}{\ALIUSFontCormorantRegular}{12.960}{https://doi.org/10.1016/j.concog.2011.01.004}
  \ALIUSPlacedTextContent{90.951}{376.576}{417.963}{ALIUSC000000}{\ALIUSFontCormorantRegular}{12.960}{Oliveira, E. C. D., Bertrand, P., Lesur, M. E. R., Palomo, P., Demarzo, M., Cebolla,}
  \ALIUSPlacedTextContent{126.951}{392.176}{382.036}{ALIUSC000000}{\ALIUSFontCormorantRegular}{12.960}{A., Baños, R., \& Tori, R. (2016). Virtual Body Swap: A New Feasible Tool to}
  \ALIUSPlacedTextContent{126.951}{408.016}{382.133}{ALIUSC000000}{\ALIUSFontCormorantRegular}{12.960}{Be Explored in Health and Education. 2016 XVIII Symposium on Virtual}
  \ALIUSPlacedTextContent{126.951}{423.856}{197.131}{ALIUSC000000}{\ALIUSFontCormorantRegular}{12.960}{and Augmented Reality (SVR), 81–89.}
  \ALIUSPlacedTextContent{324.091}{423.856}{177.202}{ALIUSC0000FF}{\ALIUSFontCormorantRegular}{12.960}{https://doi.org/10.1109/SVR.2016.23}
  \ALIUSPlacedTextContent{90.951}{445.456}{418.013}{ALIUSC000000}{\ALIUSFontCormorantRegular}{12.960}{Pamment, J., \& Aspell, J. E. (2017). Putting pain out of mind with an 'out of body'}
  \ALIUSPlacedTextContent{126.951}{461.296}{43.159}{ALIUSC000000}{\ALIUSFontCormorantRegular}{12.960}{illusion.}
  \ALIUSPlacedTextContent{194.088}{461.296}{46.796}{ALIUSC000000}{\ALIUSFontCormorantItalic}{12.960}{European}
  \ALIUSPlacedTextContent{264.868}{461.296}{37.081}{ALIUSC000000}{\ALIUSFontCormorantItalic}{12.960}{Journal}
  \ALIUSPlacedTextContent{325.933}{461.296}{11.615}{ALIUSC000000}{\ALIUSFontCormorantItalic}{12.960}{of}
  \ALIUSPlacedTextContent{361.531}{461.296}{27.385}{ALIUSC000000}{\ALIUSFontCormorantItalic}{12.960}{Pain,}
  \ALIUSPlacedTextContent{412.900}{461.296}{8.484}{ALIUSC000000}{\ALIUSFontCormorantItalic}{12.960}{21}
  \ALIUSPlacedTextContent{421.406}{461.296}{19.176}{ALIUSC000000}{\ALIUSFontCormorantRegular}{12.960}{(2),}
  \ALIUSPlacedTextContent{464.569}{461.296}{44.406}{ALIUSC000000}{\ALIUSFontCormorantRegular}{12.960}{334–342.}
  \ALIUSPlacedTextContent{126.951}{477.136}{151.982}{ALIUSC0000FF}{\ALIUSFontCormorantRegular}{12.960}{https://doi.org/10.1002/ejp.927}
  \ALIUSPlacedTextContent{90.951}{498.736}{414.987}{ALIUSC000000}{\ALIUSFontCormorantRegular}{12.960}{Park, H.-D., \& Blanke, O. (2019). Coupling Inner and Outer Body for Self-}
  \ALIUSPlacedTextContent{126.951}{514.336}{78.556}{ALIUSC000000}{\ALIUSFontCormorantRegular}{12.960}{Consciousness.}
  \ALIUSPlacedTextContent{220.781}{514.336}{34.839}{ALIUSC000000}{\ALIUSFontCormorantItalic}{12.960}{Trends}
  \ALIUSPlacedTextContent{270.870}{514.336}{12.291}{ALIUSC000000}{\ALIUSFontCormorantItalic}{12.960}{in}
  \ALIUSPlacedTextContent{298.411}{514.336}{47.114}{ALIUSC000000}{\ALIUSFontCormorantItalic}{12.960}{Cognitive}
  \ALIUSPlacedTextContent{360.775}{514.336}{43.005}{ALIUSC000000}{\ALIUSFontCormorantItalic}{12.960}{Sciences,}
  \ALIUSPlacedTextContent{419.030}{514.336}{9.146}{ALIUSC000000}{\ALIUSFontCormorantItalic}{12.960}{23}
  \ALIUSPlacedTextContent{428.200}{514.336}{19.266}{ALIUSC000000}{\ALIUSFontCormorantRegular}{12.960}{(5),}
  \ALIUSPlacedTextContent{462.706}{514.336}{46.264}{ALIUSC000000}{\ALIUSFontCormorantRegular}{12.960}{377–388.}
  \ALIUSPlacedTextContent{126.951}{530.176}{198.294}{ALIUSC0000FF}{\ALIUSFontCormorantRegular}{12.960}{https://doi.org/10.1016/j.tics.2019.02.002}
  \ALIUSPlacedTextContent{90.951}{552.016}{418.102}{ALIUSC000000}{\ALIUSFontCormorantRegular}{12.960}{Pfeiffer, C., Grivaz, P., Herbelin, B., Serino, A., \& Blanke, O. (2016). Visual gravity}
  \ALIUSPlacedTextContent{126.951}{567.616}{289.581}{ALIUSC000000}{\ALIUSFontCormorantRegular}{12.960}{contributes to subjective first-person perspective.}
  \ALIUSPlacedTextContent{425.624}{567.616}{83.346}{ALIUSC000000}{\ALIUSFontCormorantItalic}{12.960}{Neuroscience of}
  \ALIUSPlacedTextContent{126.951}{583.456}{89.533}{ALIUSC000000}{\ALIUSFontCormorantItalic}{12.960}{Consciousness, 2016}
  \ALIUSPlacedTextContent{216.518}{583.456}{61.324}{ALIUSC000000}{\ALIUSFontCormorantRegular}{12.960}{(1), niw006.}
  \ALIUSPlacedTextContent{277.838}{583.456}{170.985}{ALIUSC0000FF}{\ALIUSFontCormorantRegular}{12.960}{https://doi.org/10.1093/nc/niw006}
  \ALIUSPlacedTextContent{90.951}{605.056}{418.065}{ALIUSC000000}{\ALIUSFontCormorantRegular}{12.960}{Pfeiffer, C., Schmutz, V., \& Blanke, O. (2014). Visuospatial viewpoint manipulation}
  \ALIUSPlacedTextContent{126.951}{620.896}{382.035}{ALIUSC000000}{\ALIUSFontCormorantRegular}{12.960}{during full-body illusion modulates subjective first-person perspective.}
  \ALIUSPlacedTextContent{126.951}{636.736}{154.491}{ALIUSC000000}{\ALIUSFontCormorantItalic}{12.960}{Experimental Brain Research, 232}
  \ALIUSPlacedTextContent{281.505}{636.736}{82.617}{ALIUSC000000}{\ALIUSFontCormorantRegular}{12.960}{(12), 4021–4033.}
  \ALIUSPlacedTextContent{126.951}{652.576}{209.055}{ALIUSC0000FF}{\ALIUSFontCormorantRegular}{12.960}{https://doi.org/10.1007/s00221-014-4080-0}
  \ALIUSPlacedTextContent{90.951}{674.176}{418.083}{ALIUSC000000}{\ALIUSFontCormorantRegular}{12.960}{Preller, K. H., \& Vollenweider, F. X. (2018). Phenomenology, Structure, and}
  \ALIUSPlacedTextContent{126.951}{689.776}{381.910}{ALIUSC000000}{\ALIUSFontCormorantRegular}{12.960}{Dynamic of Psychedelic States. In A. L. Halberstadt, F. X. Vollenweider, \&}
  \ALIUSPlacedTextContent{126.951}{705.616}{378.979}{ALIUSC000000}{\ALIUSFontCormorantRegular}{12.960}{D. E. Nichols (Eds.), Behavioral Neurobiology of Psychedelic Drugs (pp. 221–}
  \ALIUSPlacedTextContent{126.951}{721.456}{169.097}{ALIUSC000000}{\ALIUSFontCormorantRegular}{12.960}{256). Springer Berlin Heidelberg.}
  \ALIUSPlacedTextContent{296.090}{721.456}{190.123}{ALIUSC0000FF}{\ALIUSFontCormorantRegular}{12.960}{https://doi.org/10.1007/7854\_2016\_459}
  \ALIUSPlacedTextContent{90.951}{793.143}{122.397}{ALIUSC767171}{\ALIUSFontCormorantLight}{12.000}{ALIUS Bulletin n°4 (2020)}
  \ALIUSPlacedTextContent{387.435}{793.143}{120.728}{ALIUSC767171}{\ALIUSFontCormorantLight}{12.000}{aliusresearch.org/bulletin}
\end{tikzpicture}
\null
\newpage

% Page 23
\thispagestyle{empty}
\noindent\begin{tikzpicture}[remember picture,overlay,x=1bp,y=1bp,shift={(current page.north west)}]
  \fill[ALIUSC0000FF] (126.951bp,-133.191bp) rectangle ++(218.640bp,-0.720bp);
  \fill[ALIUSC0000FF] (126.951bp,-202.071bp) rectangle ++(208.560bp,-0.720bp);
  \fill[ALIUSC0000FF] (126.951bp,-393.351bp) rectangle ++(200.400bp,-0.720bp);
  \fill[ALIUSC0000FF] (268.311bp,-446.631bp) rectangle ++(203.040bp,-0.720bp);
  \fill[ALIUSC0000FF] (126.951bp,-515.511bp) rectangle ++(214.560bp,-0.720bp);
  \fill[ALIUSC0000FF] (268.791bp,-584.391bp) rectangle ++(196.320bp,-0.720bp);
  \fill[ALIUSC0000FF] (126.951bp,-653.511bp) rectangle ++(189.120bp,-0.720bp);
  \ALIUSPlacedTextContent{90.951}{36.329}{258.472}{ALIUSC767171}{\ALIUSFontCormorantRegular}{13.920}{Lenggenhager – Bodily boundaries and beyond}
  \ALIUSPlacedTextContent{488.391}{36.329}{20.897}{ALIUSC000000}{\ALIUSFontCormorantRegular}{13.920}{104}
  \ALIUSPlacedTextContent{90.951}{72.976}{418.050}{ALIUSC000000}{\ALIUSFontCormorantRegular}{12.960}{Riva, G., Wiederhold, B. K., \& Mantovani, F. (2018). Neuroscience of Virtual}
  \ALIUSPlacedTextContent{126.951}{88.576}{296.540}{ALIUSC000000}{\ALIUSFontCormorantRegular}{12.960}{Reality: From Virtual Exposure to Embodied Medicine.}
  \ALIUSPlacedTextContent{425.782}{88.576}{83.188}{ALIUSC000000}{\ALIUSFontCormorantItalic}{12.960}{Cyberpsychology,}
  \ALIUSPlacedTextContent{126.951}{104.416}{171.163}{ALIUSC000000}{\ALIUSFontCormorantItalic}{12.960}{Behavior, and Social Networking, 22(}
  \ALIUSPlacedTextContent{298.144}{104.416}{53.264}{ALIUSC000000}{\ALIUSFontCormorantRegular}{12.960}{1), 82–96.}
  \ALIUSPlacedTextContent{126.951}{120.256}{218.577}{ALIUSC0000FF}{\ALIUSFontCormorantRegular}{12.960}{https://doi.org/10.1089/cyber.2017.29099.gri}
  \ALIUSPlacedTextContent{90.951}{141.856}{418.140}{ALIUSC000000}{\ALIUSFontCormorantRegular}{12.960}{Roel Lesur, M., Aicher, H., Delplanque, S., \& Lenggenhager, B. (2020). Being Short,}
  \ALIUSPlacedTextContent{126.951}{157.696}{382.079}{ALIUSC000000}{\ALIUSFontCormorantRegular}{12.960}{Sweet, and Sour: Congruent Visuo-Olfactory Stimulation Enhances Illusory}
  \ALIUSPlacedTextContent{126.951}{173.296}{71.911}{ALIUSC000000}{\ALIUSFontCormorantRegular}{12.960}{Embodiment.}
  \ALIUSPlacedTextContent{198.852}{173.296}{64.819}{ALIUSC000000}{\ALIUSFontCormorantItalic}{12.960}{Perception, 49}
  \ALIUSPlacedTextContent{263.682}{173.296}{70.618}{ALIUSC000000}{\ALIUSFontCormorantRegular}{12.960}{(6), 693–696.}
  \ALIUSPlacedTextContent{126.951}{189.136}{208.585}{ALIUSC0000FF}{\ALIUSFontCormorantRegular}{12.960}{https://doi.org/10.1177/0301006620928669}
  \ALIUSPlacedTextContent{90.951}{210.736}{418.024}{ALIUSC000000}{\ALIUSFontCormorantRegular}{12.960}{Roel Lesur, M., Gaebler, M., Bertrand, P., \& Lenggenhager, B. (2018a). Authors'}
  \ALIUSPlacedTextContent{126.951}{226.576}{381.929}{ALIUSC000000}{\ALIUSFontCormorantRegular}{12.960}{Response: On the Components and Future Experimental Setups of Bodily}
  \ALIUSPlacedTextContent{126.951}{242.416}{82.439}{ALIUSC000000}{\ALIUSFontCormorantRegular}{12.960}{Illusions/Aliefs.}
  \ALIUSPlacedTextContent{209.394}{242.416}{140.580}{ALIUSC000000}{\ALIUSFontCormorantItalic}{12.960}{Constructivist Foundations, 14}
  \ALIUSPlacedTextContent{349.999}{242.416}{57.202}{ALIUSC000000}{\ALIUSFontCormorantRegular}{12.960}{(1), 111–113.}
  \ALIUSPlacedTextContent{90.951}{264.016}{418.017}{ALIUSC000000}{\ALIUSFontCormorantRegular}{12.960}{Roel Lesur, M., Gaebler, M., Bertrand, P., \& Lenggenhager, B. (2018b). The}
  \ALIUSPlacedTextContent{126.951}{279.856}{382.041}{ALIUSC000000}{\ALIUSFontCormorantRegular}{12.960}{Plasticity of the Bodily Self: Head Movements in Bodily Illusions and Their}
  \ALIUSPlacedTextContent{126.951}{295.456}{382.024}{ALIUSC000000}{\ALIUSFontCormorantRegular}{12.960}{Relation to Gallagher's Body Image and Body Schema. Constructivist}
  \ALIUSPlacedTextContent{126.951}{311.296}{137.774}{ALIUSC000000}{\ALIUSFontCormorantRegular}{12.960}{Foundations, 14(1), 94–105.}
  \ALIUSPlacedTextContent{90.951}{333.136}{417.860}{ALIUSC000000}{\ALIUSFontCormorantRegular}{12.960}{Roel Lesur, M., Weijs, M. L., Simon, C., Kannape, O. A., \& Lenggenhager, B. (2020).}
  \ALIUSPlacedTextContent{126.951}{348.736}{382.067}{ALIUSC000000}{\ALIUSFontCormorantRegular}{12.960}{Psychometrics of Disembodiment and Its Differential Modulation by}
  \ALIUSPlacedTextContent{126.951}{364.576}{243.595}{ALIUSC000000}{\ALIUSFontCormorantRegular}{12.960}{Visuomotor and Visuotactile Mismatches.}
  \ALIUSPlacedTextContent{379.998}{364.576}{61.649}{ALIUSC000000}{\ALIUSFontCormorantItalic}{12.960}{IScience, 23}
  \ALIUSPlacedTextContent{441.681}{364.576}{67.312}{ALIUSC000000}{\ALIUSFontCormorantRegular}{12.960}{(3), 100901.}
  \ALIUSPlacedTextContent{126.951}{380.416}{200.187}{ALIUSC0000FF}{\ALIUSFontCormorantRegular}{12.960}{https://doi.org/10.1016/j.isci.2020.100901}
  \ALIUSPlacedTextContent{90.951}{402.016}{417.941}{ALIUSC000000}{\ALIUSFontCormorantRegular}{12.960}{Salomon, R., Lim, M., Pfeiffer, C., Gassert, R., \& Blanke, O. (2013). Full body}
  \ALIUSPlacedTextContent{126.951}{417.856}{337.214}{ALIUSC000000}{\ALIUSFontCormorantRegular}{12.960}{illusion is associated with widespread skin temperature reduction.}
  \ALIUSPlacedTextContent{464.183}{417.856}{44.783}{ALIUSC000000}{\ALIUSFontCormorantItalic}{12.960}{Frontiers}
  \ALIUSPlacedTextContent{126.951}{433.696}{135.748}{ALIUSC000000}{\ALIUSFontCormorantItalic}{12.960}{in Behavioral Neuroscience, 7}
  \ALIUSPlacedTextContent{262.707}{433.696}{5.594}{ALIUSC000000}{\ALIUSFontCormorantRegular}{12.960}{.}
  \ALIUSPlacedTextContent{268.310}{433.696}{202.857}{ALIUSC0000FF}{\ALIUSFontCormorantRegular}{12.960}{https://doi.org/10.3389/fnbeh.2013.00065}
  \ALIUSPlacedTextContent{90.951}{455.296}{418.048}{ALIUSC000000}{\ALIUSFontCormorantRegular}{12.960}{Serino, A., Alsmith, A., Costantini, M., Mandrigin, A., Tajadura-Jimenez, A., \&}
  \ALIUSPlacedTextContent{126.951}{470.896}{382.077}{ALIUSC000000}{\ALIUSFontCormorantRegular}{12.960}{Lopez, C. (2013). Bodily ownership and self-location: Components of bodily}
  \ALIUSPlacedTextContent{126.951}{486.736}{96.217}{ALIUSC000000}{\ALIUSFontCormorantRegular}{12.960}{self-consciousness.}
  \ALIUSPlacedTextContent{236.151}{486.736}{66.983}{ALIUSC000000}{\ALIUSFontCormorantItalic}{12.960}{Consciousness}
  \ALIUSPlacedTextContent{316.104}{486.736}{20.495}{ALIUSC000000}{\ALIUSFontCormorantItalic}{12.960}{and}
  \ALIUSPlacedTextContent{349.569}{486.736}{51.179}{ALIUSC000000}{\ALIUSFontCormorantItalic}{12.960}{Cognition,}
  \ALIUSPlacedTextContent{413.719}{486.736}{9.419}{ALIUSC000000}{\ALIUSFontCormorantItalic}{12.960}{22}
  \ALIUSPlacedTextContent{423.178}{486.736}{19.811}{ALIUSC000000}{\ALIUSFontCormorantRegular}{12.960}{(4),}
  \ALIUSPlacedTextContent{455.960}{486.736}{53.019}{ALIUSC000000}{\ALIUSFontCormorantRegular}{12.960}{1239–1252.}
  \ALIUSPlacedTextContent{126.951}{502.576}{214.421}{ALIUSC0000FF}{\ALIUSFontCormorantRegular}{12.960}{https://doi.org/10.1016/j.concog.2013.08.013}
  \ALIUSPlacedTextContent{90.951}{524.176}{417.973}{ALIUSC000000}{\ALIUSFontCormorantRegular}{12.960}{Shokur, S., O'Doherty, J. E., Winans, J. A., Bleuler, H., Lebedev, M. A., \& Nicolelis,}
  \ALIUSPlacedTextContent{126.951}{540.016}{382.007}{ALIUSC000000}{\ALIUSFontCormorantRegular}{12.960}{M. A. L. (2013). Expanding the primate body schema in sensorimotor cortex}
  \ALIUSPlacedTextContent{126.951}{555.616}{175.351}{ALIUSC000000}{\ALIUSFontCormorantRegular}{12.960}{by virtual touches of an avatar.}
  \ALIUSPlacedTextContent{305.612}{555.616}{203.433}{ALIUSC000000}{\ALIUSFontCormorantItalic}{12.960}{Proceedings of the National Academy of}
  \ALIUSPlacedTextContent{126.951}{571.456}{56.292}{ALIUSC000000}{\ALIUSFontCormorantItalic}{12.960}{Sciences, 110}
  \ALIUSPlacedTextContent{183.239}{571.456}{85.587}{ALIUSC000000}{\ALIUSFontCormorantRegular}{12.960}{(37), 15121–15126.}
  \ALIUSPlacedTextContent{268.842}{571.456}{196.228}{ALIUSC0000FF}{\ALIUSFontCormorantRegular}{12.960}{https://doi.org/10.1073/pnas.1308459110}
  \ALIUSPlacedTextContent{90.951}{593.056}{417.954}{ALIUSC000000}{\ALIUSFontCormorantRegular}{12.960}{Slater, M. (2009). Place illusion and plausibility can lead to realistic behaviour in}
  \ALIUSPlacedTextContent{126.951}{608.896}{175.631}{ALIUSC000000}{\ALIUSFontCormorantRegular}{12.960}{immersive virtual environments.}
  \ALIUSPlacedTextContent{306.728}{608.896}{202.241}{ALIUSC000000}{\ALIUSFontCormorantItalic}{12.960}{Philosophical Transactions of the Royal}
  \ALIUSPlacedTextContent{126.951}{624.736}{154.569}{ALIUSC000000}{\ALIUSFontCormorantItalic}{12.960}{Society B: Biological Sciences, 364}
  \ALIUSPlacedTextContent{281.530}{624.736}{92.848}{ALIUSC000000}{\ALIUSFontCormorantRegular}{12.960}{(1535), 3549–3557.}
  \ALIUSPlacedTextContent{126.951}{640.576}{189.119}{ALIUSC0000FF}{\ALIUSFontCormorantRegular}{12.960}{https://doi.org/10.1098/rstb.2009.0138}
  \ALIUSPlacedTextContent{90.951}{662.176}{417.986}{ALIUSC000000}{\ALIUSFontCormorantRegular}{12.960}{Smith, A. J. T. (2010). Comment: Minimal Conditions for the Simplest Form of}
  \ALIUSPlacedTextContent{126.951}{677.776}{381.997}{ALIUSC000000}{\ALIUSFontCormorantRegular}{12.960}{Self-Consciousness. In T. Fuchs, H. Sattel, \& P. Henningsen (Eds.), The}
  \ALIUSPlacedTextContent{126.951}{693.616}{310.462}{ALIUSC000000}{\ALIUSFontCormorantRegular}{12.960}{embodied self: Dimensions, coherence, disorders. Schattauer.}
  \ALIUSPlacedTextContent{90.951}{715.456}{418.001}{ALIUSC000000}{\ALIUSFontCormorantRegular}{12.960}{Timmermann, C., Roseman, L., Schartner, M., Milliere, R., Williams, L. T. J.,}
  \ALIUSPlacedTextContent{126.951}{731.056}{382.092}{ALIUSC000000}{\ALIUSFontCormorantRegular}{12.960}{Erritzoe, D., Muthukumaraswamy, S., Ashton, M., Bendrioua, A., Kaur, O.,}
  \ALIUSPlacedTextContent{126.951}{746.896}{382.042}{ALIUSC000000}{\ALIUSFontCormorantRegular}{12.960}{Turton, S., Nour, M. M., Day, C. M., Leech, R., Nutt, D. J., \& Carhart-Harris,}
  \ALIUSPlacedTextContent{90.951}{793.143}{122.397}{ALIUSC767171}{\ALIUSFontCormorantLight}{12.000}{ALIUS Bulletin n°4 (2020)}
  \ALIUSPlacedTextContent{387.435}{793.143}{120.728}{ALIUSC767171}{\ALIUSFontCormorantLight}{12.000}{aliusresearch.org/bulletin}
\end{tikzpicture}
\null
\newpage

% Page 24
\thispagestyle{empty}
\noindent\begin{tikzpicture}[remember picture,overlay,x=1bp,y=1bp,shift={(current page.north west)}]
  \fill[ALIUSC0000FF] (126.951bp,-117.351bp) rectangle ++(211.680bp,-0.720bp);
  \fill[ALIUSC0000FF] (265.191bp,-170.631bp) rectangle ++(159.360bp,-0.720bp);
  \fill[ALIUSC0000FF] (220.791bp,-223.911bp) rectangle ++(247.440bp,-0.720bp);
  \fill[ALIUSC0000FF] (126.951bp,-308.631bp) rectangle ++(205.680bp,-0.720bp);
  \fill[ALIUSC0000FF] (202.311bp,-377.511bp) rectangle ++(205.200bp,-0.720bp);
  \fill[ALIUSC0000FF] (142.071bp,-430.791bp) rectangle ++(199.680bp,-0.720bp);
  \ALIUSPlacedTextContent{90.951}{36.329}{258.472}{ALIUSC767171}{\ALIUSFontCormorantRegular}{13.920}{Lenggenhager – Bodily boundaries and beyond}
  \ALIUSPlacedTextContent{488.871}{36.329}{20.309}{ALIUSC000000}{\ALIUSFontCormorantRegular}{13.920}{105}
  \ALIUSPlacedTextContent{126.951}{72.976}{382.033}{ALIUSC000000}{\ALIUSFontCormorantRegular}{12.960}{R. L. (2019). Neural correlates of the DMT experience assessed with}
  \ALIUSPlacedTextContent{126.951}{88.576}{94.009}{ALIUSC000000}{\ALIUSFontCormorantRegular}{12.960}{multivariate EEG.}
  \ALIUSPlacedTextContent{220.977}{88.576}{91.899}{ALIUSC000000}{\ALIUSFontCormorantItalic}{12.960}{Scientific Reports, 9}
  \ALIUSPlacedTextContent{312.872}{88.576}{47.297}{ALIUSC000000}{\ALIUSFontCormorantRegular}{12.960}{(1), 1–13.}
  \ALIUSPlacedTextContent{126.951}{104.416}{211.760}{ALIUSC0000FF}{\ALIUSFontCormorantRegular}{12.960}{https://doi.org/10.1038/s41598-019-51974-4}
  \ALIUSPlacedTextContent{90.951}{126.016}{417.871}{ALIUSC000000}{\ALIUSFontCormorantRegular}{12.960}{Vollenweider, F. X., \& Kometer, M. (2010). The neurobiology of psychedelic drugs:}
  \ALIUSPlacedTextContent{126.951}{141.856}{292.893}{ALIUSC000000}{\ALIUSFontCormorantRegular}{12.960}{Implications for the treatment of mood disorders.}
  \ALIUSPlacedTextContent{426.093}{141.856}{82.879}{ALIUSC000000}{\ALIUSFontCormorantItalic}{12.960}{Nature Reviews}
  \ALIUSPlacedTextContent{126.951}{157.696}{73.133}{ALIUSC000000}{\ALIUSFontCormorantItalic}{12.960}{Neuroscience, 11}
  \ALIUSPlacedTextContent{200.100}{157.696}{65.110}{ALIUSC000000}{\ALIUSFontCormorantRegular}{12.960}{(9), 642–651.}
  \ALIUSPlacedTextContent{265.216}{157.696}{159.253}{ALIUSC0000FF}{\ALIUSFontCormorantRegular}{12.960}{https://doi.org/10.1038/nrn2884}
  \ALIUSPlacedTextContent{90.951}{179.296}{418.132}{ALIUSC000000}{\ALIUSFontCormorantRegular}{12.960}{Wada, M., Takano, K., Ora, H., Ide, M., \& Kansaku, K. (2016). The Rubber Tail}
  \ALIUSPlacedTextContent{126.951}{195.136}{261.902}{ALIUSC000000}{\ALIUSFontCormorantRegular}{12.960}{Illusion as Evidence of Body Ownership in Mice.}
  \ALIUSPlacedTextContent{390.826}{195.136}{118.099}{ALIUSC000000}{\ALIUSFontCormorantItalic}{12.960}{Journal of Neuroscience,}
  \ALIUSPlacedTextContent{126.951}{210.976}{9.888}{ALIUSC000000}{\ALIUSFontCormorantItalic}{12.960}{36}
  \ALIUSPlacedTextContent{136.856}{210.976}{83.811}{ALIUSC000000}{\ALIUSFontCormorantRegular}{12.960}{(43), 11133–11137.}
  \ALIUSPlacedTextContent{220.678}{210.976}{247.579}{ALIUSC0000FF}{\ALIUSFontCormorantRegular}{12.960}{https://doi.org/10.1523/JNEUROSCI.3006-15.2016}
  \ALIUSPlacedTextContent{90.951}{232.576}{418.033}{ALIUSC000000}{\ALIUSFontCormorantRegular}{12.960}{Walsh, E., Guilmette, D. N., Longo, M. R., Moore, J. W., Oakley, D. A., Halligan, P.}
  \ALIUSPlacedTextContent{126.951}{248.416}{382.018}{ALIUSC000000}{\ALIUSFontCormorantRegular}{12.960}{W., Mehta, M. A., \& Deeley, Q. (2015). Are You Suggesting That's My Hand?}
  \ALIUSPlacedTextContent{126.951}{264.016}{382.022}{ALIUSC000000}{\ALIUSFontCormorantRegular}{12.960}{The Relation Between Hypnotic Suggestibility and the Rubber Hand}
  \ALIUSPlacedTextContent{126.951}{279.856}{43.947}{ALIUSC000000}{\ALIUSFontCormorantRegular}{12.960}{Illusion.}
  \ALIUSPlacedTextContent{170.876}{279.856}{64.858}{ALIUSC000000}{\ALIUSFontCormorantItalic}{12.960}{Perception, 44}
  \ALIUSPlacedTextContent{235.744}{279.856}{69.019}{ALIUSC000000}{\ALIUSFontCormorantRegular}{12.960}{(6), 709–723.}
  \ALIUSPlacedTextContent{126.951}{295.696}{205.565}{ALIUSC0000FF}{\ALIUSFontCormorantRegular}{12.960}{https://doi.org/10.1177/0301006615594266}
  \ALIUSPlacedTextContent{90.951}{317.296}{418.073}{ALIUSC000000}{\ALIUSFontCormorantRegular}{12.960}{Watts, R., Day, C., Krzanowski, J., Nutt, D., \& Carhart-Harris, R. (2017). Patients'}
  \ALIUSPlacedTextContent{126.951}{333.136}{382.041}{ALIUSC000000}{\ALIUSFontCormorantRegular}{12.960}{Accounts of Increased "Connectedness" and "Acceptance" After Psilocybin}
  \ALIUSPlacedTextContent{126.951}{348.736}{199.148}{ALIUSC000000}{\ALIUSFontCormorantRegular}{12.960}{for Treatment-Resistant Depression.}
  \ALIUSPlacedTextContent{331.290}{348.736}{177.626}{ALIUSC000000}{\ALIUSFontCormorantItalic}{12.960}{Journal of Humanistic Psychology,}
  \ALIUSPlacedTextContent{126.951}{364.576}{9.877}{ALIUSC000000}{\ALIUSFontCormorantItalic}{12.960}{57}
  \ALIUSPlacedTextContent{136.843}{364.576}{65.529}{ALIUSC000000}{\ALIUSFontCormorantRegular}{12.960}{(5), 520–564.}
  \ALIUSPlacedTextContent{202.388}{364.576}{204.982}{ALIUSC0000FF}{\ALIUSFontCormorantRegular}{12.960}{https://doi.org/10.1177/0022167817709585}
  \ALIUSPlacedTextContent{90.951}{386.176}{417.951}{ALIUSC000000}{\ALIUSFontCormorantRegular}{12.960}{Weech, S., Kenny, S., \& Barnett-Cowan, M. (2019). Presence and Cybersickness in}
  \ALIUSPlacedTextContent{126.951}{402.016}{264.888}{ALIUSC000000}{\ALIUSFontCormorantRegular}{12.960}{Virtual Reality Are Negatively Related: A Review.}
  \ALIUSPlacedTextContent{393.031}{402.016}{115.953}{ALIUSC000000}{\ALIUSFontCormorantItalic}{12.960}{Frontiers in Psychology,}
  \ALIUSPlacedTextContent{126.951}{417.856}{9.460}{ALIUSC000000}{\ALIUSFontCormorantItalic}{12.960}{10}
  \ALIUSPlacedTextContent{136.427}{417.856}{5.594}{ALIUSC000000}{\ALIUSFontCormorantRegular}{12.960}{.}
  \ALIUSPlacedTextContent{142.030}{417.856}{199.655}{ALIUSC0000FF}{\ALIUSFontCormorantRegular}{12.960}{https://doi.org/10.3389/fpsyg.2019.00158}
  \ALIUSPlacedTextContent{90.951}{793.143}{122.397}{ALIUSC767171}{\ALIUSFontCormorantLight}{12.000}{ALIUS Bulletin n°4 (2020)}
  \ALIUSPlacedTextContent{387.435}{793.143}{120.728}{ALIUSC767171}{\ALIUSFontCormorantLight}{12.000}{aliusresearch.org/bulletin}
\end{tikzpicture}
\null

\ifdefined\ALIUSIssueBuild
\clearpage
\expandafter\endinput
\fi
\end{document}
